% Production d’écrits utilitaires : production d’une lettre officielle : la lettre de motivation — Français 1re Tech — Cours signé OnBuch+
% Programme officiel MINESEC. Compiler avec : tectonic N09-production-ecrits-utilitaires-production-lettre-officie.tex
\newcommand{\DOCMATIERE}{Français}
\newcommand{\DOCNIVEAU}{1re Tech}
\newcommand{\DOCMODULE}{Français – classe de Première technique}
\newcommand{\DOCLECON}{N09}
\newcommand{\DOCTITRE}{Production d’écrits utilitaires : production d’une lettre officielle : la lettre de motivation}
% OnBuch+ — préambule « cours signés » ENRICHI (compilé avec tectonic / XeLaTeX)
% Système visuel « Pop » d'OnBuch+ : fond crème (jamais blanc), bordures encre
% épaisses, ombres « dures » décalées, titres Archivo Black en capitales.
% Version MATHÉMATIQUES : graphes (pgfplots/tikz), boîtes pédagogiques
% (définition, propriété, méthode, attention, le savais-tu, exercice résolu,
% exercice, TP, bilan), page de garde et signature OnBuch+ ancrée en bas.
%
% Macros attendues dans main.tex AVANT \input{preamble} :
%   \DOCMATIERE  \DOCNIVEAU  \DOCMODULE  \DOCLECON (numéro)  \DOCTITRE
\documentclass[11pt]{article}

\usepackage{fontspec}
\usepackage[french]{babel}
\usepackage[a4paper,margin=2cm,top=3.4cm,bottom=2.8cm,headheight=1.4cm,footskip=1.3cm]{geometry}
\usepackage{amsmath,amssymb,amsfonts}
\usepackage{mathtools} % \xmapsto, \coloneqq... souvent utilisés par les modèles
\usepackage{cancel} % simplifications barrées (bilans de forces)
% \abs{z} (module, valeur absolue) et \norm{u} (norme), utilisés par les modèles
\providecommand{\abs}{}\let\abs\relax\DeclarePairedDelimiter{\abs}{\lvert}{\rvert}
\providecommand{\norm}{}\let\norm\relax\DeclarePairedDelimiter{\norm}{\lVert}{\rVert}
\usepackage[table]{xcolor} % [table] : fournit aussi \rowcolors (alternance de lignes)
\usepackage{pagecolor}
\usepackage{tikz}
\usetikzlibrary{arrows.meta,positioning,calc,shapes.geometric,shapes.misc,shapes.arrows,decorations.pathmorphing,decorations.pathreplacing,decorations.markings,patterns,shadows,mindmap,trees,fit,backgrounds,matrix,intersections,angles,quotes}
\usepackage{pgfplots}
\usepackage[version=4]{mhchem} % équations nucléaires \ce{^{235}_{92}U}
\usepackage[european,siunitx]{circuitikz} % schémas électriques
\pgfplotsset{compat=1.18}
\usepgfplotslibrary{fillbetween,groupplots} % aires entre courbes (calcul intégral)
\usepackage{enumitem}
\usepackage{fancyhdr}
% [most] charge toutes les bibliothèques tcolorbox (listings, minted, raster,
% documentation...) et gonfle inutilement la mémoire TeX sur un document long
% (~300 boîtes) : on ne charge que ce qui est réellement utilisé.
\usepackage[skins,breakable]{tcolorbox}
\usepackage{graphicx}
\usepackage{colortbl}
\usepackage{booktabs}
\usepackage{tabularx}
\usepackage{multirow}
\usepackage{diagbox} % case d'en-tête diagonale des tableaux à double entrée
% Colonnes extensibles centrée / gauche / droite (C, L, R), souvent utilisées par les modèles
\newcolumntype{C}{>{\centering\arraybackslash}X}
\newcolumntype{L}{>{\raggedright\arraybackslash}X}
\newcolumntype{R}{>{\raggedleft\arraybackslash}X}
\usepackage{array}
\usepackage{multicol}
\usepackage{siunitx}
% parse-numbers=false : accepte aussi \SI{2,5 \times 10^{-3}}{...}
\sisetup{locale=FR,output-decimal-marker={,},group-minimum-digits=4,parse-numbers=false}
\DeclareSIUnit\ohm{\ensuremath{\symup{\Omega}}} % U+2126 absent de la police
\DeclareSIUnit\division{div} % réglages d'oscilloscope (ms/div, V/div)
\DeclareSIUnit\tr{tr} % tours (vitesse de rotation, ex. \SI{1500}{\tr\per\minute})
\DeclareSIUnit\molar{M} % concentration molaire (ex. \SI{3}{\molar}, \SI{1}{\micro\molar})
\DeclareSIUnit\an{an} % année (le modèle confond parfois avec \year, un compteur TeX) — ex. \SI{5}{\kilo\meter\per\an}
\DeclareSIUnit\FCFA{FCFA} % franc CFA (ex. \SI{500}{\FCFA\per\kilo\gram})
\DeclareSIUnit\degreeBrix{\textdegree Brix} % teneur en sucre d'un sirop/jus (réfractométrie)
\DeclareSIUnit\month{mois} % le modèle utilise parfois \month, un compteur TeX, comme unité de durée
\DeclareSIUnit\microliter{\micro\liter} % le modèle écrit parfois \microliter en un seul token
\DeclareSIUnit\million{millions} % dénombrement (ex. \SI{15}{\million\per\mL} spermatozoïdes)
\usepackage{qrcode}
% \degree : commande fréquemment utilisée par les modèles pour un angle ou une
% température en dehors de \SI{}{} (non fournie par les paquets chargés).
\providecommand{\degree}{^{\circ}}
% \qed (fin de démonstration), utilisé par les modèles sans amsthm
\providecommand{\qed}{\hfill$\square$}
% \barre : double barre des valeurs interdites dans un tableau de variations (tkz-tab)
\providecommand{\barre}{\ensuremath{\|}}
% environnement proof (amsthm non chargé) utilisé par les modèles
\ifdefined\proof\else\newenvironment{proof}[1][Démonstration]{\par\noindent\textit{#1.}\ }{\qed\par}\fi
\usepackage{lastpage}
\usepackage{hyperref}
\hypersetup{hidelinks}

% ============================================================
% Palette « Pop » OnBuch+ — mêmes hex que la classe P (plus_ui.dart)
% ============================================================
\definecolor{poporange}{HTML}{F59321}
\definecolor{popdark}{HTML}{E07A0C}
\definecolor{popcream}{HTML}{FFF6E8}
\definecolor{popink}{HTML}{1C1714}
\definecolor{poppurple}{HTML}{7A5AE0}
\definecolor{popgreen}{HTML}{2E9E6A}
\definecolor{poppink}{HTML}{E0457B}
\definecolor{popblue}{HTML}{2F7DE1}
\definecolor{popgold}{HTML}{C98A12}
\definecolor{popmuted}{HTML}{8A7F76}
\definecolor{popline}{HTML}{EADFD2}
\definecolor{popgray}{HTML}{8A7F76}
\colorlet{popgrey}{popgray}
% Alias tolérants (le modèle nomme parfois les couleurs différemment)
\colorlet{popwhite}{white}
\colorlet{popblack}{popink}
\colorlet{popred}{poppink}
\colorlet{popyellow}{popgold}
% Teintes claires (fonds de tableaux/encadrés) — le modèle nomme parfois
% les couleurs avec un suffixe L (ex. popblueL) sans les avoir définies.
\colorlet{poporangeL}{poporange!20}
\colorlet{popdarkL}{popdark!20}
\colorlet{poppurpleL}{poppurple!20}
\colorlet{popgreenL}{popgreen!20}
\colorlet{poppinkL}{poppink!20}
\colorlet{popblueL}{popblue!20}
\colorlet{popgoldL}{popgold!20}
\colorlet{popredL}{popred!20}
\colorlet{popgrayL}{popgray!20}
% Teintes claires pour fonds de tableaux / zones
\colorlet{popblueL}{popblue!10!white}
\colorlet{poporangeL}{poporange!14!white}
\colorlet{popgreenL}{popgreen!12!white}
\colorlet{poppurpleL}{poppurple!12!white}
\colorlet{poppinkL}{poppink!12!white}
\colorlet{popgoldL}{popgold!14!white}

\pagecolor{popcream}

% ============================================================
% Typographie
% ============================================================
\defaultfontfeatures{Ligatures=TeX}
\setmainfont{PlusJakartaSans-Medium.ttf}[Path=../../../fonts/,BoldFont=PlusJakartaSans-Bold.ttf]
% Maths en sans-serif (Fira Math) avec chiffres et lettres droites de Plus
% Jakarta, pour que les formules se fondent dans le texte.
\usepackage[math-style=ISO,bold-style=ISO]{unicode-math}
\setmathfont{FiraMath-Regular.otf}[Path=../../../fonts/]
\setmathfont{PlusJakartaSans-Medium.ttf}[Path=../../../fonts/,range={up/num,up/{Latin,latin},\mathup}]
\usepackage{icomma} % virgule décimale sans espace en mode math
% Caractères mathématiques Unicode absents des polices de texte (Plus Jakarta,
% Archivo Black) : rendus par leur équivalent mathématique, en texte comme en
% formule — sinon ils s'affichent en carré vide (ex. « Arithmétique dans ℤ »).
\usepackage{newunicodechar}
\newunicodechar{ℝ}{\ensuremath{\mathbb{R}}}
\newunicodechar{ℤ}{\ensuremath{\mathbb{Z}}}
\newunicodechar{ℂ}{\ensuremath{\mathbb{C}}}
\newunicodechar{ℕ}{\ensuremath{\mathbb{N}}}
\newunicodechar{ℚ}{\ensuremath{\mathbb{Q}}}
\newunicodechar{𝔻}{\ensuremath{\mathbb{D}}}
\newunicodechar{α}{\ensuremath{\alpha}}
\newunicodechar{β}{\ensuremath{\beta}}
\newunicodechar{γ}{\ensuremath{\gamma}}
\newunicodechar{δ}{\ensuremath{\delta}}
\newunicodechar{ε}{\ensuremath{\varepsilon}}
\newunicodechar{θ}{\ensuremath{\theta}}
\newunicodechar{λ}{\ensuremath{\lambda}}
\newunicodechar{μ}{\ensuremath{\mu}}
\newunicodechar{σ}{\ensuremath{\sigma}}
\newunicodechar{τ}{\ensuremath{\tau}}
\newunicodechar{φ}{\ensuremath{\varphi}}
\newunicodechar{ψ}{\ensuremath{\psi}}
\newunicodechar{ω}{\ensuremath{\omega}}
\newunicodechar{Σ}{\ensuremath{\Sigma}}
% majuscules produites par \MakeUppercase dans les titres (α-aminés -> Α)
\newunicodechar{Α}{\ensuremath{\alpha}}
\newunicodechar{Β}{\ensuremath{\beta}}
\newunicodechar{Γ}{\ensuremath{\Gamma}}
\newunicodechar{Δ}{\ensuremath{\Delta}}
\newunicodechar{Ε}{\ensuremath{\varepsilon}}
\newunicodechar{Θ}{\ensuremath{\Theta}}
\newunicodechar{Λ}{\ensuremath{\Lambda}}
\newunicodechar{Μ}{\ensuremath{\mu}}
\newunicodechar{Τ}{\ensuremath{\tau}}
\newunicodechar{Φ}{\ensuremath{\Phi}}
\newunicodechar{Ψ}{\ensuremath{\Psi}}
\newunicodechar{Ω}{\ensuremath{\Omega}}
\newunicodechar{↦}{\ensuremath{\mapsto}}
\newunicodechar{∈}{\ensuremath{\in}}
\newunicodechar{∉}{\ensuremath{\notin}}
\newunicodechar{∧}{\ensuremath{\wedge}}
\newunicodechar{⇔}{\ensuremath{\Leftrightarrow}}
\newunicodechar{⟺}{\ensuremath{\Longleftrightarrow}}
\newunicodechar{⇒}{\ensuremath{\Rightarrow}}
\newunicodechar{⟹}{\ensuremath{\Longrightarrow}}
\newunicodechar{∘}{\ensuremath{\circ}}
\newunicodechar{∪}{\ensuremath{\cup}}
\newunicodechar{∩}{\ensuremath{\cap}}
\newunicodechar{≡}{\ensuremath{\equiv}}
\newunicodechar{⊂}{\ensuremath{\subset}}
\newunicodechar{⊥}{\ensuremath{\perp}}
\newunicodechar{∃}{\ensuremath{\exists}}
\newunicodechar{∀}{\ensuremath{\forall}}
\newunicodechar{∛}{\ensuremath{\sqrt[3]{\,}}}
\newunicodechar{∜}{\ensuremath{\sqrt[4]{\,}}}
\newunicodechar{⃗}{\ensuremath{{}^{\to}}}
\newunicodechar{ⁿ}{\ifmmode^{n}\else\textsuperscript{\ensuremath{n}}\fi}
\newunicodechar{ˣ}{\ifmmode^{x}\else\textsuperscript{\ensuremath{x}}\fi}
\newunicodechar{ᵇ}{\ifmmode^{b}\else\textsuperscript{\ensuremath{b}}\fi}
\newunicodechar{ʳ}{\ifmmode^{r}\else\textsuperscript{\ensuremath{r}}\fi}
\newunicodechar{ᵘ}{\ifmmode^{u}\else\textsuperscript{\ensuremath{u}}\fi}
\newunicodechar{ʸ}{\ifmmode^{y}\else\textsuperscript{\ensuremath{y}}\fi}
\newunicodechar{ᵃ}{\ifmmode^{a}\else\textsuperscript{\ensuremath{a}}\fi}
\newunicodechar{ᵉ}{\ifmmode^{e}\else\textsuperscript{\ensuremath{e}}\fi}
\newunicodechar{ᵏ}{\ifmmode^{k}\else\textsuperscript{\ensuremath{k}}\fi}
\newunicodechar{ᵖ}{\ifmmode^{p}\else\textsuperscript{\ensuremath{p}}\fi}
\newunicodechar{ᴺ}{\ifmmode^{N}\else\textsuperscript{\ensuremath{N}}\fi}
\newunicodechar{ᶜ}{\ifmmode^{c}\else\textsuperscript{\ensuremath{c}}\fi}
\newunicodechar{ʰ}{\ifmmode^{h}\else\textsuperscript{\ensuremath{h}}\fi}
\newunicodechar{ᵠ}{\ifmmode^{q}\else\textsuperscript{\ensuremath{q}}\fi}
\newunicodechar{ₙ}{\ifmmode_{n}\else\textsubscript{\ensuremath{n}}\fi}
\newunicodechar{ₐ}{\ifmmode_{a}\else\textsubscript{\ensuremath{a}}\fi}
\newunicodechar{ᵢ}{\ifmmode_{i}\else\textsubscript{\ensuremath{i}}\fi}
\newunicodechar{ᵣ}{\ifmmode_{r}\else\textsubscript{\ensuremath{r}}\fi}
\newunicodechar{ₖ}{\ifmmode_{k}\else\textsubscript{\ensuremath{k}}\fi}
\newunicodechar{ᵦ}{\ifmmode_{\beta}\else\textsubscript{\ensuremath{\beta}}\fi}
\newunicodechar{ₓ}{\ifmmode_{x}\else\textsubscript{\ensuremath{x}}\fi}
\newfontfamily\popdisplay{ArchivoBlack-Regular.ttf}[Path=../../../fonts/]
\newfontfamily\poplogo{SpaceGrotesk-Bold.ttf}[Path=../../../fonts/]
\newfontfamily\popsemi{PlusJakartaSans-SemiBold.ttf}[Path=../../../fonts/]
\newfontfamily\popxbold{PlusJakartaSans-ExtraBold.ttf}[Path=../../../fonts/]
\newfontfamily\popxboldls{PlusJakartaSans-ExtraBold.ttf}[Path=../../../fonts/,LetterSpace=3.0]

\setlist[enumerate]{leftmargin=1.7em,itemsep=5pt,topsep=4pt}
\setlist[itemize]{leftmargin=1.7em,itemsep=4pt,topsep=4pt,label={\color{poporange}\textbullet}}
\setlength{\parindent}{0pt}
\setlength{\parskip}{5pt}
\renewcommand{\arraystretch}{1.25}

% pgfplots : style Pop par défaut
\pgfplotsset{
  every axis/.append style={axis line style={popink,line width=1pt},tick style={popink},
    grid style={popline,line width=0.5pt},label style={font=\small},tick label style={font=\footnotesize}},
  popcurve/.style={poporange,line width=1.8pt,no markers,smooth},
}
% Même style disponible dans un \draw TikZ simple (hors pgfplots)
\tikzset{popcurve/.style={poporange,line width=1.8pt}}
\tikzset{popgrid/.style={popline,very thin}}
\colorlet{popcurve}{poporange} % le nom du style est parfois utilisé comme couleur

% ============================================================
% Pill / squiggle
% ============================================================
\newcommand{\poppill}[3]{%
  \tikz[baseline=-0.55ex]\node[draw=popink,line width=1.2pt,rounded corners=2.6pt,%
    fill=#2,inner xsep=6.5pt,inner ysep=3pt]{%
    \popxboldls\fontsize{7}{7}\selectfont\color{#3}\MakeUppercase{#1}};%
}
\newcommand{\popsquiggle}[1]{%
  \tikz[baseline=-0.4ex]{
    \draw[#1,line width=2.6pt,line cap=round]
      plot [smooth] coordinates {(0,0.09) (0.34,0.01) (0.68,0.21) (1.02,0.01) (1.36,0.10)};
  }%
}
% Mot-clé surligné (marqueur orange sous le texte)
\newcommand{\cle}[1]{{\popxbold\color{popdark}#1}}

% ============================================================
% En-tête / pied
% ============================================================
\pagestyle{fancy}
\fancyhf{}
\renewcommand{\headrulewidth}{0pt}
\renewcommand{\footrulewidth}{0pt}
\lhead{\raisebox{-0.15ex}{\poplogo\fontsize{13}{13}\selectfont\color{popink}OnBuch\color{poporange}+}}
\rhead{\poppill{\DOCMATIERE\ \textbullet\ \DOCNIVEAU\ \textbullet\ Leçon \DOCLECON}{white}{popink}}
\lfoot{\popxbold\fontsize{7.6}{7.6}\selectfont\color{popmuted}\MakeUppercase{Cours signé OnBuch+}\ \textbullet\ {\popsemi\DOCTITRE}}
\rfoot{\tikz[baseline=-0.6ex]\node[rounded corners=8.5pt,fill=popink,minimum height=17pt,minimum width=17pt,inner xsep=5pt,inner sep=0]{\popxbold\fontsize{8}{8}\selectfont\color{popcream}\thepage\,/\,\pageref*{LastPage}};}

% ============================================================
% Titres
% ============================================================
\newcommand{\chaptitre}[2]{%
  \begin{tcolorbox}[enhanced,colback=poporange,colframe=popink,boxrule=2.6pt,arc=11pt,
      drop shadow=popink,left=15pt,right=15pt,top=12pt,bottom=11pt,before skip=3pt,after skip=18pt]
    \poppill{#2}{popink}{popcream}\par\vspace{9pt}
    {\raggedright\hyphenpenalty=10000\exhyphenpenalty=10000\popdisplay\fontsize{22}{24}\selectfont\color{popcream}\MakeUppercase{#1}\par}
  \end{tcolorbox}
}
% Section numérotée : pastille orange avec le numéro + titre Archivo
\newcounter{popsec}
\newcommand{\coursec}[1]{%
  \stepcounter{popsec}\setcounter{popsub}{0}%
  \par\vspace{16pt}\noindent
  \begin{minipage}{\linewidth}
  \tikz[baseline=-0.7ex]\node[draw=popink,line width=1.6pt,fill=poporange,rounded corners=4pt,minimum size=22pt,inner sep=0]{\popdisplay\fontsize{12}{12}\selectfont\color{popcream}\thepopsec};%
  \hspace{8pt}{\popdisplay\fontsize{15}{17}\selectfont\color{popink}\MakeUppercase{#1}}\par
  \vspace{4pt}\noindent{\color{poporange}\rule{36pt}{3.6pt}}
  \end{minipage}\par\nopagebreak\vspace{8pt}
}
\newcounter{popsub}
\newcommand{\courssub}[1]{%
  \stepcounter{popsub}%
  \par\vspace{9pt}\noindent{\popxbold\fontsize{11.5}{13}\selectfont\color{popink}{\color{poporange}\thepopsec.\thepopsub}\hspace{6pt}#1}\par\nopagebreak\vspace{4pt}
}

% ============================================================
% Boîtes pédagogiques — même recette : fond blanc, bordure épaisse colorée,
% ombre dure encre, badge plein. Titre optionnel [#1].
% ============================================================
\tcbset{popbase/.style={breakable,enhanced,colback=white,boxrule=2.3pt,arc=9pt,
  shadow={3pt}{-3pt}{0pt}{popink},left=14pt,right=14pt,top=11pt,bottom=11pt,before skip=11pt,after skip=15pt,
  fontupper=\popsemi\fontsize{10.6}{14.5}\selectfont\color{popink},
  fontlower=\popsemi\fontsize{10.6}{14.5}\selectfont\color{popink}}}
% Coche dessinée (le glyphe ✓ n'existe pas dans les polices du document)
\newcommand{\popcheck}[1]{\tikz[baseline=-0.5ex]\draw[#1,line width=1.6pt,line cap=round,line join=round](-0.13,0.0)--(-0.04,-0.09)--(0.13,0.1);}
\providecommand{\coche}{\popcheck{popgreen}}
\providecommand{\resultat}[1]{\par\smallskip\noindent\fcolorbox{popgreen}{popgreenL}{\parbox{\dimexpr\linewidth-2\fboxsep-2\fboxrule\relax}{\textbf{Résultat.} #1}}\par\smallskip}
\newcommand{\popboxhead}[3]{\poppill{#1}{#2}{white}\ifx\relax#3\relax\else\hspace{6pt}{\popxbold\fontsize{10.6}{12}\selectfont\color{popink}#3}\fi\par\vspace{7pt}}

\newtcolorbox{aretenir}[1][]{popbase,colframe=popgreen,before upper={\popboxhead{À retenir}{popgreen}{#1}}}
\newtcolorbox{exemplebox}[1][]{popbase,colframe=poporange,before upper={\popboxhead{Exemple}{poporange}{#1}}}
\newtcolorbox{definition}[1][]{popbase,colframe=popblue,before upper={\popboxhead{Définition}{popblue}{#1}}}
\newtcolorbox{propriete}[1][]{popbase,colframe=poppurple,before upper={\popboxhead{Propriété}{poppurple}{#1}}}
\newtcolorbox{methode}[1][]{popbase,colframe=popgold,before upper={\popboxhead{Méthode}{popgold}{#1}}}
\newtcolorbox{attention}[1][]{popbase,colframe=poppink,before upper={\popboxhead{Attention}{poppink}{#1}}}
\newtcolorbox{experience}[1][]{popbase,colframe=popblue,colback=popblueL,before upper={\popboxhead{Expérience}{popblue}{#1}}}
% « Le savais-tu ? » : carte violette pleine (contexte, culture, Cameroun)
\newtcolorbox{savaistu}[1][]{popbase,colback=poppurple,colframe=popink,
  fontupper=\popsemi\fontsize{10.4}{14.2}\selectfont\color{white},
  before upper={\poppill{Le savais-tu\,?}{white}{poppurple}\ifx\relax#1\relax\else\hspace{6pt}{\popxbold\color{white}#1}\fi\par\vspace{7pt}}}
% Exercice résolu : énoncé en haut, \tcblower puis la solution sur fond crème
\newcounter{popexo}\newcounter{popexr}
\newtcolorbox{exoresolu}[1][]{popbase,colframe=poporange,colbacklower=popcream,
  segmentation style={popink,line width=1.2pt,dashed},
  before upper={\stepcounter{popexr}\popboxhead{Exercice résolu \thepopexr}{poporange}{#1}},
  before lower={\poppill{Solution}{popink}{popcream}\par\vspace{6pt}}}
% Exercice (non résolu en place) — la correction est dans la partie Corrigés
\newtcolorbox{exercice}[1][]{popbase,colframe=popink,
  before upper={\stepcounter{popexo}\popboxhead{Exercice \thepopexo}{popink}{#1}}}
% Corrigé
\newtcolorbox{corrige}[1][]{popbase,colframe=popgreen,colback=popgreenL,
  before upper={\popboxhead{Corrigé}{popgreen}{#1}}}
% Objectifs / prérequis / bilan
\newtcolorbox{objectifs}{popbase,colframe=popink,colback=poporangeL,
  before upper={\popboxhead{Objectifs de la leçon}{popink}{}}}
\newtcolorbox{prerequis}{popbase,colframe=popink,colback=popblueL,
  before upper={\popboxhead{Prérequis}{popblue}{}}}
\newtcolorbox{bilan}{popbase,colframe=popink,colback=white,boxrule=2.8pt,
  before upper={\popboxhead{Fiche bilan}{poporange}{L'essentiel à connaître pour l'examen}}}
% Figure encadrée avec légende
\newtcolorbox{popfigure}[1][]{enhanced,breakable=false,colback=white,colframe=popline,boxrule=1.4pt,arc=8pt,
  left=10pt,right=10pt,top=10pt,bottom=8pt,before skip=10pt,after skip=12pt,halign=center,
  lower separated=false,halign lower=center,
  fontlower=\popsemi\fontsize{9}{11.5}\selectfont\color{popmuted},
  #1}
\newcommand{\legende}[1]{\par\vspace{6pt}{\popsemi\fontsize{9}{11.5}\selectfont\color{popmuted}#1}}

% ============================================================
% Signature OnBuch+
% ============================================================
\newcommand{\poplogoinline}[1]{{\poplogo\fontsize{#1}{#1}\selectfont\color{popink}OnBuch\color{poporange}+}}
\newcommand{\signatureonbuch}{%
  \begin{tcolorbox}[enhanced,colback=popink,colframe=popink,boxrule=2.2pt,arc=10pt,
      drop shadow=poporange,left=14pt,right=14pt,top=10pt,bottom=10pt,before skip=0pt,after skip=0pt]
    \tikz[baseline=-0.7ex]\node[circle,fill=popgreen,draw=popcream,line width=1.3pt,minimum size=22pt,inner sep=0]{\popcheck{white}};%
    \hspace{9pt}\begin{minipage}[c]{0.62\linewidth}
      {\popxbold\fontsize{12}{13}\selectfont\color{popcream}Signé\ \textbullet\ {\poplogo OnBuch\color{poporange}+}}\par\vspace{2pt}
      {\raggedright\popsemi\fontsize{8}{10.5}\selectfont\color{popline}\MakeUppercase{Équipe pédagogique OnBuch+}\\\MakeUppercase{Conforme au programme officiel MINESEC}\par}
    \end{minipage}\hfill
    {\poplogo\fontsize{22}{22}\selectfont\color{popcream}OnBuch\color{poporange}+}
  \end{tcolorbox}%
}

% ============================================================
% Page de garde
% #1 titre · #2 matière · #3 niveau · #4 module · #5 n° leçon · #6 durée ·
% #7 description / objectifs (texte ou itemize)
% ============================================================
\newcommand{\pagegarde}[7]{%
  \thispagestyle{empty}
  \newgeometry{a4paper,margin=1.7cm,top=1.6cm,bottom=1.4cm}%
  \begin{tikzpicture}[remember picture,overlay]
    \fill[poporange!18] ([xshift=-3.2cm,yshift=-3.6cm]current page.north east) circle (4.6cm);
    \fill[poppurple!14] ([xshift=2.4cm,yshift=7.6cm]current page.south west) circle (3.8cm);
  \end{tikzpicture}%
  {\poplogo\fontsize{30}{30}\selectfont\color{popink}OnBuch\color{poporange}+}\hfill
  \raisebox{0.6ex}{\poppill{Cours signé \textbullet\ Premium}{popink}{popcream}}\par
  \vspace{1pt}\popsquiggle{poppurple}\par
  \vspace{8pt}
  \begin{tcolorbox}[enhanced,colback=poporange,colframe=popink,boxrule=2.8pt,arc=13pt,
      drop shadow=popink,left=18pt,right=18pt,top=12pt,bottom=13pt,after skip=10pt]
    \poppill{#2 \textbullet\ #3}{popink}{popcream}\hspace{5pt}\poppill{#4}{white}{popink}\par\vspace{8pt}
    {\popdisplay\fontsize{14}{14}\selectfont\color{popink}LEÇON #5}\par\vspace{4pt}
    {\raggedright\hyphenpenalty=10000\exhyphenpenalty=10000\popdisplay\fontsize{25}{27}\selectfont\color{popcream}\MakeUppercase{#1}\par}
    \vspace{8pt}
    {\popxbold\fontsize{9.5}{9.5}\selectfont\color{popink}\MakeUppercase{Durée conseillée : #6}}
  \end{tcolorbox}
  \begin{tcolorbox}[enhanced,colback=white,colframe=popink,boxrule=2.2pt,arc=10pt,
      drop shadow=popink,left=16pt,right=16pt,top=10pt,bottom=10pt,after skip=8pt]
    \poppill{Dans cette leçon}{poppurple}{white}\par\vspace{5pt}
    \popsemi\fontsize{9.3}{12.6}\selectfont\color{popink}#7
  \end{tcolorbox}
  \noindent\begin{minipage}[t]{0.487\linewidth}
    \begin{tcolorbox}[enhanced,colback=popgreen,colframe=popink,boxrule=2.2pt,arc=10pt,
        drop shadow=popink,left=11pt,right=11pt,top=10pt,bottom=10pt,height=3.1cm,valign=center]
      \tcbox[colback=white,colframe=popink,boxrule=1.3pt,arc=5pt,boxsep=0pt,nobeforeafter,
        left=3pt,right=3pt,top=3pt,bottom=3pt,valign=center]{\qrcode[height=1.9cm]{https://play.google.com/store/apps/details?id=com.luvvix.onbuch&hl=fr}}%
      \hspace{8pt}\begin{minipage}[c]{0.5\linewidth}
        {\popdisplay\fontsize{11}{12.5}\selectfont\color{white}\MakeUppercase{Télécharge OnBuch}}\par\vspace{4pt}
        {\raggedright\popsemi\fontsize{8.4}{11}\selectfont\color{white}Gratuit \textbullet\ Google Play\\Tuteur IA, annales, concours, résultats\par}
      \end{minipage}
    \end{tcolorbox}
  \end{minipage}\hfill
  \begin{minipage}[t]{0.487\linewidth}
    \begin{tcolorbox}[enhanced,colback=poppurple,colframe=popink,boxrule=2.2pt,arc=10pt,
        drop shadow=popink,left=11pt,right=11pt,top=10pt,bottom=10pt,height=3.1cm,valign=center]
      \tikz[baseline=-0.7ex]\node[circle,fill=white,draw=popink,line width=1.3pt,minimum size=1.6cm,inner sep=0]{\popdisplay\fontsize{24}{24}\selectfont\color{poppurple}+};%
      \hspace{8pt}\begin{minipage}[c]{0.6\linewidth}
        {\popdisplay\fontsize{11}{12.5}\selectfont\color{white}\MakeUppercase{Passe à OnBuch+}}\par\vspace{4pt}
        {\raggedright\popsemi\fontsize{8.4}{11}\selectfont\color{white}Tous les cours signés, Tuteur IA illimité, fiches PDF — onglet OnBuch+ de l'app\par}
      \end{minipage}
    \end{tcolorbox}
  \end{minipage}
  \vspace{8pt}
  \popcontenu
  \vfill
  \signatureonbuch
  \restoregeometry
  \newpage
}

% Rangée « Ce cours contient » (page de garde)
\newcommand{\poptile}[3]{% #1 couleur #2 gros texte #3 légende
  \begin{tcolorbox}[enhanced,colback=#1,colframe=popink,boxrule=2pt,arc=9pt,shadow={2.5pt}{-2.5pt}{0pt}{popink},
      left=6pt,right=6pt,top=9pt,bottom=9pt,halign=center,valign=center,height=1.75cm,width=\linewidth,nobeforeafter]
    {\popdisplay\fontsize{12.5}{13}\selectfont\color{white}\MakeUppercase{#2}}\par\vspace{3pt}
    {\centering\popsemi\fontsize{7.8}{9.5}\selectfont\color{white}#3\par}
  \end{tcolorbox}}
\newcommand{\popcontenu}{%
  \par\noindent{\popxboldls\fontsize{7.5}{7.5}\selectfont\color{popmuted}CE COURS SIGNÉ CONTIENT}\par\vspace{7pt}
  \noindent\begin{minipage}[t]{0.235\linewidth}\poptile{popblue}{Cours}{complet, illustré, pas à pas}\end{minipage}\hfill
  \begin{minipage}[t]{0.235\linewidth}\poptile{poporange}{Méthodes}{et exercices résolus}\end{minipage}\hfill
  \begin{minipage}[t]{0.235\linewidth}\poptile{poppink}{Exercices}{type examen + corrigés}\end{minipage}\hfill
  \begin{minipage}[t]{0.235\linewidth}\poptile{popgreen}{Fiche bilan}{l'essentiel à retenir}\end{minipage}\par}

% Sommaire
\newtcolorbox{sommaire}{enhanced,colback=white,colframe=popink,boxrule=2.2pt,arc=10pt,
  drop shadow=popink,left=18pt,right=18pt,top=13pt,bottom=13pt,before skip=6pt,after skip=20pt,
  before upper={\poppill{Sommaire}{poppurple}{white}\par\vspace{9pt}\popsemi\fontsize{10.6}{14.5}\selectfont\color{popink}}}

% Signature de fin de cours — poussée tout en bas de la dernière page
\newcommand{\findecours}{%
  \par\vfill\nopagebreak
  \begin{center}\popsquiggle{poporange}\end{center}\vspace{4pt}
  \signatureonbuch
}
\AtBeginDocument{\def\llbracket{\mathopen{[\mkern-3.5mu[}}\def\rrbracket{\mathclose{]\mkern-3.5mu]}}} % intervalles d'entiers (glyphe ⟦ absent de la police)
% Glyphes absents de Fira Math : redéfinis à partir de glyphes disponibles
\AtBeginDocument{%
  \def\setminus{\mathbin{\backslash}}\let\smallsetminus\setminus
  \def\vdots{\mathord{\vbox{\baselineskip3.2pt\lineskiplimit0pt\kern4pt\hbox{.}\hbox{.}\hbox{.}}}}%
  \def\ddots{\mathinner{\mkern1mu\raise6.5pt\hbox{.}\mkern2mu\raise3.6pt\hbox{.}\mkern2mu\raise0.7pt\hbox{.}\mkern1mu}}%
  \def\triangle{\mathord{\scalebox{0.9}{$\symup{\Delta}$}}}%
  \def\nabla{\mathord{\raisebox{\depth}{\scalebox{1}[-1]{$\symup{\Delta}$}}}}%
  \def\checkmark{\popcheck{popdark}}%
  \def\star{\mathbin{\ast}}\def\bigstar{\mathord{\ast}}%
  \def\diamond{\mathbin{\scalebox{0.6}{$\Diamond$}}}%
  \def\bigcup{\mathop{\mathchoice{\raisebox{-0.3em}{\scalebox{1.7}{$\cup$}}}{\scalebox{1.25}{$\cup$}}{\cup}{\cup}}\displaylimits}%
  \def\bigcap{\mathop{\mathchoice{\raisebox{-0.3em}{\scalebox{1.7}{$\cap$}}}{\scalebox{1.25}{$\cap$}}{\cap}{\cap}}\displaylimits}%
  \def\lesseqqgtr{\mathrel{\lessgtr}}\def\bigoplus{\mathop{\oplus}\displaylimits}%
}
\providecommand{\ssi}{si et seulement si} % abréviation fréquente
\AtBeginDocument{\providecommand{\wideparen}{\overparen}} % arc de cercle

%% FIGFIX-BEGIN v5 — correctifs globaux des figures (docs/onbuch-generation/tools/figfix)
\usepackage{adjustbox}\usepackage{etoolbox}\tcbuselibrary{raster}
% toute figure TikZ/pgfplots plus large que la ligne est réduite à la largeur disponible
\BeforeBeginEnvironment{tikzpicture}{\begin{adjustbox}{max width=\linewidth}}
\AfterEndEnvironment{tikzpicture}{\end{adjustbox}}
% légendes pgfplots lisibles par-dessus les courbes
\pgfplotsset{every axis/.append style={legend style={fill=white,fill opacity=0.92,text opacity=1,draw=black!15,inner sep=3pt}}}
% étiquettes d'arêtes (syntaxe edge["..."]) avec fond blanc
\tikzset{every edge quotes/.append style={fill=white,inner sep=1.2pt}}
%% FIGFIX-END
\begin{document}

\pagegarde{Production d’écrits utilitaires : production d’une lettre officielle : la lettre de motivation}{Français}{1re Tech}{Français – classe de Première technique}{N09}{15 séances (fiche de progression harmonisée)}{Cette leçon outille l'élève de Première technique pour produire une lettre de motivation conforme aux normes de la correspondance officielle : en-tête, corps argumenté, formules de conclusion. Elle mobilise en appui la phrase complexe, l'emprunt lexical, les codes de la communication par internet et les tonalités dramatique et comique.
\vspace{4pt}
\begin{multicols}{2}\small
\begin{itemize}
\item À la fin de cette leçon, je sais identifier les caractéristiques d'un texte fonctionnel et d'un texte iconique.
\item À la fin de cette leçon, je sais rédiger l'en-tête d'une lettre de motivation (coordonnées, objet, lieu et date).
\item À la fin de cette leçon, je sais construire le corps d'une lettre de motivation en trois paragraphes (vous, moi, nous).
\item À la fin de cette leçon, je sais employer les formules de conclusion et de politesse adaptées au destinataire.
\item À la fin de cette leçon, je sais construire et analyser une phrase complexe (subordination relative, conjonctive, participiale).
\item À la fin de cette leçon, je sais reconnaître et employer correctement les emprunts dans la lettre officielle.
\end{itemize}\end{multicols}}

\begin{sommaire}
\begin{multicols}{2}
\begin{enumerate}
\item Textes fonctionnels et communication par internet
\item L'en-tête de la lettre de motivation
\item La phrase complexe et l'emprunt
\item Le corps de la lettre et les tonalités
\item Formules de conclusion et textes iconiques
\item Rédaction, confrontation et amélioration de la lettre
\item Activité d'intégration
\item Méthodes et exercices résolus
\item Exercices
\item Corrigés des exercices
\item Fiche bilan
\end{enumerate}
\end{multicols}
\end{sommaire}

\begin{objectifs}
\begin{itemize}
\item À la fin de cette leçon, je sais identifier les caractéristiques d'un texte fonctionnel et d'un texte iconique.
\item À la fin de cette leçon, je sais rédiger l'en-tête d'une lettre de motivation (coordonnées, objet, lieu et date).
\item À la fin de cette leçon, je sais construire le corps d'une lettre de motivation en trois paragraphes (vous, moi, nous).
\item À la fin de cette leçon, je sais employer les formules de conclusion et de politesse adaptées au destinataire.
\item À la fin de cette leçon, je sais construire et analyser une phrase complexe (subordination relative, conjonctive, participiale).
\item À la fin de cette leçon, je sais reconnaître et employer correctement les emprunts dans la lettre officielle.
\item À la fin de cette leçon, je sais distinguer les tonalités dramatique et comique dans un texte.
\item À la fin de cette leçon, je sais appliquer les codes de la communication par internet à la correspondance professionnelle (courriel de candidature).
\item À la fin de cette leçon, je sais rédiger, confronter et améliorer une lettre de motivation complète en justifiant ma démarche.
\end{itemize}
\end{objectifs}

\begin{prerequis}
\begin{itemize}
\item La lettre administrative et ses éléments de base (classe de Seconde)
\item La phrase simple et la phrase complexe : notions de proposition (classes de 4e et 3e)
\item Les classes de mots et les expansions du nom (début d'année de Première)
\item Les registres de langue : soutenu, courant, familier
\item La lecture méthodique d'un texte (situation d'énonciation, thème, thèse)
\end{itemize}
\end{prerequis}

\newpage

\coursec{Textes fonctionnels et communication par internet}

\courssub{Situation d'accroche et problématique}

À Douala, quartier Bonanjo, l'entreprise de BTP \emph{Bâtir Cameroun SARL} publie dans le journal \emph{Cameroon Tribune} une offre d'emploi : elle recherche un technicien en génie civil, titulaire d'un Baccalauréat technique ou d'un BTS, pour un chantier de construction d'un immeuble de quatre étages. Trois jeunes diplômés, tous sortis du même lycée technique, décident de postuler.

Le premier, pressé, envoie un simple SMS au numéro indiqué dans l'annonce : \og slt je suis interessé par le boulot, appelez moi \fg{}. Le deuxième rédige une lettre, mais sans en-tête, sans adresse du destinataire, sans objet, et il la termine par un simple \og merci \fg{}. Le troisième prend le temps d'écrire une lettre de motivation claire, structurée, avec un en-tête complet, un objet précis, un corps de lettre qui met en valeur sa formation et son stage pratique, et une formule de politesse respectueuse.

Une semaine plus tard, un seul candidat est convoqué à l'entretien : le troisième. Pourquoi ? Parce que, dans le monde professionnel, \cle{la forme compte autant que le fond}. Un texte mal présenté, malpoli ou rédigé en langage SMS décourage immédiatement le lecteur, avant même qu'il n'ait découvert les qualités du candidat.

\textbf{Problématique :} qu'est-ce qu'un texte fonctionnel ? Quelles règles faut-il maîtriser pour rédiger une lettre de motivation efficace, capable d'ouvrir les portes du monde professionnel au Cameroun ? Et comment la communication par internet (le courriel) transforme-t-elle nos habitudes d'écriture ?

\courssub{Qu'est-ce qu'un texte fonctionnel ?}

Commençons par une observation simple de la vie quotidienne. Autour de nous, nous lisons chaque jour des textes qui ne racontent pas une histoire et ne cherchent pas à nous émouvoir : le mode d'emploi d'un téléphone, l'affiche d'une réunion de parents d'élèves, le formulaire d'inscription à un concours, l'offre d'emploi publiée dans un journal, la facture d'électricité, le CV d'un candidat. Tous ces textes ont un point commun : ils visent une \cle{action concrète}. On les lit \emph{pour faire quelque chose} : acheter, s'inscrire, postuler, réparer, se déplacer.

\begin{definition}[Texte fonctionnel]
Un \cle{texte fonctionnel} (ou texte utilitaire) est un texte qui vise une action concrète et immédiate dans la vie courante ou professionnelle. Il informe, demande, ordonne ou propose, dans un but pratique précis. Exemples : la lettre officielle, la lettre de motivation, le curriculum vitae (CV), l'offre d'emploi, le formulaire administratif, l'affiche, le mode d'emploi, la convocation.
\end{definition}

Il faut bien distinguer le texte fonctionnel du \cle{texte littéraire}. Le roman, le poème ou le conte visent le plaisir esthétique, l'émotion, la réflexion ; leur langue est souvent riche d'images et de figures de style. Le texte fonctionnel, lui, doit être compris \emph{immédiatement et sans ambiguïté}.

\begin{center}
\begin{tabularx}{\linewidth}{|l|X|X|}
\hline
\rowcolor{popblueL} \textbf{Critère} & \textbf{Texte littéraire} & \textbf{Texte fonctionnel} \\
\hline
Objectif & Plaire, émouvoir, faire réfléchir & Informer, demander, obtenir une action \\
\hline
Langue & Images, figures de style, subjectivité & Clarté, précision, neutralité \\
\hline
Lecture & Plaisir, relecture possible & Lecture rapide, efficacité immédiate \\
\hline
Exemples & Roman, poème, conte, fable & Lettre, CV, affiche, mode d'emploi \\
\hline
\end{tabularx}
\end{center}

\begin{aretenir}[Caractéristiques du texte fonctionnel]
Un bon texte fonctionnel respecte quatre qualités :
\begin{enumerate}
\item la \cle{clarté} : chaque phrase a un seul sens possible ;
\item la \cle{concision} : on va à l'essentiel, sans détour ;
\item une \cle{mise en page aérée} : paragraphes courts, rubriques, titres, espaces ;
\item des \cle{informations hiérarchisées} : l'essentiel (poste, date, lieu, contact) est visible en premier.
\end{enumerate}
\end{aretenir}

\courssub{Étude du texte fonctionnel n°1 : une offre d'emploi}

Appliquons ces notions à un document réel de la vie économique camerounaise. Voici une offre d'emploi inspirée des annonces publiées par de grandes entreprises nationales (Camtel, Eneo, ou des PME du BTP).

\begin{exemplebox}[Offre d'emploi — Eneo Cameroon S.A.]
\textbf{ENEO CAMEROON S.A. — RECRUTEMENT}\par
\emph{Poste :} Technicien de maintenance des réseaux électriques (1 poste), Yaoundé.\par
\emph{Profil :} Baccalauréat technique (F3) ou BTS en électrotechnique ; moins de 30 ans ; permis B souhaité ; sens de la rigueur et du travail en équipe.\par
\emph{Conditions :} contrat à durée indéterminée, période d'essai de 6 mois ; salaire selon grille en vigueur.\par
\emph{Dossier de candidature :} lettre de motivation, CV, photocopie des diplômes, à déposer au plus tard le \textbf{30 novembre} au siège d'Eneo, ou par courriel : recrutement@eneo.cm.\par
\emph{Contact :} tél. 222 23 45 67.
\end{exemplebox}

\begin{exoresolu}[Repérer les rubriques d'une offre d'emploi]
\textbf{Énoncé.} Relevez dans l'offre d'emploi Eneo ci-dessus les cinq rubriques essentielles et indiquez pour chacune l'information qu'elle apporte.
\tcblower
\textbf{Solution détaillée.} Une offre d'emploi bien rédigée présente toujours des rubriques repérables :
\begin{enumerate}
\item \textbf{Le poste} : technicien de maintenance des réseaux électriques, à Yaoundé (1 poste) ;
\item \textbf{Le profil recherché} : diplôme (Bac technique F3 ou BTS électrotechnique), âge (moins de 30 ans), qualités (rigueur, esprit d'équipe) ;
\item \textbf{Les conditions} : type de contrat (CDI), période d'essai (6 mois), salaire (selon grille) ;
\item \textbf{Le délai et le dossier} : lettre de motivation, CV, diplômes, à déposer avant le 30 novembre ;
\item \textbf{Le contact} : adresse physique (siège d'Eneo), courriel (recrutement@eneo.cm), téléphone.
\end{enumerate}
Ces cinq rubriques montrent la \cle{hiérarchisation de l'information} : le lecteur trouve en quelques secondes si l'offre le concerne et comment postuler.
\end{exoresolu}

\courssub{La situation de communication d'un texte fonctionnel}

Pour comprendre un texte fonctionnel, il ne suffit pas de lire les mots : il faut identifier \emph{qui parle, à qui, pour obtenir quoi, par quel moyen}. C'est ce qu'on appelle la \cle{situation de communication}.

\begin{definition}[Situation de communication]
La situation de communication est l'ensemble des éléments qui entourent l'échange d'un message : l'\cle{émetteur} (celui qui envoie), le \cle{destinataire} ou récepteur (celui qui reçoit), le \cle{message} (le contenu), le \cle{canal} (le moyen de transmission : papier, téléphone, internet) et le \cle{code} (la langue et les conventions utilisées).
\end{definition}

\begin{popfigure}
\begin{center}
\begin{tikzpicture}[font=\small, node distance=1.1cm]
\node[draw=popblue, fill=popblueL, rounded corners, thick, minimum width=2.6cm, minimum height=1cm, align=center] (em){\textbf{Émetteur}\\ l'entreprise Eneo};
\node[draw=popgreen, fill=popgreenL, rounded corners, thick, minimum width=2.6cm, minimum height=1cm, right=4.2cm of em, align=center] (rec){\textbf{Récepteur}\\ le candidat};
\node[draw=poporange, fill=poporangeL, rounded corners, thick, minimum width=4.6cm, minimum height=1cm] at ($(em)!0.5!(rec)+(0,1.8)$) (msg) {\textbf{Message :} l'offre d'emploi};
\draw[-{Stealth[length=3mm]}, very thick, popblue] (em) -- (msg);
\draw[-{Stealth[length=3mm]}, very thick, popgreen] (msg) -- (rec);
\node[below=0.25cm of em, text=popdark, align=center] (canal) {\textbf{Canal :} journal, affiche,\\ internet (site, courriel)};
\node[below=0.25cm of rec, text=popdark, align=center] (code) {\textbf{Code :} français\\ formel, rubriques};
\draw[dashed, popmuted] (em) -- (canal);
\draw[dashed, popmuted] (rec) -- (code);
\end{tikzpicture}
\end{center}
\legende{Schéma de la situation de communication appliquée à l'offre d'emploi Eneo : l'entreprise (émetteur) transmet un message (l'offre) au candidat (récepteur) par un canal (journal ou internet) en utilisant un code (le français formel organisé en rubriques).}
\end{popfigure}

\begin{methode}[Analyser la situation de communication d'un texte fonctionnel]
Face à un texte fonctionnel à l'examen ou dans la vie courante, posez-vous cinq questions, dans l'ordre :
\begin{enumerate}
\item \textbf{Qui émet ?} Identifiez l'émetteur (une entreprise, une administration, un particulier) grâce à l'en-tête, au logo, à la signature.
\item \textbf{À qui ?} Identifiez le destinataire (un candidat, le public, un client, un chef de service).
\item \textbf{Quel message ?} Résumez le contenu en une phrase (une proposition d'emploi, une demande, une information).
\item \textbf{Quel objectif ?} Précisez l'action attendue (postuler, payer, se présenter, s'inscrire).
\item \textbf{Quel canal et quel code ?} Repérez le support (papier, courriel, affiche) et le niveau de langue (formel, technique).
\end{enumerate}
Concluez en une phrase complète : \og Cette offre d'emploi est émise par Eneo à l'attention de jeunes techniciens, pour les inviter à postuler avant le 30 novembre, par courrier ou courriel, dans un français formel. \fg{}
\end{methode}

\courssub{Les codes de la communication par internet}

Aujourd'hui, de nombreuses candidatures ne passent plus par le courrier papier mais par le \cle{courriel} (courrier électronique, ou \emph{e-mail}). Ce nouveau canal a ses propres codes, qu'il faut connaître pour ne pas être éliminé d'office.

Un courriel professionnel comprend :
\begin{itemize}
\item l'\textbf{adresse du destinataire} (champ \og À \fg{}) : elle doit être exacte et professionnelle ;
\item l'\textbf{objet du message} : une ligne brève et précise, par exemple \og Candidature au poste de technicien de maintenance \fg{}. Un courriel sans objet passe souvent pour du courrier indésirable ;
\item le \textbf{corps du message} : il reprend la structure de la lettre classique (formule d'appel, paragraphes, formule de politesse) ;
\item les \textbf{pièces jointes} : CV et diplômes scannés, nommés clairement (\emph{CV\_Ngo\_Bella.pdf} et non \emph{document1.pdf}) ;
\item la \textbf{signature} : nom, prénom, téléphone, adresse.
\end{itemize}

\begin{definition}[Netiquette]
La \cle{netiquette} (mot formé de \emph{net}, réseau, et \emph{étiquette}, règles de politesse) désigne l'ensemble des règles de bonne conduite et de courtoisie à respecter dans la communication par internet : politesse des formules, clarté de l'objet, réponse dans un délai raisonnable, respect de l'orthographe, interdiction du langage SMS et des majuscules agressives dans un cadre professionnel.
\end{definition}

\begin{attention}[Erreur fréquente : le langage SMS dans un courriel de candidature]
N'écrivez jamais \og slt \fg{}, \og bjr \fg{}, \og je c pa \fg{}, \og stp \fg{} ou \og mci \fg{} dans un courriel de candidature ou adressé à une administration. Ce langage, acceptable entre amis sur un téléphone, est perçu par un recruteur comme un manque de respect et de sérieux : il suffit souvent à éliminer un dossier. Écrivez en toutes lettres : \og Bonjour Monsieur, je vous remercie de votre réponse. \fg{} De même, évitez d'écrire TOUT EN MAJUSCULES, ce qui équivaut, sur internet, à crier.
\end{attention}

\courssub{Du courrier papier au courriel de candidature : ce qui change, ce qui demeure}

Le passage du papier à l'écran ne supprime pas les règles de la lettre officielle : il les déplace. Observons ce qui demeure et ce qui change.

\begin{popfigure}
\begin{center}
\begin{tikzpicture}[font=\small]
\node[draw=popblue, fill=popblueL, rounded corners, thick, minimum width=5.2cm, minimum height=0.9cm] (t1) at (0,0) {\textbf{Lettre papier}};
\node[draw=popgreen, fill=popgreenL, rounded corners, thick, minimum width=5.2cm, minimum height=0.9cm] (t2) at (8.2,0) {\textbf{Courriel professionnel}};
\node[draw=poporange, fill=poporangeL, rounded corners, thick, minimum width=4.2cm, minimum height=0.9cm] (t3) at (4.1,-2.1) {\textbf{Éléments communs}};
\node[align=left, text width=5.2cm] at (0,-2.3) {%
\textbullet{} en-tête rédigé (adresses)\\[2pt]
\textbullet{} lieu et date écrits\\[2pt]
\textbullet{} timbre et enveloppe\\[2pt]
\textbullet{} délai : plusieurs jours};
\node[align=left, text width=5.2cm] at (8.2,-2.3) {%
\textbullet{} adresses électroniques\\[2pt]
\textbullet{} date automatique\\[2pt]
\textbullet{} pièces jointes (PDF)\\[2pt]
\textbullet{} délai : quelques secondes};
\node[align=left, text width=4.6cm] at (4.1,-4.6) {%
\textbullet{} formule d'appel (\og Monsieur \fg{})\\[2pt]
\textbullet{} objet précis\\[2pt]
\textbullet{} corps structuré en paragraphes\\[2pt]
\textbullet{} formule de politesse finale\\[2pt]
\textbullet{} orthographe soignée};
\draw[-{Stealth[length=2.5mm]}, thick, popmuted] (t1.south) -- (t3.north west);
\draw[-{Stealth[length=2.5mm]}, thick, popmuted] (t2.south) -- (t3.north east);
\end{tikzpicture}
\end{center}
\legende{Comparaison entre la lettre papier et le courriel professionnel : le support et la rapidité changent, mais les formules de politesse, l'objet, la structure et le soin orthographique demeurent identiques.}
\end{popfigure}

\begin{aretenir}[Ce qui change, ce qui demeure]
\textbf{Ce qui change :} le support (écran au lieu du papier), la rapidité (quelques secondes), l'objet qui devient une ligne obligatoire, les documents qui deviennent des pièces jointes.\par
\textbf{Ce qui demeure :} la politesse des formules, la clarté de la structure, le respect de l'orthographe et de la grammaire, l'identification précise de l'émetteur et du destinataire. Le courriel professionnel est une \emph{lettre}, pas un message de conversation.
\end{aretenir}

\begin{savaistu}[La lettre manuscrite, toujours exigée au Cameroun]
Au Cameroun, malgré le développement d'internet, de nombreuses administrations exigent encore la \cle{lettre manuscrite} (écrite à la main) pour les grands concours de la fonction publique : ENS (École Normale Supérieure), ENAM (École Nationale d'Administration et de Magistrature), IRIC (Institut des Relations Internationales du Cameroun), concours de la police et de la gendarmerie. La lettre manuscrite permet notamment de vérifier l'écriture et l'authenticité du candidat. Savoir rédiger à la main une lettre de motivation propre et correcte reste donc une compétence décisive pour l'élève de Première technique.
\end{savaistu}

\courssub{Exercice d'application}

\begin{exercice}[Analyse d'une situation de communication]
Votre lycée reçoit l'affiche suivante : \og La Camtel recrute 10 techniciens en réseaux de télécommunication. Envoyez lettre de motivation et CV avant le 15 décembre à recrutement@camtel.cm. \fg{}
\begin{enumerate}
\item Identifiez l'émetteur, le destinataire, le message, l'objectif, le canal et le code de ce texte.
\item Rédigez l'objet du courriel qu'enverra un candidat.
\item Expliquez pourquoi il serait fautif d'écrire dans ce courriel : \og bjr, je ve le poste, mci \fg{}.
\end{enumerate}
\end{exercice}

\begin{corrige}[Exercice d'application]
\begin{enumerate}
\item \textbf{Émetteur :} l'entreprise Camtel. \textbf{Destinataire :} les jeunes techniciens en télécommunication (le public des diplômés). \textbf{Message :} une offre de 10 postes de techniciens en réseaux. \textbf{Objectif :} inciter les candidats à envoyer leur dossier avant le 15 décembre. \textbf{Canal :} l'affiche et le courriel (internet). \textbf{Code :} le français formel, concis, organisé en informations hiérarchisées.
\item Objet possible : \og Candidature au poste de technicien en réseaux de télécommunication \fg{}. L'objet est bref, précis et reprend l'intitulé exact du poste.
\item Ce message est fautif parce qu'il utilise le langage SMS (\og bjr \fg{}, \og ve \fg{}, \og mci \fg{}), interdit par la netiquette dans un cadre professionnel. Il manque de formule d'appel, de formule de politesse et de structure ; il donne une image négligée du candidat et peut entraîner l'élimination immédiate de sa candidature.
\end{enumerate}
\end{corrige}

\begin{aretenir}[Bilan de la section et transition]
Un texte fonctionnel vise une action concrète ; il se reconnaît à sa clarté, sa concision, sa mise en page aérée et ses informations hiérarchisées. Pour l'analyser, on identifie la situation de communication : émetteur, destinataire, message, objectif, canal, code. La communication par internet obéit à la netiquette : objet précis, formules de politesse, orthographe soignée, refus du langage SMS. Ces compétences préparent directement la rédaction de la lettre de motivation. \textbf{La prochaine étape consiste à maîtriser la structure rigoureuse de son en-tête}, première marque de sérieux aux yeux du recruteur.
\end{aretenir}

\coursec{L'en-tête de la lettre de motivation}

Dans la section précédente, nous avons vu que la communication par internet obéit à des codes précis : un courriel professionnel, par exemple, exige un objet clair, une formule d'appel et une signature. La lettre de motivation, document utilitaire par excellence, obéit à des codes encore plus stricts. Avant même d'écrire le corps de la lettre, le candidat doit soigner sa première partie : \cle{l'en-tête}. C'est elle que le recruteur lit en premier, et c'est souvent elle qui décide si la lettre sera lue jusqu'au bout ou écartée. Cette section étudie donc en détail la définition de la lettre de motivation, puis chacun des éléments qui composent son en-tête.

\courssub{Définition et fonction de la lettre de motivation}

\begin{definition}[La lettre de motivation]
La \cle{lettre de motivation} (ou lettre de candidature) est une lettre officielle, courte et structurée, qu'un candidat adresse à un recruteur (directeur, chef de service, responsable des ressources humaines) pour solliciter un emploi, un stage ou une formation. Sa fonction est de \textbf{convaincre} le destinataire de retenir la candidature : elle présente le candidat, montre l'adéquation entre son profil et le poste, et donne envie de le rencontrer en entretien.
\end{definition}

La lettre de motivation accompagne généralement le curriculum vitae (CV). Le CV énumère les diplômes et les expériences ; la lettre, elle, les \emph{met en valeur} et les relie au poste visé. On distingue deux situations :

\begin{itemize}
\item la \textbf{candidature en réponse à une offre} : le candidat répond à une annonce publiée (journal, affichage, internet) ; il doit alors citer la référence de l'offre ;
\item la \textbf{candidature spontanée} : le candidat propose ses services sans annonce préalable ; il doit alors cibler précisément l'entreprise et le poste souhaité.
\end{itemize}

\begin{savaistu}[Une minute pour convaincre]
Des études menées auprès de cabinets de recrutement montrent qu'un recruteur consacre en moyenne \textbf{moins d'une minute} à la première lecture d'une lettre de motivation. Son regard se porte d'abord sur l'en-tête : qui écrit ? à qui ? pour quel poste ? Si ces informations sont absentes, confuses ou fautives, la lettre est souvent écartée avant même la lecture du corps. Soigner l'en-tête, c'est donc déjà gagner la moitié du combat.
\end{savaistu}

\courssub{Observation : que contient un en-tête ?}

Avant d'énoncer les règles, observons un document réel. Voici l'en-tête d'une lettre adressée par un jeune technicien au Directeur des Ressources Humaines d'une entreprise de Douala :

\begin{exemplebox}[En-tête d'une lettre de candidature]
\begin{tabular}{p{0.45\linewidth}p{0.5\linewidth}}
\textbf{MBALLA Etienne} & \\
12, rue des Palmiers, B.P. 4021 Yaoundé & \\
Tél. : 690 12 34 56 & \\
Courriel : e.mballa@mail.com & \\
 & \textbf{Monsieur le Directeur}\\
 & \textbf{des Ressources Humaines}\\
 & Société CAMBAT S.A.\\
 & Douala\\[0.4em]
\multicolumn{2}{r}{Yaoundé, le 15 octobre 2026}\\[0.4em]
\multicolumn{2}{l}{\textbf{Objet : candidature au poste de technicien comptable}}\\
\multicolumn{2}{l}{Réf. : annonce n° 245 du 02/10/2026}\\[0.4em]
\multicolumn{2}{l}{Monsieur le Directeur,}\\
\end{tabular}
\end{exemplebox}

L'observation de ce document permet de dégager une règle générale.

\begin{exemplebox}[Les 6 éléments obligatoires de l'en-tête]
Un en-tête complet comporte \textbf{six éléments obligatoires}, toujours placés dans le même ordre :
\begin{enumerate}
\item les \textbf{coordonnées de l'expéditeur} (celui qui écrit), en haut à gauche ;
\item les \textbf{coordonnées du destinataire} (celui qui reçoit), en haut à droite ;
\item le \textbf{lieu et la date} de rédaction, à droite, sous le destinataire ;
\item l'\textbf{objet} de la lettre, en gras ;
\item les \textbf{références} de l'offre, le cas échéant ;
\item la \textbf{formule d'appel}, qui ouvre la lettre.
\end{enumerate}
\end{exemplebox}

\courssub{Les coordonnées de l'expéditeur}

L'\cle{expéditeur} est la personne qui écrit et envoie la lettre : ici, le candidat. Ses coordonnées se placent \textbf{en haut à gauche} de la page. Elles comprennent, dans l'ordre :

\begin{itemize}
\item le \textbf{nom} (souvent en majuscules) et le \textbf{prénom} ;
\item l'\textbf{adresse complète} (rue, boîte postale, ville) ;
\item le \textbf{numéro de téléphone} ;
\item le \textbf{courriel} (adresse électronique sérieuse, de préférence de la forme prénom.nom).
\end{itemize}

Ces informations permettent au recruteur de répondre ou de convoquer le candidat. Un numéro de téléphone faux ou un courriel fantaisiste (du type « petitprince237@... ») dessert la candidature.

\courssub{Les coordonnées du destinataire}

Le \cle{destinataire} est la personne à qui la lettre est adressée. Ses coordonnées se placent \textbf{en haut à droite}, légèrement plus bas que celles de l'expéditeur. On indique :

\begin{itemize}
\item la \textbf{qualité} : Monsieur, Madame ;
\item la \textbf{fonction} : le Directeur des Ressources Humaines, le Chef du Personnel, le Principal du lycée ;
\item le \textbf{nom de l'entreprise} ou de l'institution ;
\item la \textbf{ville} où elle se trouve.
\end{itemize}

On n'écrit jamais le nom personnel du destinataire dans l'en-tête d'une lettre officielle, même si on le connaît : c'est sa \emph{fonction} qui est visée, non sa personne.

\courssub{Le lieu et la date}

Sous les coordonnées du destinataire, à droite, on indique le \textbf{lieu} de rédaction suivi de la \textbf{date complète}, selon le modèle :

\begin{center}
\emph{Yaoundé, le 15 octobre 2026}
\end{center}

La date doit être \textbf{complète} : jour en chiffres, mois en lettres minuscules, année en entier. On écrit « le 15 octobre 2026 », et non « 15/10/26 » ni « octobre 2026 ». Le mois ne prend pas de majuscule en français.

\courssub{L'objet et les références}

L'\cle{objet} est une phrase nominale (sans verbe conjugué) qui résume en quelques mots la raison de la lettre. Il est précédé de la mention « Objet : » et se met \textbf{en gras}. Il doit être \textbf{bref} (une ligne au maximum) et \textbf{précis} : il nomme le poste ou la demande exacte.

\begin{methode}[Rédiger un objet précis et percutant]
\begin{enumerate}
\item Identifier la \textbf{nature} de la démarche : candidature, demande de stage, demande d'emploi.
\item Nommer \textbf{exactement} le poste ou la formation visés, avec les mots de l'annonce : « technicien comptable », « stage de maintenance industrielle ».
\item Assembler les deux en une phrase nominale courte : « Objet : candidature au poste de technicien comptable ».
\item Vérifier que l'objet tient sur \textbf{une seule ligne} et ne contient ni verbe conjugué, ni point final, ni détail superflu (le salaire, les motivations se discutent dans le corps).
\end{enumerate}
\end{methode}

Lorsque la candidature répond à une annonce, on ajoute sous l'objet les \cle{références} de l'offre, précédées de « Réf. : ». Cela permet au recruteur de classer immédiatement la lettre dans le bon dossier :

\begin{center}
\emph{Réf. : annonce n° 245 du 02/10/2026}
\end{center}

\courssub{La formule d'appel}

La \cle{formule d'appel} est la première ligne du texte proprement dit : elle s'adresse directement au destinataire. Elle reprend sa qualité et, éventuellement, sa fonction :

\begin{itemize}
\item \emph{Monsieur le Directeur,}
\item \emph{Madame,}
\item \emph{Monsieur le Principal,}
\end{itemize}

Deux règles sont absolues : on ne connaît jamais le destinataire par son nom (« Monsieur Ngo Bell » est interdit dans l'appel), et on ne le tutoie ni ne le familiarise jamais (« Cher Monsieur » est à proscrire dans une lettre de candidature).

\begin{attention}[Erreur fréquente : la redondance dans l'appel]
Il est fautif d'écrire : « \emph{Monsieur le Directeur Général de la société X, Monsieur,} ». Cette formule double inutilement la qualité : la fonction complète suffit. On écrira simplement « \emph{Monsieur le Directeur Général,} ». Autres erreurs à éviter : « Monsieur le Directeur de la société X, Monsieur le Directeur, » (répétition), « Cher Directeur, » (trop familier), ou « Monsieur Etoua, » (nom personnel interdit dans l'appel).
\end{attention}

\begin{attention}[Les trois erreurs classiques de l'en-tête]
\begin{itemize}
\item \textbf{Inverser les coordonnées} : l'expéditeur à gauche, le destinataire à droite — jamais l'inverse ;
\item \textbf{Dater de façon incomplète} : « le 15/10 » ou « octobre 2026 » ne sont pas acceptables ;
\item \textbf{Formuler un objet vague} : « Objet : demande » ou « Objet : lettre de motivation » ne renseignent pas le recruteur ; l'objet doit nommer le poste précis.
\end{itemize}
\end{attention}

\courssub{Schéma récapitulatif : la maquette de l'en-tête}

\begin{popfigure}
\begin{tikzpicture}
% Page
\draw[thick, popdark, fill=popcream] (0,0) rectangle (9,11);
% Expéditeur
\node[anchor=north west, text width=4.2cm, align=left, font=\small] at (0.4,10.6)
{\textbf{MBALLA Etienne}\\12, rue des Palmiers\\B.P. 4021 Yaoundé\\Tél. : 690 12 34 56\\e.mballa@mail.com};
% Destinataire
\node[anchor=north west, text width=4.2cm, align=left, font=\small] at (5.0,8.9)
{\textbf{Monsieur le Directeur}\\\textbf{des Ressources Humaines}\\Société CAMBAT S.A.\\Douala};
% Lieu et date
\node[anchor=north west, font=\small] at (5.0,6.6) {Yaoundé, le 15 octobre 2026};
% Objet
\node[anchor=north west, font=\small] at (0.4,5.7) {\textbf{Objet : candidature au poste de technicien comptable}};
% Référence
\node[anchor=north west, font=\small] at (0.4,5.1) {Réf. : annonce n° 245 du 02/10/2026};
% Appel
\node[anchor=north west, font=\small] at (0.4,4.3) {Monsieur le Directeur,};
% Corps suggéré
\foreach \y in {3.5,3.1,2.7,2.3,1.9,1.5} {\draw[popline] (0.4,\y) -- (8.6,\y);}
% Flèches d'annotation
\draw[-{Stealth}, thick, popblue] (2.2,9.6) -- (2.2,11.6) node[above, font=\footnotesize\bfseries, popblue]{1. Expéditeur (gauche)};
\draw[-{Stealth}, thick, poporange] (7.0,8.2) -- (7.0,11.0) node[above, font=\footnotesize\bfseries, poporange, align=center]{2. Destinataire\\(droite)};
\draw[-{Stealth}, thick, popgreen] (7.0,6.5) -- (10.3,6.5) node[right, font=\footnotesize\bfseries, popgreen]{3. Lieu et date};
\draw[-{Stealth}, thick, poppurple] (4.6,5.55) -- (10.3,5.55) node[right, font=\footnotesize\bfseries, poppurple]{4. Objet (en gras)};
\draw[-{Stealth}, thick, poppink] (4.2,4.95) -- (10.3,4.95) node[right, font=\footnotesize\bfseries, poppink]{5. Références};
\draw[-{Stealth}, thick, popgold] (2.6,4.15) -- (10.3,4.15) node[right, font=\footnotesize\bfseries, popgold]{6. Formule d'appel};
\end{tikzpicture}
\legende{Maquette annotée de l'en-tête d'une lettre de motivation. Les flèches numérotées désignent les six éléments obligatoires, dans leur ordre et leur position sur la page.}
\end{popfigure}

\courssub{Exercice résolu et entraînement}

\begin{exoresolu}[Corriger un en-tête fautif]
\textbf{Énoncé.} Voici l'en-tête rédigé par un élève de Première technique. Relevez les erreurs et réécrivez l'en-tête correctement.

\smallskip
\emph{Monsieur le Directeur de la SODECOTON, Douala.\\
Objet : lettre\\
Douala, 03/11/26\\
NGOUMOU Larissa, B.P. 55 Garoua, tél. 699 87 65 43, larissa.ngoumou@mail.com\\
Monsieur le Directeur de la SODECOTON, Monsieur,}
\tcblower
\textbf{Solution détaillée.} On relève quatre erreurs :
\begin{enumerate}
\item \textbf{Inversion des coordonnées} : le destinataire (la SODECOTON) est placé en premier, à la place de l'expéditrice. Or l'expéditeur doit figurer en haut à gauche, le destinataire en haut à droite.
\item \textbf{Objet vague} : « lettre » ne précise ni la nature de la démarche ni le poste visé.
\item \textbf{Date incomplète} : « 03/11/26 » doit s'écrire en toutes lettres : « le 3 novembre 2026 ».
\item \textbf{Redondance dans l'appel} : « Monsieur le Directeur de la SODECOTON, Monsieur, » répète la qualité.
\end{enumerate}
En-tête corrigé :

\smallskip
\begin{tabular}{p{0.45\linewidth}p{0.5\linewidth}}
\textbf{NGOUMOU Larissa} & \\
B.P. 55 Garoua & \\
Tél. : 699 87 65 43 & \\
larissa.ngoumou@mail.com & \\
 & \textbf{Monsieur le Directeur}\\
 & \textbf{de la SODECOTON}\\
 & Douala\\[0.3em]
\multicolumn{2}{r}{Garoua, le 3 novembre 2026}\\[0.3em]
\multicolumn{2}{l}{\textbf{Objet : candidature au poste d'agent de maintenance}}\\[0.3em]
\multicolumn{2}{l}{Monsieur le Directeur,}\\
\end{tabular}
\end{exoresolu}

\begin{exercice}[À vous de rédiger]
Vous êtes ESSOMBA Jean, demeurant au quartier Nkolbisson, B.P. 1200 Yaoundé (tél. 677 45 67 89, courriel : j.essomba@mail.com). Vous répondez à l'annonce n° 88 du 10 janvier 2027 publiée par l'entreprise TECHNIPLAST de Douala, qui recherche un technicien en électromécanique. Rédigez l'en-tête complet de votre lettre de motivation, datée du 20 janvier 2027, en respectant les six éléments obligatoires et leur disposition sur la page.
\end{exercice}

\begin{corrige}[Exercice : rédaction de l'en-tête]
\begin{tabular}{p{0.45\linewidth}p{0.5\linewidth}}
\textbf{ESSOMBA Jean} & \\
Quartier Nkolbisson & \\
B.P. 1200 Yaoundé & \\
Tél. : 677 45 67 89 & \\
j.essomba@mail.com & \\
 & \textbf{Monsieur le Directeur}\\
 & \textbf{des Ressources Humaines}\\
 & Entreprise TECHNIPLAST\\
 & Douala\\[0.3em]
\multicolumn{2}{r}{Yaoundé, le 20 janvier 2027}\\[0.3em]
\multicolumn{2}{l}{\textbf{Objet : candidature au poste de technicien en électromécanique}}\\
\multicolumn{2}{l}{Réf. : annonce n° 88 du 10/01/2027}\\[0.3em]
\multicolumn{2}{l}{Monsieur le Directeur,}\\
\end{tabular}

\smallskip
Points de vigilance : expéditeur à gauche, destinataire à droite ; date complète avec le mois en minuscules ; objet précis reprenant l'intitulé exact de l'annonce ; référence citée ; appel sans nom personnel ni redondance.
\end{corrige}

\begin{aretenir}[L'essentiel de la section]
\begin{itemize}
\item La lettre de motivation vise à \textbf{convaincre} un recruteur de retenir une candidature ; l'en-tête est lu en premier et doit être irréprochable.
\item L'en-tête comporte \textbf{six éléments obligatoires} : expéditeur (gauche), destinataire (droite), lieu et date complets, objet bref et précis en gras, références de l'offre, formule d'appel.
\item L'objet est une phrase nominale d'une ligne qui nomme exactement le poste visé.
\item L'appel vise la fonction, jamais le nom : « Monsieur le Directeur, » — sans redondance ni familiarité.
\end{itemize}
\end{aretenir}

L'en-tête ainsi construit, la lettre peut s'ouvrir. Mais pour rédiger le corps de cette lettre avec des phrases variées et un vocabulaire riche, encore faut-il maîtriser deux outils de langue : c'est l'objet de la section suivante, consacrée à la phrase complexe et à l'emprunt.

\coursec{La phrase complexe et l'emprunt}

Dans la section précédente, nous avons appris à construire l'en-tête d'une lettre de motivation : expéditeur, destinataire, lieu et date, objet. Mais une lettre ne se juge pas seulement à sa présentation : elle se juge aussi à la \emph{qualité de son écriture}. Deux candidats peuvent avoir le même diplôme ; celui qui sait \cle{lier ses idées avec élégance} et choisir un \cle{vocabulaire soigné} obtiendra l'entretien. C'est précisément l'objet de cette section : d'une part la \cle{phrase complexe}, outil grammatical qui permet de relier plusieurs idées dans une même phrase ; d'autre part l'\cle{emprunt}, phénomène lexical qu'il faut connaître pour éviter les mots étrangers malvenus dans une lettre officielle.

\courssub{De la phrase simple à la phrase complexe}

\textbf{Observation de départ.} Lisons ces deux versions d'un même passage de lettre de motivation :

\begin{itemize}
\item \emph{Version A :} « J'ai obtenu mon diplôme. J'ai fait un stage. Le stage a duré six mois. Je suis disponible. »
\item \emph{Version B :} « Titulaire d'un diplôme de technicien, j'ai effectué un stage de six mois, ce qui me rend immédiatement disponible. »
\end{itemize}

Les deux versions disent la même chose. Pourtant, la version A ressemble à une liste d'enfant, tandis que la version B donne l'image d'un candidat maîtrisant sa langue. La différence tient à un seul fait : la version A empile quatre \cle{phrases simples}, la version B construit une \cle{phrase complexe}.

\begin{definition}[Phrase simple et phrase complexe]
Une \cle{phrase simple} ne contient qu'\textbf{une seule proposition}, c'est-à-dire un seul verbe conjugué (ou un seul verbe à un mode personnel) autour duquel s'organisent les groupes de mots. Une \cle{phrase complexe} contient \textbf{au moins deux propositions}, donc au moins deux verbes conjugués, associés par différents procédés.
\end{definition}

\begin{exemplebox}[Compter les verbes pour compter les propositions]
\begin{itemize}
\item « Je vous adresse ma candidature. » : un seul verbe (\emph{adresse}) $\rightarrow$ phrase \textbf{simple}.
\item « Je vous adresse ma candidature car votre entreprise me passionne. » : deux verbes (\emph{adresse}, \emph{passionne}) $\rightarrow$ phrase \textbf{complexe}.
\end{itemize}
\textbf{Réflexe :} pour analyser une phrase, soulignez d'abord tous les verbes conjugués ; leur nombre donne le nombre de propositions.
\end{exemplebox}

\courssub{Les trois procédés d'association des propositions}

Lorsqu'une phrase contient plusieurs propositions, celles-ci peuvent être reliées de trois manières.

\begin{definition}[Juxtaposition, coordination, subordination]
\begin{itemize}
\item La \cle{juxtaposition} : les propositions sont placées côte à côte, séparées par une virgule ou un point-virgule, \textbf{sans mot de liaison}. \emph{Ex. : « Je suis ponctuel, je suis rigoureux. »}
\item La \cle{coordination} : les propositions sont reliées par une \textbf{conjonction de coordination} (\emph{mais, ou, et, donc, or, ni, car}). Elles restent sur le même plan. \emph{Ex. : « Je suis jeune, mais je suis expérimenté. »}
\item La \cle{subordination} : une proposition (la \textbf{subordonnée}) dépend d'une autre (la \textbf{principale}) et ne peut pas exister seule. Elle est introduite par un mot subordonnant (pronom relatif, conjonction de subordination) ou mise à un mode impersonnel (participe, infinitif).
\end{itemize}
\end{definition}

\begin{attention}[Ne pas confondre juxtaposition et coordination]
Beaucoup d'élèves appellent « coordination » toute suite de propositions. Le critère est simple : \textbf{y a-t-il une conjonction de coordination ?} Dans « Il pleut, je reste », il n'y a aucun mot de liaison : c'est une \textbf{juxtaposition}. Dans « Il pleut, donc je reste », la conjonction \emph{donc} fait de la seconde proposition une \textbf{coordonnée}. Autre piège fréquent : employer le mot anglais \emph{job} dans une lettre officielle ; nous verrons plus bas pourquoi il faut lui préférer \emph{emploi} ou \emph{poste}.
\end{attention}

\courssub{La subordination : les trois types de subordonnées}

\begin{definition}[Subordination]
La \cle{subordination} est le procédé par lequel une proposition, dite \textbf{subordonnée}, est intégrée à une proposition \textbf{principale} dont elle dépend grammaticalement. La subordonnée occupe une fonction dans la principale (épithète, complément, circonstancielle), exactement comme le ferait un simple groupe de mots.
\end{definition}

\textbf{1. La subordonnée relative.} Elle est introduite par un \cle{pronom relatif} (\emph{qui, que, où, dont}) et complète un nom, appelé \emph{antécédent}.

\begin{exemplebox}[La relative dans la lettre de motivation]
« Je suis très intéressé par le poste \textbf{que} vous proposez. »
\begin{itemize}
\item Proposition principale : « Je suis très intéressé par le poste. »
\item Subordonnée relative : « que vous proposez » ; \emph{que} est pronom relatif, COD de \emph{proposez}, antécédent : \emph{le poste}.
\end{itemize}
Autres exemples utiles : « l'entreprise \textbf{qui} recrute » ; « la ville \textbf{où} je réside » ; « le diplôme \textbf{dont} je suis titulaire ».
\end{exemplebox}

\textbf{2. La subordonnée conjonctive.} Elle est introduite par une \cle{conjonction de subordination} (\emph{que, parce que, si, bien que, afin que, lorsque\ldots}). On distingue :
\begin{itemize}
\item les \textbf{complétives}, qui complètent un verbe : « Je pense \textbf{que} mon profil correspond à vos attentes » ;
\item les \textbf{circonstancielles}, qui expriment une circonstance :
\begin{itemize}
\item but : « Je me tiens à votre disposition \textbf{afin que} nous puissions nous rencontrer » ;
\item cause : « Je postule \textbf{parce que} votre entreprise est un leader du secteur » ;
\item condition : « \textbf{Si} vous retenez ma candidature, je m'engagerai pleinement » ;
\item concession : « \textbf{Bien que} jeune diplômé, je possède une solide expérience de stage ».
\end{itemize}
\end{itemize}

\textbf{3. La subordonnée participiale.} Elle est construite avec un \cle{participe présent} (en \emph{-ant}) ou un \cle{participe passé}, sans conjonction. Elle est très appréciée dans les lettres officielles car elle est brève et élégante.

\begin{exemplebox}[La participiale, marque de style soutenu]
« \textbf{Titulaire d'un diplôme de technicien}, je maîtrise la maintenance industrielle. » La proposition participiale « titulaire d'un diplôme de technicien » renseigne sur le sujet \emph{je} sans alourdir la phrase. Variante au participe présent : « \textbf{Ayant} effectué plusieurs stages, je connais le milieu de l'entreprise. »
\end{exemplebox}

\begin{popfigure}
\begin{tikzpicture}[
  grow=down,
  level distance=1.6cm,
  level 1/.style={sibling distance=5.2cm},
  level 2/.style={sibling distance=2.6cm},
  every node/.style={draw=popblue, rounded corners=2pt, fill=popblueL, text width=3.1cm, align=center, font=\small, inner sep=3pt},
  edge from parent/.style={draw=popdark, thick}
]
\node[fill=poporangeL, draw=poporange, text width=4.5cm]{\textbf{Propositions subordonnées}}
  child { node[align=center]{\textbf{Relative}\\ introduite par \emph{qui, que, où, dont}}
    child { node[fill=popgreenL, draw=popgreen, text width=4.2cm]{« le poste \emph{que} vous proposez »}} }
  child { node[align=center]{\textbf{Conjonctive}\\ introduite par \emph{que, parce que, si, bien que, afin que}}
    child { node[fill=popgreenL, draw=popgreen, text width=4.2cm]{« Je pense \emph{que} mon profil convient »}}
    child { node[fill=popgreenL, draw=popgreen, text width=4.2cm]{« \emph{Bien que} jeune, je suis expérimenté »}} }
  child { node[align=center]{\textbf{Participiale}\\ participe présent ou passé, sans conjonction}
    child { node[fill=popgreenL, draw=popgreen, text width=4.2cm]{« \emph{Titulaire} d'un diplôme, je\ldots »}} };
\end{tikzpicture}
\legende{Classification des trois types de propositions subordonnées, avec un exemple représentatif pour chacune, tiré du contexte de la lettre de motivation.}
\end{popfigure}

\courssub{Méthode : transformer des phrases simples en une phrase complexe}

\begin{methode}[Fusionner deux phrases simples en une phrase complexe]
\begin{enumerate}
\item \textbf{Repérer} l'élément commun aux deux phrases (un nom répété, un lien logique de cause, de but, de condition).
\item \textbf{Choisir} le procédé adapté : pronom relatif si un nom est répété ; conjonction de subordination s'il existe un rapport logique ; participe si la seconde phrase décrit le sujet de la première.
\item \textbf{Réécrire} en supprimant la répétition : le mot répété est remplacé par le pronom relatif, ou la seconde phrase devient participiale.
\item \textbf{Vérifier} que la phrase obtenue est correcte (accord du participe, ponctuation) et qu'elle garde exactement le même sens.
\end{enumerate}
\end{methode}

\begin{exoresolu}[Trois phrases maladroites fusionnées en une phrase élégante]
Énoncé. Un candidat a écrit dans sa lettre : « J'ai obtenu le Probatoire technique. J'ai ensuite préparé un diplôme de technicien. Ce diplôme m'a permis de faire un stage chez SODECOTON. » Transformez ces trois phrases simples en une seule phrase complexe.
\tcblower
Solution détaillée.
\begin{enumerate}
\item Éléments communs : \emph{diplôme} est répété (phrases 2 et 3) $\rightarrow$ subordonnée relative avec \emph{qui} ou \emph{que}. La phrase 1 est une étape préalable $\rightarrow$ proposition participiale avec \emph{ayant} + participe passé.
\item Choix des procédés : participiale pour la première idée (« ayant obtenu\ldots »), relative pour la seconde (« qui m'a permis\ldots »).
\item Réécriture : « \textbf{Ayant obtenu le Probatoire technique, j'ai préparé un diplôme de technicien qui m'a permis d'effectuer un stage chez SODECOTON.} »
\item Vérification : deux verbes conjugués (\emph{ai préparé}, \emph{a permis}) et un participe (\emph{ayant obtenu}) ; le sens est identique ; le style est nettement plus soutenu.
\end{enumerate}
Bilan chiffré : \textbf{3 phrases simples} (3 verbes isolés, 2 répétitions) sont devenues \textbf{1 phrase complexe} à 3 propositions, sans aucune perte d'information.
\end{exoresolu}

\begin{aretenir}[Pourquoi la phrase complexe est un atout dans la lettre de motivation]
La phrase complexe permet de \textbf{lier les arguments} au lieu de les empiler : elle montre les rapports logiques (cause, but, concession) entre les idées. Une lettre bien écrite alterne phrases simples (pour les informations brèves et fermes) et phrases complexes (pour développer et nuancer). Attention toutefois : une phrase de plus de trois lignes devient illisible ; l'élégance n'exclut pas la clarté.
\end{aretenir}

\courssub{L'emprunt : définition et origines}

\textbf{Observation de départ.} Dans la cour de récréation, on entend : « J'ai un \emph{job} le week-end, mon \emph{boss} est cool, on fait un \emph{briefing} le matin. » Ces mots ne sont pas français d'origine : ils viennent de l'anglais. Le français, comme toutes les langues vivantes, \emph{emprunte} des mots aux autres langues.

\begin{definition}[L'emprunt (lexicologie)]
Un \cle{emprunt} est un mot qu'une langue prend à une autre langue et intègre à son vocabulaire, en conservant ou en adaptant sa forme et sa prononciation. L'emprunt est un des principaux \textbf{procédés d'enrichissement du lexique}, avec la dérivation et la composition.
\end{definition}

Le français a emprunté à de nombreuses langues :
\begin{itemize}
\item à l'\textbf{anglais} : \emph{week-end, parking, planning, challenge, interview} ;
\item à l'\textbf{arabe} : \emph{algèbre, chiffre, sucre, café, douane} ;
\item aux \textbf{langues nationales du Cameroun} : des mots comme \emph{ndolé} (plat traditionnel), \emph{mbol} (en bafia), \emph{ashia} (marque de compassion, en pidgin), qui enrichissent le français parlé au Cameroun ;
\item à l'\textbf{italien} : \emph{piano, opera, spaghetti} ; à l'\textbf{espagnol} : \emph{cigare, embargo}.
\end{itemize}

\begin{savaistu}[Le français du Cameroun, une richesse]
Le français parlé au Cameroun intègre naturellement des emprunts aux langues nationales (ewondo, douala, fulfuldé, pidgin\ldots). C'est une richesse culturelle dans la conversation quotidienne. Mais dans une \textbf{lettre officielle}, on exige le français standard : on écrira « je vous serais reconnaissant de\ldots » et non une formule empruntée à une langue nationale, aussi expressive soit-elle.
\end{savaistu}

\courssub{Emprunts intégrés et emprunts à éviter dans une lettre officielle}

Tous les emprunts ne se valent pas. Il faut distinguer :
\begin{itemize}
\item les \textbf{emprunts intégrés} : anciens, admis par les dictionnaires, souvent sans équivalent français courant. Ils sont acceptables, voire inévitables : \emph{planning} (pour « calendrier prévisionnel »), \emph{challenge} (pour « défi »), \emph{parking}, \emph{interview} ;
\item les \textbf{emprunts à éviter} : récents, familiers ou snobs, alors qu'un mot français simple existe. Dans une lettre officielle, ils donnent une image de négligence : \emph{job} (dire \emph{emploi, poste}), \emph{staff} (dire \emph{personnel, équipe}), \emph{boss} (dire \emph{chef, supérieur hiérarchique}), \emph{business} (dire \emph{affaires, entreprise}).
\end{itemize}

\begin{popfigure}
\begin{tikzpicture}
\node[draw=popblue, fill=popblueL, rounded corners=3pt, inner sep=6pt, text width=0.94\linewidth, font=\small] {%
\begin{tabular}{@{}p{0.2\linewidth}p{0.18\linewidth}p{0.3\linewidth}p{0.2\linewidth}@{}}
\textbf{Emprunt} & \textbf{Origine} & \textbf{Équivalent recommandé} & \textbf{Statut en lettre officielle} \\[2pt]
\hline
\rule{0pt}{10pt}\emph{planning} & anglais & calendrier, programme & toléré (intégré) \\
\emph{challenge} & anglais & défi & toléré (intégré) \\
\emph{interview} & anglais & entretien (d'embauche) & toléré, préférer \emph{entretien} \\
\emph{job} & anglais & emploi, poste & \textbf{à éviter} \\
\emph{staff} & anglais & personnel, équipe & \textbf{à éviter} \\
\emph{boss} & anglais & chef, supérieur & \textbf{à éviter} \\
\emph{sucre} & arabe & --- (pas d'équivalent) & intégré depuis longtemps \\
\emph{ndolé} & langue nationale & --- (mets camerounais) & hors lettre officielle \\
\end{tabular}};
\end{tikzpicture}
\legende{Emprunts courants, leur langue d'origine, l'équivalent recommandé en français soigné et leur statut dans une lettre officielle.}
\end{popfigure}

\begin{attention}[Le piège de l'anglicisme familier]
Écrire « Je cherche un \emph{job} dans votre \emph{staff} » dans une lettre de motivation est une faute de registre : le mot existe, mais il appartient à la langue familière. Le correcteur du Probatoire ou du Baccalauréat technique, comme un employeur, y verra un manque de maîtrise du registre soutenu. Écrivez : « Je souhaite occuper un \emph{emploi} au sein de votre \emph{personnel} » ou, mieux : « Je souhaite intégrer votre \emph{équipe} ».
\end{attention}

\courssub{Entraînement au type d'exercice du Bac : la réécriture}

Au Probatoire et au Baccalauréat technique, un exercice fréquent consiste à \textbf{transformer des phrases simples en une phrase complexe} en respectant une consigne précise (utiliser une relative, une conjonctive, une participiale).

\begin{exercice}[Réécritures]
\begin{enumerate}
\item Transformez en une phrase complexe avec une subordonnée relative : « Votre entreprise recrute des techniciens. Votre entreprise est réputée dans toute la région. »
\item Transformez avec une subordonnée circonstancielle de cause : « Je vous écris. J'ai lu votre annonce dans \emph{Cameroon Tribune}. »
\item Transformez avec une proposition participiale : « Je possède cinq ans d'expérience. Je peux occuper immédiatement le poste. »
\item Corrigez le registre : « Je veux faire partie de votre staff et avoir un job stable. »
\end{enumerate}
\end{exercice}

\begin{corrige}[Exercice de réécriture]
\begin{enumerate}
\item « Votre entreprise, \textbf{qui} est réputée dans toute la région, recrute des techniciens. » (relative introduite par \emph{qui}, antécédent \emph{entreprise}).
\item « Je vous écris \textbf{parce que} j'ai lu votre annonce dans \emph{Cameroon Tribune}. » (subordonnée conjonctive circonstancielle de cause).
\item « \textbf{Possédant} cinq ans d'expérience, je peux occuper immédiatement le poste. » (proposition participiale au participe présent ; on accepte aussi « Fort de cinq ans d'expérience\ldots »).
\item « Je souhaite intégrer votre \textbf{équipe} et occuper un \textbf{emploi} stable. » (\emph{staff} et \emph{job}, anglicismes familiers, sont remplacés par leurs équivalents soutenus).
\end{enumerate}
\end{corrige}

\begin{aretenir}[L'essentiel de la section]
\begin{itemize}
\item Une phrase \textbf{simple} contient une seule proposition ; une phrase \textbf{complexe} en contient au moins deux.
\item Trois procédés d'association : \textbf{juxtaposition} (sans mot de liaison), \textbf{coordination} (\emph{mais, ou, et, donc, or, ni, car}), \textbf{subordination} (relative, conjonctive, participiale).
\item La phrase complexe lie les arguments avec élégance : c'est l'outil stylistique majeur de la lettre de motivation.
\item Un \textbf{emprunt} est un mot pris à une autre langue ; dans une lettre officielle, on utilise les emprunts intégrés (\emph{planning, challenge}) et on bannit les anglicismes familiers (\emph{job, staff, boss}).
\end{itemize}
Nous sommes désormais prêts, dans la section suivante, à construire le \textbf{corps de la lettre} en organisant ces phrases complexes en paragraphes argumentés, et à y jouer des tonalités adaptées.
\end{aretenir}

\coursec{Le corps de la lettre et les tonalités}

Dans la section précédente, nous avons appris à construire des phrases complexes et à reconnaître les emprunts, deux outils qui donnent à notre écriture souplesse et précision. Nous allons maintenant mettre ces outils au service du cœur de la lettre de motivation : son \cle{corps}. C'est là que tout se joue. Un recruteur lit souvent une lettre en moins d'une minute : si le corps est mal organisé, la candidature est écartée. Nous verrons ensuite les \cle{tonalités} (dramatique et comique), notions de stylistique utiles pour lire les textes de réception, mais dont nous apprendrons aussi pourquoi elles sont à proscrire dans une lettre officielle.

\courssub{La structure du corps : le plan « vous – moi – nous »}

Imaginez un entretien d'embauche. Le candidat qui ne parle que de lui pendant vingt minutes fatigue son interlocuteur ; celui qui ne parle que de l'entreprise paraît flatteur et vide. Le bon candidat suit un mouvement naturel en trois temps : il montre d'abord qu'il connaît l'entreprise, puis il présente ce qu'il peut lui apporter, enfin il propose une rencontre. Le corps de la lettre de motivation reproduit exactement ce mouvement, en trois paragraphes.

\begin{definition}[Le plan « vous – moi – nous »]
Le corps d'une lettre de motivation se construit en \textbf{trois paragraphes} :
\begin{itemize}
\item le paragraphe du \cle{vous} : on parle de l'\textbf{entreprise} et du poste (ses activités, ses valeurs, ce qui nous attire chez elle) ;
\item le paragraphe du \cle{moi} : on parle de \textbf{soi} (formation, diplômes, stages, compétences, qualités) en lien avec le poste ;
\item le paragraphe du \cle{nous} : on évoque la \textbf{collaboration} possible, on exprime sa disponibilité et l'on sollicite un entretien.
\end{itemize}
\end{definition}

\begin{popfigure}
\begin{center}
\begin{tikzpicture}[font=\small]
\draw[fill=popblueL, draw=popblue, rounded corners, thick] (0,0) rectangle (3.6,4.2);
\node[font=\bfseries\large, text=popblue] at (1.8,3.7) {VOUS};
\node[align=left, text width=3.1cm] at (1.8,2)
  {$\bullet$ l'entreprise\\[2pt] $\bullet$ ses activités\\[2pt] $\bullet$ ses valeurs\\[2pt] $\bullet$ le poste visé};
\draw[fill=poporangeL, draw=poporange, rounded corners, thick] (4.1,0) rectangle (7.7,4.2);
\node[font=\bfseries\large, text=poporange] at (5.9,3.7) {MOI};
\node[align=left, text width=3.1cm] at (5.9,2)
  {$\bullet$ ma formation\\[2pt] $\bullet$ mes diplômes\\[2pt] $\bullet$ mes stages\\[2pt] $\bullet$ mes qualités};
\draw[fill=popgreenL, draw=popgreen, rounded corners, thick] (8.2,0) rectangle (11.8,4.2);
\node[font=\bfseries\large, text=popgreen] at (10,3.7) {NOUS};
\node[align=left, text width=3.1cm] at (10,2)
  {$\bullet$ la collaboration\\[2pt] $\bullet$ ma disponibilité\\[2pt] $\bullet$ la demande\\ d'entretien};
\draw[-{Stealth[length=3mm]}, very thick, popdark] (3.7,2.1) -- (4.0,2.1);
\draw[-{Stealth[length=3mm]}, very thick, popdark] (7.8,2.1) -- (8.1,2.1);
\end{tikzpicture}
\end{center}
\legende{Le corps de la lettre de motivation : trois paragraphes ordonnés, de l'entreprise vers la rencontre.}
\end{popfigure}

\begin{savaistu}[Une norme internationale]
Le plan « vous – moi – nous » n'est pas une invention scolaire camerounaise : c'est une \textbf{norme internationale} de la lettre de motivation, enseignée dans les écoles de commerce et les services de ressources humaines du monde entier (on parle parfois du modèle anglo-saxon \emph{you -- me -- us}). Un recruteur de Douala, de Paris ou de Montréal attend la même logique. Le maîtriser, c'est parler la langue des employeurs partout.
\end{savaistu}

\courssub{Le paragraphe du « vous » : prouver qu'on s'est informé}

Le premier paragraphe répond à une question que le recruteur se pose toujours : \emph{pourquoi nous ?} Une lettre qui commence par « Je vous écris parce que je cherche du travail » est faible : elle parle déjà de soi. Une lettre forte commence par l'entreprise.

\begin{methode}[Construire le paragraphe du « vous »]
\begin{enumerate}
\item \textbf{Se renseigner} avant d'écrire : site internet de l'entreprise, pages Facebook, articles de presse, bouche-à-oreille, visite des lieux.
\item \textbf{Choisir deux ou trois informations précises} : secteur d'activité, produits ou services, valeurs affichées, implantation au Cameroun (ville, quartier, filiales).
\item \textbf{Relier ces informations au poste} : montrer que l'on a compris ce que l'entreprise attend du candidat.
\item \textbf{Rédiger en deux à quatre phrases}, sans flatterie excessive.
\end{enumerate}
\end{methode}

\begin{exemplebox}[Deux débuts comparés]
\textbf{Faible :} « Je viens par la présente solliciter un emploi dans votre entreprise. » (Aucune information : la lettre pourrait être envoyée à cent entreprises.)

\textbf{Fort :} « Leader de l'hôtellerie sur le littoral camerounais, votre établissement accueille chaque année des milliers de touristes à Kribi et mise sur le confort de ses clients, notamment grâce à des installations climatisées. C'est pour contribuer à cette exigence de qualité que je me permets de vous proposer mes services. » (L'entreprise se reconnaît : la lettre lui est destinée.)
\end{exemplebox}

\courssub{Le paragraphe du « moi » : argumenter sans mentir}

C'est le paragraphe le plus délicat. Il faut y présenter sa \textbf{formation} (Baccalauréat technique, Brevet de Technicien, DUT), ses \textbf{expériences} (stages en entreprise, travaux pratiques, petits boulots) et ses \textbf{qualités} (ponctualité, rigueur, sens du travail en équipe). Deux écueils menacent : la \emph{prétention} (« je suis le meilleur technicien de ma génération ») et la \emph{fausse modestie} (« je ne sais rien faire, mais je veux bien essayer »).

\begin{methode}[Écrire le paragraphe du « moi » avec justesse]
\begin{enumerate}
\item \textbf{Énoncer un fait vérifiable} : diplôme obtenu, stage effectué (lieu, durée), tâche réalisée.
\item \textbf{En tirer la compétence} : « ce stage m'a permis de maîtriser l'entretien des groupes froids ».
\item \textbf{Relier au poste} : « compétence directement utile pour votre service technique ».
\item \textbf{Ne citer que des qualités illustrées} : une qualité sans preuve ne vaut rien. « Rigoureux » ne convainc que si l'on ajoute « comme en témoigne la tenue de mon carnet de maintenance pendant mon stage ».
\item \textbf{Ne jamais mentir} : tout ce qui est écrit peut être vérifié à l'entretien (dates, références, gestes techniques).
\end{enumerate}
\end{methode}

\begin{attention}[Erreur fréquente : recopier son CV]
Beaucoup d'élèves transforment le paragraphe du « moi » en liste sèche : « 2023 : BEPC ; 2025 : Probatoire ; 2026 : stage à la SODECAO ; 2027 : Bac technique. » C'est recopier son CV, et c'est une erreur. Le CV \textbf{énumère}, la lettre \textbf{argumente} : elle choisit deux ou trois éléments du CV, les commente et montre en quoi ils répondent aux besoins du poste. Une lettre qui répète le CV ne sert à rien ; une lettre qui l'explique convainc.
\end{attention}

\courssub{Le paragraphe du « nous » : vers l'entretien}

Le dernier paragraphe projette l'avenir : l'entreprise et le candidat travaillant ensemble. On y exprime sa \textbf{disponibilité} (« je suis disponible immédiatement », « je peux me déplacer à Kribi ») et l'on \textbf{sollicite un entretien}, car le but réel de la lettre n'est pas d'être embauché sur-le-champ, mais d'obtenir un rendez-vous.

\begin{exemplebox}[Formulations du « nous »]
\begin{itemize}
\item « Intégrer votre équipe technique me permettrait de mettre mes compétences au service de la satisfaction de votre clientèle. »
\item « Je reste à votre entière disposition pour un entretien, au cours duquel je pourrai vous exposer plus en détail mes motivations. »
\item « Disponible dès à présent, je serais heureux de vous rencontrer à votre convenance. »
\end{itemize}
\end{exemplebox}

\courssub{Articuler l'argumentation : les connecteurs logiques}

Un corps de lettre n'est pas une juxtaposition de phrases : c'est un raisonnement. Les \cle{connecteurs logiques} en sont les articulations. Pour ordonner les idées, on utilise \textbf{d'abord}, \textbf{ensuite}, \textbf{enfin} ; pour ajouter un argument, \textbf{en outre}, \textbf{de plus}, \textbf{également} ; pour conclure, \textbf{c'est pourquoi}, \textbf{ainsi}, \textbf{par conséquent}. Ils s'appuient sur la phrase complexe étudiée dans la section précédente : subordonnées de cause (\emph{puisque}, \emph{parce que}) et de conséquence (\emph{si bien que}) renforcent la démonstration.

\begin{exoresolu}[Transformer une liste en argumentation]
\textbf{Énoncé.} Voici trois phrases brutes d'un brouillon d'élève : « J'ai un Bac technique F3. J'ai fait un stage de deux mois à la Camwater de Bafoussam. Je connais la maintenance des pompes. » Réécrivez-les en un paragraphe argumenté avec des connecteurs.
\tcblower
\textbf{Solution.} « Titulaire d'un Baccalauréat technique F3 (Électrotechnique), j'ai \textbf{d'abord} acquis une solide formation théorique en installations électriques. J'ai \textbf{ensuite} effectué un stage de deux mois à la Camwater de Bafoussam, qui m'a permis de pratiquer la maintenance des pompes et des tableaux de commande. \textbf{En outre}, ce stage m'a appris la rigueur des consignes de sécurité. \textbf{C'est pourquoi} je pense pouvoir être rapidement opérationnel au sein de votre service technique. » Chaque fait est relié au suivant et l'ensemble conclut vers le poste.
\end{exoresolu}

\begin{exemplebox}[Corps complet : un technicien en froid et climatisation]
Candidature d'un élève de Première/Terminale technique (série F) à un poste de technicien frigoriste dans un hôtel de Kribi.

\emph{Vous :} « Votre hôtel, l'un des établissements phares du littoral sud-camerounais, accueille une clientèle nationale et internationale exigeante, pour qui la climatisation des chambres et la conservation des denrées sont essentielles. »

\emph{Moi :} « Titulaire d'un Baccalauréat technique F4 (Génie civil — option froid et climatisation), j'ai effectué un stage de trois mois dans un atelier de réfrigération de Douala, où j'ai participé à l'installation et à la maintenance de douze climatiseurs et de deux chambres froides. Rigoureux et ponctuel, je tiens à jour un carnet de suivi des interventions, comme l'a noté mon maître de stage dans son attestation. »

\emph{Nous :} « Rejoindre votre équipe technique me permettrait de mettre ces compétences au service du confort de vos clients. Disponible immédiatement et mobile sur toute la région du Sud, je reste à votre disposition pour un entretien à votre convenance. »
\end{exemplebox}

\courssub{Les tonalités : dramatique et comique}

Changeons de regard. En \textbf{réception de textes} (théâtre, roman, récit), nous devons savoir identifier la \cle{tonalité} d'un texte, c'est-à-dire l'émotion dominante que l'auteur cherche à faire naître chez le lecteur ou le spectateur.

\begin{definition}[Tonalité dramatique et tonalité comique]
La \cle{tonalité dramatique} (ou registre \textbf{pathétique}) vise à susciter une \textbf{émotion forte} : pitié, terreur, tristesse, angoisse. Elle repose sur la \textbf{tension} (un conflit, un danger, un destin qui s'accomplit), l'accumulation des malheurs, l'expression des sentiments (exclamations, apostrophes, gradation) et souvent l'annonce d'une catastrophe.

La \cle{tonalité comique} vise à provoquer le \textbf{rire} ou le sourire. On distingue :
\begin{itemize}
\item le \textbf{comique de mots} : le rire naît du langage (calembours, mots inventés, contrepèteries, jargon ridicule) ;
\item le \textbf{comique de situation} : le rire naît des circonstances (quiproquo, malentendu, rebondissement inattendu) ;
\item le \textbf{comique de caractère} : le rire naît d'un défaut exagéré d'un personnage (l'avare, le vaniteux, le jaloux) ;
\item le \textbf{comique de répétition} : le rire naît de la répétition mécanique d'un geste, d'une réplique ou d'une situation.
\end{itemize}
\end{definition}

\begin{propriete}[Les procédés du comique]
Les principaux procédés qui créent la tonalité comique sont :
\begin{itemize}
\item le \cle{quiproquo} : un personnage en prend un autre (ou une chose) pour ce qu'il n'est pas, et parle à côté de la réalité ;
\item l'\cle{ironie} : dire le contraire de ce que l'on pense, pour se moquer sans avoir l'air d'y toucher ;
\item l'\cle{exagération} (hyperbole comique) : grossir un défaut ou une situation jusqu'à l'absurde ;
\item le \cle{décalage} : contraste entre le ton et le sujet (parler d'une bagatelle avec un ton tragique, ou d'un drame avec légèreté), entre ce qu'un personnage croit et ce que le public sait.
\end{itemize}
\end{propriete}

\begin{popfigure}
\begin{center}
\begin{tikzpicture}[font=\small]
\draw[fill=popblue, text=white, font=\bfseries] (0,3.6) rectangle (11.5,4.4);
\node[text=white, font=\bfseries] at (1.9,4) {Type de comique};
\node[text=white, font=\bfseries] at (5.6,4) {Principe};
\node[text=white, font=\bfseries] at (9.3,4) {Exemple};
\draw[fill=popblueL, draw=popblue] (0,2.4) rectangle (3.8,3.6);
\node[align=center, text width=3.4cm] at (1.9,3) {\textbf{Comique de mots}\\ (calembour, jargon)};
\draw[fill=popcream, draw=popline] (3.8,2.4) rectangle (7.4,3.6);
\node[align=center, text width=3.2cm] at (5.6,3) {Le langage lui-même fait rire};
\draw[fill=popcream, draw=popline] (7.4,2.4) rectangle (11.5,3.6);
\node[align=center, text width=3.7cm] at (9.45,3) {« Il a des \emph{poules} de luxe » pour « des \emph{poules}... » (calembour sur un mot)};
\draw[fill=popblueL, draw=popblue] (0,1.2) rectangle (3.8,2.4);
\node[align=center, text width=3.4cm] at (1.9,1.8) {\textbf{Comique de situation}\\ (quiproquo)};
\draw[fill=popcream, draw=popline] (3.8,1.2) rectangle (7.4,2.4);
\node[align=center, text width=3.2cm] at (5.6,1.8) {Malentendu, confusion d'identité};
\draw[fill=popcream, draw=popline] (7.4,1.2) rectangle (11.5,2.4);
\node[align=center, text width=3.7cm] at (9.45,1.8) {Un valet prend son maître déguisé pour un étranger};
\draw[fill=popblueL, draw=popblue] (0,0) rectangle (3.8,1.2);
\node[align=center, text width=3.4cm] at (1.9,0.6) {\textbf{Comique de caractère}\\ (défaut exagéré)};
\draw[fill=popcream, draw=popline] (3.8,0) rectangle (7.4,1.2);
\node[align=center, text width=3.2cm] at (5.6,0.6) {Un trait poussé à l'extrême};
\draw[fill=popcream, draw=popline] (7.4,0) rectangle (11.5,1.2);
\node[align=center, text width=3.7cm] at (9.45,0.6) {L'avare qui compte et re-compte ses sous (Harpagon)};
\end{tikzpicture}
\end{center}
\legende{Les trois grandes sources du comique, avec un exemple pour chacune.}
\end{popfigure}

\begin{exemplebox}[Observer les tonalités dans les textes de réception]
Dans un récit dramatique, repérez : les exclamations (« Hélas ! », « Ô douleur ! »), le champ lexical du malheur (larmes, sang, destin, nuit), la gradation qui précipite le héros vers la catastrophe. Dans une scène comique, repérez : le malentendu initial, les répliques qui se répètent, le personnage qui se contredit sans s'en rendre compte, la chute inattendue qui dégonfle la prétention.
\end{exemplebox}

\begin{attention}[Tonalité et lettre de motivation : ne pas confondre les genres]
La tonalité dramatique et la tonalité comique s'\textbf{observent} dans les textes de réception (théâtre, roman), mais elles sont \textbf{interdites} dans la lettre de motivation, qui exige une tonalité \textbf{sérieuse et respectueuse}. Supplier (« Je vous en conjure, ayez pitié d'un pauvre orphelin ! ») relève du pathétique déplacé ; plaisanter (« Embauchez-moi, vous ne le regretterez pas... ou presque ! ») relève du comique déplacé. L'employeur attend de la mesure : ni larmes, ni blagues. Réservez l'analyse des tonalités aux textes littéraires, et la sobriété à vos écrits utilitaires.
\end{attention}

\begin{exercice}[À vous de jouer]
1. Un élève écrit dans sa lettre : « Sans ce stage, ma vie est finie, je vous supplie de me sauver ! » Quelle tonalité emploie-t-il ? Pourquoi est-ce inadapté ? Réécrivez la phrase correctement.

2. Classez les procédés suivants (quiproquo, ironie, exagération, décalage) : a) « Quel génie ! » dit à quelqu'un qui vient de casser un vase ; b) un personnage parle à un portrait en le prenant pour son rival ; c) « Il a ri si fort que les murs de Yaoundé ont tremblé » ; d) un personnage annonce la fin du monde parce qu'il a perdu son parapluie.

3. Rédigez le paragraphe « moi » d'une lettre de motivation pour un stage d'été dans un garage de Bafoussam (vous êtes en Première F3), en trois ou quatre phrases avec au moins deux connecteurs.
\end{exercice}

\begin{corrige}[Exercice]
1. C'est la tonalité \textbf{dramatique} (pathétique) : supplication, exclamation, dramatisation. Inadaptée car une lettre officielle exige la mesure. Correction possible : « Ce stage représente pour moi une étape importante de ma formation ; c'est pourquoi je sollicite votre bienveillance. »

2. a) ironie ; b) quiproquo ; c) exagération ; d) décalage (ton tragique pour un sujet futile).

3. Exemple : « Actuellement en classe de Première F3 (Électrotechnique) au lycée technique, j'ai d'abord acquis des bases solides en électricité automobile. J'ai ensuite participé, pendant les vacances, à la réparation du système de démarrage d'un véhicule familial, sous la conduite d'un mécanicien confirmé. En outre, je suis ponctuel et soigneux, comme en témoigne mon carnet d'atelier. C'est pourquoi je serais heureux d'effectuer mon stage d'été dans votre garage. »
\end{corrige}

\begin{aretenir}[L'essentiel de la section]
\begin{itemize}
\item Le corps de la lettre suit le plan « \cle{vous – moi – nous} » : l'entreprise, puis le candidat, puis la collaboration et la demande d'entretien.
\item Le « moi » \textbf{argumente} à partir de faits vérifiables (diplômes, stages) ; il ne recopie pas le CV et ne ment jamais.
\item Les connecteurs (\emph{d'abord, ensuite, en outre, c'est pourquoi}) transforment une liste en raisonnement.
\item La tonalité \textbf{dramatique} (pathétique, tension, émotion forte) et la tonalité \textbf{comique} (de mots, de situation, de caractère, de répétition ; quiproquo, ironie, exagération, décalage) s'analysent dans les textes littéraires.
\item La lettre de motivation exige une tonalité \textbf{sérieuse} : ni drame, ni plaisanterie.
\end{itemize}
\end{aretenir}

\coursec{Formules de conclusion et textes iconiques}

Dans la section précédente, nous avons appris à construire le \cle{corps} de la lettre de motivation : l'accroche, la présentation du candidat, l'argumentation et la tonalité adaptée au contexte. Mais une lettre, comme une phrase, ne doit pas s'arrêter brusquement. Imaginez un candidat qui, après un excellent entretien, quitte le bureau du directeur sans le saluer : toute sa prestation serait gâchée. De même, une lettre de motivation mal terminée laisse une mauvaise impression, même si le corps est bien rédigé. Cette section vous apprendra à \emph{conclure} correctement votre lettre, puis à lire un autre type de texte que vous rencontrerez à l'examen et dans la vie professionnelle : le \cle{texte iconique}, c'est-à-dire l'image porteuse de sens.

\courssub{La formule de demande d'entretien}

La lettre de motivation a un but précis : obtenir un \cle{entretien}. Il faut donc le dire clairement, sans arrogance et sans supplication. C'est le rôle de la formule de demande d'entretien, placée juste avant la formule de politesse finale.

\begin{definition}[Formule de demande d'entretien]
La \cle{formule de demande d'entretien} est une phrase brève, placée à la fin du corps de la lettre, par laquelle le candidat propose au destinataire de le rencontrer. Elle témoigne de la disponibilité et de la motivation du candidat.
\end{definition}

Modèles corrects :
\begin{itemize}
\item \og Je reste à votre disposition pour un entretien à votre convenance. \fg
\item \og Je me tiens à votre entière disposition pour tout entretien que vous voudrez bien m'accorder. \fg
\item \og Disponible immédiatement, je serais heureux de vous exposer ma motivation lors d'un entretien. \fg
\end{itemize}

\begin{attention}[Piège fréquent]
Ne terminez jamais le corps de la lettre par \og Merci d'avance \fg ou \og Merci de me répondre vite \fg : ces formules sont familières, impatientes et impolies. De même, évitez de mendier (\og Je vous supplie de m'accepter \fg) : la demande d'entretien doit rester digne et confiante.
\end{attention}

\courssub{Les formules de politesse finales}

La formule de politesse finale est la \emph{signature verbale} de votre lettre. Elle obéit à un code social strict : plus le destinataire est haut placé dans la hiérarchie, plus la formule est solennelle.

\begin{definition}[Formule de politesse finale]
La \cle{formule de politesse finale} est une phrase conventionnelle qui clôt la lettre officielle. Elle exprime le respect du candidat envers le destinataire. Sa structure type est : \emph{Je vous prie d'agréer} + \emph{l'appellation du destinataire} + \emph{l'expression de mes salutations distinguées / de ma considération distinguée / de ma haute considération}.
\end{definition}

\begin{exemplebox}[Trois formules hiérarchisées selon le destinataire]
\begin{itemize}
\item À un \textbf{directeur des ressources humaines} (DRH) : \og Je vous prie d'agréer, Monsieur le Directeur, l'expression de mes salutations distinguées. \fg
\item À un \textbf{maire} : \og Je vous prie d'agréer, Monsieur le Maire, l'expression de ma considération distinguée. \fg
\item À un \textbf{ministre} : \og Je vous prie d'agréer, Monsieur le Ministre, l'expression de ma haute considération. \fg
\end{itemize}
\end{exemplebox}

\begin{attention}[Erreur fréquente]
N'écrivez jamais \og Cordialement \fg ou \og Amitiés \fg à un ministre, un préfet ou un directeur : ce sont des formules réservées aux courriels entre collègues ou amis. À l'examen, une formule familière dans une lettre officielle fait perdre des points de correction et de présentation. Autre faute classique : oublier l'appellation du destinataire au milieu de la formule (\og Monsieur le Directeur, \fg).
\end{attention}

\courssub{La signature et les pièces jointes}

Après la formule de politesse, la lettre se termine par la \cle{signature} :
\begin{enumerate}
\item le \textbf{nom et le prénom} du candidat, écrits lisiblement ;
\item la \textbf{signature manuscrite}, apposée au-dessus ou à côté du nom ;
\item la mention des \textbf{pièces jointes}, notée \og P. J. \fg ou \og Pièces jointes \fg, en bas à gauche : elle annonce les documents accompagnant la lettre (curriculum vitae, photocopies des diplômes, attestations de stage).
\end{enumerate}

\begin{exemplebox}[Modèle de fin de lettre]
Je vous prie d'agréer, Monsieur le Directeur, l'expression de ma considération distinguée.\par
\vspace{0.5em}
\hfill \emph{Signature}\par
\hfill NGA Marie-Claire\par
\vspace{0.5em}
P. J. : un curriculum vitae ; une photocopie du Probatoire technique ; une attestation de stage.
\end{exemplebox}

\courssub{Le texte iconique : définition}

Toutes les lettres ne se lisent pas : certaines informations passent par l'\emph{image}. Une affiche de recrutement, une caricature dans le journal, une photographie en première page : ce sont des textes à part entière, que l'on appelle \cle{textes iconiques}.

\begin{definition}[Texte iconique]
Un \cle{texte iconique} est une image porteuse de sens, produite pour communiquer un message. Il peut s'agir d'une \textbf{photographie}, d'une \textbf{affiche}, d'une \textbf{caricature}, d'un \textbf{dessin de presse}, d'un \textbf{logo} ou d'une \textbf{publicité}. Comme un texte écrit, il possède un auteur, un destinataire, un contexte et une intention (informer, convaincre, dénoncer, faire rire).
\end{definition}

\courssub{Lecture méthodique d'un texte iconique}

Lire une image ne consiste pas à dire \og je vois une photo \fg. Il faut procéder en \textbf{trois étapes}, exactement comme pour un texte écrit.

\begin{methode}[Analyser une image en trois étapes]
\begin{enumerate}
\item \textbf{Décrire} : présenter objectivement l'image (nature du document, source, date, éléments visibles : personnages, objets, décor, couleurs, légende éventuelle). On répond aux questions : quoi ? qui ? où ?
\item \textbf{Analyser les codes} : étudier les procédés visuels : les \emph{couleurs} (le rouge signale le danger ou l'urgence, le vert l'espoir ou la nature), le \emph{cadrage} (gros plan = émotion ; plan large = situation d'ensemble), la \emph{lumière}, les \emph{symboles}, la place des personnages, les mots écrits dans l'image.
\item \textbf{Interpréter} : dégager le message et l'intention de l'auteur (informer, sensibiliser, critiquer, recruter), puis identifier la \emph{tonalité} (comique, dramatique, polémique).
\end{enumerate}
\end{methode}

\begin{popfigure}
\begin{center}
\begin{tikzpicture}[node distance=0.6cm, every node/.style={font=\small}]
\node[draw=popblue, fill=popblueL, rounded corners, minimum width=3.4cm, minimum height=1.1cm, align=center] (e1) {\textbf{1. Décrire}\\ ce que l'on voit};
\node[draw=poporange, fill=poporangeL, rounded corners, minimum width=3.4cm, minimum height=1.1cm, align=center, right=of e1] (e2) {\textbf{2. Analyser}\\ couleurs, cadrage,\\ symboles};
\node[draw=popgreen, fill=popgreenL, rounded corners, minimum width=3.4cm, minimum height=1.1cm, align=center, right=of e2] (e3) {\textbf{3. Interpréter}\\ message, intention,\\ tonalité};
\draw[-{Stealth}, very thick, popdark] (e1) -- (e2);
\draw[-{Stealth}, very thick, popdark] (e2) -- (e3);
\end{tikzpicture}
\end{center}
\legende{La démarche d'analyse d'un texte iconique en trois étapes : de l'observation objective à l'interprétation du message.}
\end{popfigure}

\courssub{Étude des trois textes iconiques de l'unité}

\textbf{Texte iconique n° 1 : une affiche de recrutement d'une entreprise camerounaise.} Description : l'affiche porte le logo de l'entreprise en haut, un titre en grosses lettres (\og Nous recrutons ! \fg), la liste des postes offerts, les conditions requises et les contacts. Analyse : les couleurs vives (bleu, orange) attirent le regard ; le titre en caractères gras hiérarchise l'information ; le logo rassure sur le sérieux de l'entreprise. Interprétation : l'affiche veut \emph{informer} et \emph{convaincre} des candidats de postuler ; elle montre que la lettre de motivation que nous rédigeons répond à un besoin réel du monde professionnel.

\textbf{Texte iconique n° 2 : une caricature de presse camerounaise.} Description : le dessin représente un personnage aux traits exagérés (énorme ventre, petit chapeau), assis sur un fauteuil pendant qu'une file de jeunes diplômés attend devant sa porte. Analyse : l'\emph{exagération} des traits (charge), la disproportion entre le fauteuil confortable et la file d'attente, la bulle de parole ironique. Interprétation : le caricaturiste dénonce, par le rire, le favoritisme dans les recrutements. La tonalité est \cle{comique} : on rit, mais pour mieux critiquer.

\begin{savaistu}[Le pouvoir de la caricature]
On dit qu'une caricature peut dire en une image ce qu'un éditorial dit en mille mots. Dans la presse camerounaise comme dans la presse française, les dessinateurs satiriques jouent un rôle de vigie : en une seule charge, ils résument une injustice que tout le monde reconnaît immédiatement. C'est pourquoi la caricature est un genre redouté par les puissants et protégé par la liberté de la presse.
\end{savaistu}

\textbf{Texte iconique n° 3 : une photographie dramatique.} Description : la photo montre une foule de personnes fuyant une zone sinistrée (inondation, incendie), des bagages sur la tête, des enfants dans les bras. Analyse : les tons sombres, le cadrage serré sur les visages marqués par la peur, la boue et la fumée en arrière-plan. Interprétation : le photographe veut \emph{témoigner} et \emph{émouvoir} l'opinion ; la tonalité est \cle{dramatique}. Une telle image peut accompagner un appel à l'aide humanitaire.

\begin{exoresolu}[Analyse guidée d'une affiche de recrutement]
Énoncé : une entreprise de Douala publie une affiche : fond bleu, titre jaune \og Jeunes diplômés, rejoignez-nous ! \fg, photo de jeunes employés souriants en uniforme, liste de trois postes, adresse électronique en bas. Décrivez, analysez et interprétez cette affiche.
\tcblower
\textbf{Solution.} \emph{Description} : il s'agit d'une affiche publicitaire de recrutement ; on y voit un titre, une photographie de jeunes employés et des informations pratiques. \emph{Analyse} : le fond bleu inspire confiance et sérieux ; le titre jaune sur fond bleu crée un contraste qui attire l'œil ; le sourire des employés et le tutoiement implicite (\og rejoignez-nous \fg) créent une proximité avec le lecteur jeune ; les informations pratiques sont placées en bas, selon le parcours naturel du regard. \emph{Interprétation} : l'entreprise veut attirer des candidats jeunes en présentant une image dynamique et accueillante ; la tonalité est publicitaire et chaleureuse. Le candidat qui répondra à cette offre devra rédiger une lettre de motivation soignée, conforme aux codes étudiés dans cette unité.
\end{exoresolu}

\courssub{Complément utile au Bac : l'image comme argument}

Dans une campagne de communication (publicité, sensibilisation à la santé ou à la citoyenneté), l'image fonctionne comme un \cle{argument} : elle prouve, elle touche, elle convainc plus vite qu'un long discours. Une photographie d'enfants vaccinés plaide pour la vaccination ; une affiche montrant un élève assidu plaide pour l'école. À l'examen, face à un texte iconique, appliquez toujours la méthode en trois étapes et n'oubliez pas de nommer la tonalité : c'est elle qui révèle l'intention de l'auteur.

\begin{aretenir}[L'essentiel de la section]
\begin{itemize}
\item La lettre de motivation se termine par une \cle{demande d'entretien} (\og Je reste à votre disposition pour un entretien\ldots \fg), une \cle{formule de politesse} adaptée au destinataire (salutations distinguées, considération distinguée, haute considération), la \cle{signature} et les \cle{pièces jointes}.
\item Jamais de \og Merci d'avance \fg ni de \og Cordialement \fg dans une lettre officielle.
\item Un \cle{texte iconique} est une image porteuse de sens (photographie, affiche, caricature, logo).
\item On l'analyse en trois étapes : \textbf{décrire}, \textbf{analyser les codes} (couleurs, cadrage, symboles), \textbf{interpréter} (message, intention, tonalité comique ou dramatique).
\end{itemize}
\end{aretenir}

\begin{exercice}[À vous de jouer]
1. Rédigez la fin complète d'une lettre de motivation adressée au maire de votre commune (demande d'entretien, formule de politesse, signature, deux pièces jointes).\par
2. Votre journal publie la photographie d'un marché incendié : décrivez-la, analysez deux codes visuels probables et dégagez la tonalité.\par
3. Pourquoi la formule \og Cordialement \fg est-elle inadaptée dans une lettre adressée à un ministre ? Proposez la formule correcte.\par
4. En binôme, échangez vos productions de l'exercice 1, repérez les éventuelles maladresses (formule de politesse, pièces jointes, présentation) et proposez une version améliorée.
\end{exercice}

\coursec{Rédaction, confrontation et amélioration de la lettre}

Dans la section précédente, nous avons appris à conclure une lettre de motivation par des formules de politesse adaptées et à lire des textes iconiques. Nous disposons désormais de toutes les pièces du puzzle : l'en-tête, le corps en trois paragraphes, les formules de conclusion, la phrase complexe et le lexique soutenu. Il reste à assembler ces pièces, c'est-à-dire à \cle{rédiger} la lettre complète, puis à la \cle{confronter} au regard des autres et à l'\cle{améliorer}. C'est toute la démarche de production d'écrits que nous allons suivre pas à pas.

\courssub{Le cercle vertueux de la production d'écrit}

\begin{definition}[Produire un écrit]
\cle{Produire un écrit}, ce n'est pas seulement écrire d'un seul jet. C'est suivre une démarche en cinq étapes : \textbf{préparer} (analyser la situation et rassembler ses idées), \textbf{rédiger} un brouillon, \textbf{relire} pour corriger, \textbf{confronter} sa production à celle des pairs, puis \textbf{améliorer} pour aboutir à la version définitive.
\end{definition}

L'intuition est simple : aucun écrivain, aucun secrétaire, aucun candidat sérieux n'envoie jamais son premier jet. Une lettre de motivation efficace est le résultat d'un \emph{processus}, pas d'un instant d'inspiration. Chaque étape nourrit la suivante, et la dernière peut ramener à la première si l'on découvre un défaut majeur : c'est un cercle, non une ligne droite.

\begin{popfigure}
\begin{center}
\begin{tikzpicture}[scale=1, every node/.style={align=center}]
\foreach \angle/\txt/\col in {90/{1. Préparer\\analyser l'offre}/popblue, 18/{2. Rédiger\\le brouillon}/popgreen, -54/{3. Relire\\corriger}/popgold, -126/{4. Confronter\\échanger}/poporange, 162/{5. Améliorer\\version finale}/poppurple}{
  \node[draw, circle, fill=\col!35, minimum size=2.2cm, font=\small\bfseries] at (\angle:3.2) {\txt};
}
\foreach \a/\b in {90/18, 18/-54, -54/-126, -126/162, 162/90}{
  \draw[-{Stealth[length=3mm]}, very thick, popdark] (\a:2.15) arc[start angle=\a, end angle=\b, radius=2.15];
}
\node[font=\small\itshape, popmuted, align=center] at (0,0){Le cercle vertueux\\de l'écriture};
\end{tikzpicture}
\end{center}
\legende{Schéma : le cercle vertueux de la production d'écrit. Les cinq étapes s'enchaînent ; l'amélioration peut renvoyer à la préparation si un problème de fond est détecté.}
\end{popfigure}

\begin{methode}[Rédiger une lettre de motivation en 5 étapes]
\begin{enumerate}
\item \textbf{Préparer.} Lis attentivement l'offre d'emploi ou de stage : souligne le poste, l'entreprise, les qualités exigées. Liste ensuite tes \cle{atouts} (diplômes, stages, qualités personnelles) et choisis deux ou trois \cle{arguments} qui répondent exactement à la demande.
\item \textbf{Rédiger le brouillon.} Respecte la structure : en-tête (coordonnées, destinataire, lieu et date, objet), corps en trois paragraphes (vous — moi — nous), formules de conclusion, signature. Ne te soucie pas encore de la perfection : fais sortir les idées.
\item \textbf{Relire.} Vérifie l'orthographe, les accords, la construction des phrases complexes, le niveau de langue. Remplace tout mot familier par un mot soutenu.
\item \textbf{Confronter.} Échange ta lettre avec un camarade et évalue la sienne à l'aide d'une grille critériée. Note les remarques reçues sans te vexer : elles sont un cadeau.
\item \textbf{Améliorer.} Intègre les remarques pertinentes et réécris la version définitive, propre et complète.
\end{enumerate}
\end{methode}

\courssub{Étape 1 — Préparation : analyser l'offre et choisir ses arguments}

Avant d'écrire la moindre phrase, il faut \emph{comprendre à qui l'on écrit et pourquoi}. Prenons une offre concrète.

\begin{exemplebox}[Offre de stage à analyser]
« La SODECOTON de Garoua recherche des stagiaires titulaires d'un Probatoire technique, sérieux, ponctuels et ayant le sens du travail en équipe, pour la période de juillet à septembre. Adresser lettre de motivation et CV au Chef du personnel. »

Analyse : le \textbf{destinataire} est le Chef du personnel ; le \textbf{poste} est un stage ; les \textbf{qualités exigées} sont le sérieux, la ponctualité, l'esprit d'équipe. Mes arguments devront donc prouver ces trois qualités par des faits (assiduité en atelier, projets de groupe réussis), et non les affirmer gratuitement.
\end{exemplebox}

\begin{attention}[Ne pas confondre atout et argument]
Un \cle{atout} est une qualité que tu possèdes ; un \cle{argument} est un atout \emph{présenté comme une réponse au besoin de l'employeur}. Dire « je suis ponctuel » est un atout ; écrire « assidu à tous les cours et travaux pratiques depuis trois ans, je mettrai cette ponctualité au service de votre équipe » est un argument. C'est l'argument qui convainc.
\end{attention}

\courssub{Étape 2 — Le brouillon : assembler les pièces}

Le brouillon respecte la structure étudiée dans les sections précédentes. Rappelons-la sous forme de tableau.

\begin{center}
\begin{tabularx}{\linewidth}{|l|X|}
\hline
\rowcolor{popblueL} \textbf{Partie} & \textbf{Contenu attendu} \\
\hline
En-tête & Coordonnées de l'expéditeur (en haut à gauche), destinataire (à droite), lieu et date, objet de la lettre. \\
\hline
Paragraphe 1 : « vous » & Montrer qu'on connaît l'entreprise et l'intérêt qu'on lui porte. \\
\hline
Paragraphe 2 : « moi » & Présenter sa formation et ses atouts en lien avec le poste. \\
\hline
Paragraphe 3 : « nous » & Proposer la collaboration, exprimer sa disponibilité pour un entretien. \\
\hline
Conclusion & Formule de politesse soutenue, signature. \\
\hline
\end{tabularx}
\end{center}

\begin{exemplebox}[Extrait de brouillon (premier jet)]
« Je vous écris parce que je veux faire un stage chez vous. J'ai mon probatoire technique et je suis quelqu'un de bien. Je peux venir quand vous voulez. »

Ce premier jet contient les idées, mais le ton est familier (« quelqu'un de bien »), les phrases sont simples et pauvres, et le paragraphe « vous » est absent. C'est normal : le brouillon sert précisément à repérer ces faiblesses.
\end{exemplebox}

\courssub{Étape 3 — Relecture orthographique et syntaxique}

La relecture porte sur trois niveaux, toujours dans cet ordre : le \textbf{fond} (les arguments sont-ils là ?), la \textbf{syntaxe} (les phrases complexes sont-elles bien construites ?), l'\textbf{orthographe} (accords, conjugaisons).

\begin{attention}[Erreur fréquente : négliger la relecture]
Beaucoup d'élèves recopient leur brouillon sans le relire. Or une seule faute d'orthographe dans une lettre de motivation peut \textbf{disqualifier une candidature} : l'employeur y voit un manque de sérieux. Relis ta lettre au moins deux fois, dont une fois à voix haute, et fais-la relire par un camarade : l'œil de l'autre voit ce que le tien ne voit plus.
\end{attention}

Points de vigilance pour la relecture :
\begin{itemize}
\item \textbf{Accords} : sujet-verbe (« les qualités que je possède »), participe passé (« la lettre que j'ai écrite »).
\item \textbf{Phrase complexe} : chaque subordonnée doit avoir son verbe conjugué et être introduite par un mot de liaison correct (« bien que je sois jeune », et non « bien que je suis jeune »).
\item \textbf{Lexique soutenu} : remplacer « je veux » par « je souhaite », « bosser » par « travailler », « je suis quelqu'un de bien » par « je suis rigoureux et ponctuel ».
\end{itemize}

\courssub{Étape 4 — Confrontation : l'évaluation croisée entre pairs}

La \cle{confrontation} consiste à échanger les productions au sein de la classe : chaque élève évalue la lettre d'un camarade à l'aide d'une \cle{grille critériée}, et reçoit en retour l'évaluation de la sienne. Ce regard extérieur est irremplaçable : il révèle les maladresses que l'auteur ne perçoit plus.

\begin{methode}[Évaluer la production d'un pair avec une grille critériée]
\begin{enumerate}
\item Lis la lettre entièrement une première fois, sans rien noter, pour saisir l'impression générale.
\item Relis en t'arrêtant sur chacun des cinq critères de la grille : présentation, respect des normes de la lettre, qualité argumentative, correction de la langue, niveau de langue.
\item Pour chaque critère, attribue un niveau de 1 (insuffisant) à 4 (excellent) et justifie ton choix par une remarque précise et bienveillante.
\item Formule au moins un \textbf{conseil d'amélioration} concret (« remplace ce mot familier », « ajoute un argument au paragraphe 2 »).
\item Restitue la grille à ton camarade et discutez ensemble des remarques.
\end{enumerate}
\end{methode}

\begin{popfigure}
\begin{center}
\begin{tikzpicture}[x=2.9cm, y=0.85cm]
\fill[popblue!30] (0,5) rectangle (5,6);
\node[font=\small\bfseries] at (0.85,5.5) {Critère};
\node[font=\small\bfseries] at (2.0,5.5) {1. Insuffisant};
\node[font=\small\bfseries] at (3.05,5.5) {2. Passable};
\node[font=\small\bfseries] at (4.05,5.5) {3. Bien};
\node[font=\small\bfseries] at (4.9,5.5) {4. Excellent};
\foreach \i/\txt in {0/{Présentation matérielle},1/{Respect des normes (en-tête, objet, formules)},2/{Qualité argumentative},3/{Correction de la langue},4/{Niveau de langue soutenu}}{
  \fill[popcream] (0,4-\i) rectangle (5,5-\i);
  \node[anchor=west, font=\small] at (0.05,4.5-\i) {\txt};
}
\foreach \x in {0,1.7,2.7,3.7,4.7,5}{\draw (\x,0) -- (\x,6);}
\foreach \y in {0,1,2,3,4,5,6}{\draw (0,\y) -- (5,\y);}
\end{tikzpicture}
\end{center}
\legende{Grille d'évaluation croisée : cinq critères cotés sur quatre niveaux. L'évaluateur coche un niveau par critère et ajoute une remarque justificative.}
\end{popfigure}

\begin{savaistu}[L'évaluation entre pairs et l'APC]
La confrontation des écrits entre pairs est une pratique recommandée par l'\cle{approche par les compétences} (APC) du MINESEC : en évaluant la production d'un camarade, l'élève développe son esprit critique, intériorise les critères de qualité et apprend à mieux évaluer ses propres écrits. Évaluer, c'est déjà progresser.
\end{savaistu}

\courssub{Étape 5 — Amélioration : de la version corrigée à la version définitive}

L'amélioration consiste à \emph{intégrer} les remarques reçues et ses propres corrections, puis à \emph{réécrire} entièrement la lettre. Il ne s'agit pas de raturer le brouillon, mais de produire une version nouvelle, propre et aboutie.

\begin{exoresolu}[Avant / après : les 8 erreurs corrigées d'une lettre d'élève]
Énoncé. Voici un extrait de la première version d'une lettre d'élève, contenant 8 erreurs. Repère-les et propose la version améliorée.

\emph{Version 1 :} « Monsieur le Directeur, je vous écris pour que vous me donné un stage dans votre entreprise. J'ai eu mon probatoire technique l'année dernière et j'ai fait de bon résultats. Je suis quelqu'un qui travail bien et qui est toujours à l'heure. Je peux commencer quand vous voulez. Salutations. »

\tcblower
Solution détaillée. Les 8 erreurs : (1) « donné » au lieu de « donniez » (subjonctif après « pour que ») ; (2) « de bon résultats » au lieu de « de bons résultats » ; (3) « travail » au lieu de « travaille » ; (4) « quelqu'un qui » : tournure lourde à remplacer ; (5) « quand vous voulez » : registre familier ; (6) « Salutations » : formule de conclusion trop familière ; (7) absence du paragraphe « vous » sur l'entreprise ; (8) absence de proposition d'entretien.

\emph{Version 2 améliorée :} « Monsieur le Directeur, votre entreprise, reconnue pour la qualité de sa production, forme chaque année de jeunes techniciens : c'est pourquoi je sollicite un stage au sein de vos ateliers. Titulaire du Probatoire technique, j'ai obtenu de bons résultats en travaux pratiques. Rigoureux et ponctuel, je mettrai ces qualités au service de votre équipe. Je reste à votre disposition pour tout entretien que vous jugerez utile. Dans cette attente, je vous prie d'agréer, Monsieur le Directeur, l'expression de ma haute considération. »
\end{exoresolu}

\courssub{Compte rendu et remédiation}

En fin de séquence, l'enseignant établit un \cle{compte rendu} : bilan des acquis (structure de la lettre maîtrisée, formules de politesse connues) et recensement des difficultés persistantes (accords du participe passé, subjonctif, objet de la lettre trop vague). Sur cette base, la \cle{remédiation} propose des exercices ciblés.

\begin{exercice}[Remédiation 1 : les formules de politesse]
Complète avec la formule soutenue convenable : « Je vous prie d'\_\_\_\_\_\_, Monsieur le Chef du personnel, l'expression de ma \_\_\_\_\_\_\_\_\_\_\_\_\_\_. » Puis transforme la formule familière « Amitiés » en formule adaptée à une lettre officielle.
\end{exercice}

\begin{exercice}[Remédiation 2 : la subordination]
Relie les propositions par le mot de subordination indiqué : a) « Je suis jeune. / J'ai déjà une expérience de stage. » (bien que) ; b) « Je vous écris. / Vous me proposiez un entretien. » (afin que).
\end{exercice}

\begin{exercice}[Remédiation 3 : l'objet de la lettre]
Rédige un objet précis et court pour une lettre adressée au Directeur de la CAMTEL en vue d'un stage de deux mois.
\end{exercice}

\begin{corrige}[Exercices de remédiation]
\textbf{1.} « Je vous prie d'\emph{agréer}, Monsieur le Chef du personnel, l'expression de ma \emph{haute considération} » (ou « de mes salutations distinguées »). « Amitiés » devient « Veuillez agréer, Madame, Monsieur, mes salutations distinguées ».

\textbf{2.} a) « Bien que je sois jeune, j'ai déjà une expérience de stage. » (subjonctif après \emph{bien que}). b) « Je vous écris afin que vous me proposiez un entretien. » (subjonctif après \emph{afin que}).

\textbf{3.} « Objet : demande de stage académique de deux mois. » L'objet est bref, précis (nature de la demande, durée) et sans verbe conjugué.
\end{corrige}

\begin{aretenir}[L'essentiel de la section]
Rédiger une lettre de motivation suit cinq étapes : préparer, rédiger, relire, confronter, améliorer. La relecture est obligatoire : une faute peut disqualifier une candidature. La confrontation entre pairs, guidée par une grille à cinq critères (présentation, normes, argumentation, langue, niveau de langue), permet de repérer ses propres faiblesses. La version définitive intègre toutes les remarques. Compte rendu et remédiation consolident enfin les acquis de toute la séquence.
\end{aretenir}

\newpage

\coursec{Activité d'intégration}

\coursec{Leçon 09 : Production d'écrits utilitaires — La lettre de motivation}

\courssub{Objectifs de l'activité}

À l'issue de cette activité d'intégration, l'élève sera capable de :

\begin{itemize}
\item analyser un \cle{texte fonctionnel} (une offre d'emploi) et un \cle{texte iconique} (une affiche de recrutement) pour en dégager les informations utiles ;
\item rédiger une \cle{lettre de motivation} complète et conforme aux normes : en-tête, corps organisé selon le schéma \emph{vous -- moi -- nous}, formules de conclusion adaptées ;
\item employer correctement la \cle{phrase complexe} (subordonnées relatives et conjonctives) et repérer les \cle{emprunts} dans un texte professionnel ;
\item choisir une \cle{tonalité} adaptée à la situation de communication (éviter le comique ou le dramatique inopportun) ;
\item évaluer la production d'un camarade à l'aide d'une \cle{grille critériée}, puis améliorer sa propre lettre et justifier ses corrections par écrit.
\end{itemize}

\courssub{Situation}

Le Lycée technique de Bafoussam organise, à l'approche des examens, sa traditionnelle \emph{Semaine de l'insertion professionnelle}. À cette occasion, le proviseur a affiché au tableau d'information trois offres d'emploi reconstituées à partir d'annonces réellement parues dans la presse camerounaise, accompagnées d'une affiche de recrutement.

\textbf{Offre n° 1 -- Technicien agricole (Bafoussam, région de l'Ouest).}
La SODECOPAL, coopérative de production maraîchère de l'Ouest, recrute un technicien agricole titulaire d'un Baccalauréat technique ou d'un BTS en productions végétales. Missions : encadrement de 12 exploitants, suivi des cultures de pommes de terre et de choux sur 25 hectares, tenue du carnet de rendements. Type de contrat : CDI (\emph{job} permanent). Salaire : \num{120000} FCFA mensuels. Candidature à adresser avant le 30 juin au Directeur de la SODECOPAL, BP 214, Bafoussam.

\textbf{Offre n° 2 -- Comptable (Douala, région du Littoral).}
L'entreprise TRANSLOG SARL, société de transit située à Bonanjo, recrute un comptable débutant. Profil : Probatoire ou Baccalauréat technique série G (comptabilité), rigueur, maîtrise de l'outil informatique. Missions : saisie des factures (environ 150 par mois), rapprochements bancaires, classement des pièces. Salaire : \num{150000} FCFA mensuels. Dossier à déposer au 45, rue Joffre, Akwa, Douala, ou à envoyer par \emph{e-mail} à recrutement@translog.cm.

\textbf{Offre n° 3 -- Électrotechnicien (Garoua, région du Nord).}
La société NORDENERGIE, spécialisée dans l'installation solaire, recrute un électrotechnicien pour la pose et la maintenance de \emph{kits} solaires dans les villages du Mayo-Rey. Profil : Baccalauréat technique F3 ou équivalent, permis B souhaité, disponibilité pour les déplacements. Salaire : \num{130000} FCFA mensuels plus primes de mission de \num{15000} FCFA par déplacement. Écrire au Chef du personnel, NORDENERGIE, BP 87, Garoua.

\textbf{L'affiche iconique} qui accompagne les offres présente, sur fond vert-jaune-rouge, un jeune diplômé en tenue de travail serrant la main d'un employeur, avec le slogan : \emph{« Un diplôme technique, un métier, un avenir ! »} et, en bas, les coordonnées du stand d'orientation du lycée.

Tu es élève en classe de Première technique. Le professeur de français te demande de te projeter : \emph{« Imagine que tu viens d'obtenir ton Baccalauréat technique. Choisis l'une des trois offres et rédige la lettre de motivation que tu enverrais réellement. »}

\courssub{Consignes}

\textbf{Première phase : analyse des documents (travail de groupe).}

\begin{enumerate}
\item Relevez dans l'offre choisie : l'expéditeur (qui recrute ?), le destinataire (qui doit postuler ?), le poste, les missions, le lieu, le salaire et la date limite. Que remarquez-vous sur la mise en page de ce texte fonctionnel ?
\item Observez l'affiche : décrivez l'image, les couleurs et le slogan. Quel message l'image ajoute-t-elle au texte ? Pourquoi dit-on qu'un texte iconique et un texte écrit se complètent ?
\item Relevez dans les trois offres deux mots empruntés à une langue étrangère (par exemple à l'anglais) et proposez pour chacun un équivalent français.
\end{enumerate}

\textbf{Deuxième phase : rédaction individuelle.}

\begin{enumerate}
\setcounter{enumi}{3}
\item Rédigez l'en-tête complet de la lettre : vos coordonnées, celles du destinataire, le lieu et la date, l'objet de la lettre.
\item Rédigez le corps de la lettre en trois paragraphes selon le schéma \emph{vous -- moi -- nous}. Vous utiliserez \textbf{au moins deux phrases complexes} : une avec une subordonnée relative, une avec une subordonnée conjonctive.
\item Terminez par la formule de politesse adaptée à un destinataire que vous ne connaissez pas personnellement, puis signez. Vérifiez que la tonalité reste sérieuse et respectueuse : aucune plaisanterie (comique), aucun pathos exagéré (dramatique).
\end{enumerate}

\textbf{Troisième phase : confrontation et amélioration.}

\begin{enumerate}
\setcounter{enumi}{6}
\item Échangez votre lettre avec celle d'un camarade qui a choisi une offre différente. Évaluez sa production à l'aide de la grille critériée fournie (en-tête complet sur 4 points ; schéma vous--moi--nous sur 6 points ; phrases complexes sur 4 points ; formule de conclusion sur 2 points ; orthographe et présentation sur 4 points ; total sur 20).
\item Réécrivez votre version définitive en tenant compte des remarques reçues, puis justifiez \textbf{par écrit deux améliorations} que vous avez apportées (par exemple : « J'ai remplacé une phrase simple par une phrase complexe pour mieux relier mes compétences au poste »).
\end{enumerate}

\courssub{Ressources à mobiliser}

\begin{aretenir}[Les outils de la leçon à réutiliser]
\begin{itemize}
\item \textbf{Structure de la lettre officielle} : coordonnées de l'expéditeur en haut à gauche, coordonnées du destinataire à droite, lieu et date, objet, formule d'appel (\emph{Monsieur le Directeur,} / \emph{Madame, Monsieur,}), corps, formule de politesse, signature.
\item \textbf{Schéma du corps} : \emph{vous} (montrer qu'on connaît l'entreprise et le poste) ; \emph{moi} (présenter son diplôme, ses qualités, son expérience) ; \emph{nous} (proposer une collaboration, demander un entretien).
\item \textbf{Phrase complexe} : au moins deux propositions ; subordonnée relative introduite par \emph{qui, que, où, dont} ; subordonnée conjonctive introduite par \emph{que, parce que, puisque, afin que}.
\item \textbf{Emprunt} : mot venant d'une langue étrangère (\emph{job, e-mail, kit}) ; en lettre officielle, on préfère l'équivalent français (\emph{emploi, courriel, équipement}).
\item \textbf{Tonalité} : la lettre de motivation exige un ton soutenu, courtois et confiant ; les tonalités comique et dramatique sont ici hors de propos.
\item \textbf{Texte fonctionnel} : texte utilitaire à mise en page aérée (titres, listes, informations chiffrées) qui vise à informer et à faire agir.
\end{itemize}
\end{aretenir}

\courssub{Démarche et corrigé commenté}

\begin{exoresolu}[Rédaction complète d'une lettre de motivation (offre n° 3, Garoua)]
Énoncé : à partir de l'offre de NORDENERGIE, rédiger la lettre de motivation définitive d'un candidat fictif, puis justifier deux améliorations apportées après évaluation croisée.
\tcblower
\textbf{Étape 1 : analyse de l'offre (texte fonctionnel).}
L'expéditeur est la société NORDENERGIE ; le destinataire est un électrotechnicien débutant ; le poste porte sur la pose et la maintenance de kits solaires ; le lieu est Garoua et le Mayo-Rey ; le salaire est de \num{130000} FCFA plus \num{15000} FCFA de prime par mission ; le dossier va au Chef du personnel, BP 87, Garoua. La mise en page (titre en gras, rubriques \emph{Profil}, \emph{Missions}, \emph{Salaire}) permet une lecture rapide : c'est la marque du texte fonctionnel.

\textbf{Étape 2 : analyse de l'affiche (texte iconique).}
L'image du jeune diplômé serrant la main de l'employeur symbolise l'embauche ; les couleurs nationales valorisent le travail au Cameroun ; le slogan résume le message en trois étapes : diplôme, métier, avenir. L'image motive, le texte informe : les deux se complètent.

\textbf{Étape 3 : repérage des emprunts.}
Dans les offres, on relève \emph{« job »} (offre n° 1, anglais) et \emph{« kit »} (offre n° 3, anglais) ; en lettre officielle, on préférera \emph{« emploi »} (ou \emph{poste}) et \emph{« équipement solaire »} (ou \emph{installation}).

\textbf{Étape 4 : rédaction de la lettre.}

\begin{quote}
Amadou BELLO\\
Quartier Plateau, rue des Artisans\\
Tél. : 6 90 00 00 00\\
Garoua\\[4pt]
\hspace*{8cm}À Monsieur le Chef du personnel\\
\hspace*{8cm}NORDENERGIE, BP 87, Garoua\\[4pt]
\hspace*{8cm}Garoua, le 12 juin 2026\\[4pt]
\textbf{Objet :} candidature au poste d'électrotechnicien\\[4pt]
Monsieur le Chef du personnel,\\[4pt]
Votre société, qui équipe les villages du Mayo-Rey en énergie solaire, recherche un électrotechnicien disponible pour les déplacements. Cette mission, qui associe technique et service aux populations, correspond exactement au métier que je souhaite exercer.\\[4pt]
Titulaire du Baccalauréat technique F3 obtenu au Lycée technique de Garoua, j'ai effectué un stage de trois mois dans un atelier d'installation électrique, où j'ai appris à poser et à entretenir des équipements solaires. Je suis rigoureux, endurant, et je possède le permis B, parce que je savais que ce poste exigeait de fréquents déplacements.\\[4pt]
Intégrer votre équipe me permettrait de mettre mes compétences au service du développement de notre région, afin que davantage de villages aient accès à l'électricité. Je me tiens à votre disposition pour un entretien à la date qui vous conviendra.\\[4pt]
Je vous prie d'agréer, Monsieur le Chef du personnel, l'expression de ma considération distinguée.\\[4pt]
\hspace*{8cm}Amadou BELLO (signature)
\end{quote}

\textbf{Étape 5 : vérification des critères.}
L'en-tête est complet (expéditeur, destinataire, lieu et date, objet, appel). Le corps suit le schéma \emph{vous -- moi -- nous}. Deux phrases complexes sont présentes : une subordonnée relative (\emph{« qui équipe les villages du Mayo-Rey »}) et deux subordonnées conjonctives (\emph{« parce que je savais que... »}, \emph{« afin que davantage de villages aient accès... »}). La formule de politesse est adaptée à un supérieur hiérarchique inconnu. Le ton reste courtois et confiant, sans comique ni dramatique.

\textbf{Étape 6 : justification de deux améliorations (après évaluation croisée).}
\begin{enumerate}
\item \emph{Amélioration 1} : dans ma première version, j'avais écrit « Je veux ce job » ; je l'ai remplacé par « le métier que je souhaite exercer » afin d'éviter l'emprunt à l'anglais et de soutenir le niveau de langue exigé par une lettre officielle.
\item \emph{Amélioration 2} : j'avais écrit deux phrases simples (« J'ai le permis B. Ce poste exige des déplacements. ») ; je les ai fondues en une phrase complexe avec \emph{parce que}, ce qui montre la logique de ma candidature et enrichit la syntaxe.
\end{enumerate}
\end{exoresolu}

\courssub{Bilan de l'activité}

Cette activité a permis de mettre en évidence plusieurs acquis essentiels de la leçon. D'abord, la lecture croisée d'un \cle{texte fonctionnel} et d'un \cle{texte iconique} montre que l'information professionnelle circule à la fois par l'écrit (données chiffrées, rubriques) et par l'image (slogan, symboles), y compris dans la communication par internet où l'offre d'emploi est souvent accompagnée d'un visuel. Ensuite, la rédaction a imposé le respect strict des normes de la \cle{lettre de motivation} : un en-tête complet, un corps construit sur le schéma \emph{vous -- moi -- nous}, une formule de conclusion adaptée au destinataire. La grammaire a été mobilisée de façon concrète : la \cle{phrase complexe} n'est pas un exercice scolaire abstrait, mais un outil qui relie les idées et donne de la force à l'argumentation ; le choix du vocabulaire a montré que les \cle{emprunts}, fréquents dans le langage professionnel, doivent être maîtrisés et souvent remplacés par leurs équivalents français dans un écrit officiel. Enfin, l'évaluation croisée par grille critériée a développé l'autonomie : chaque élève a appris à lire la production d'autrui avec des critères objectifs, à accepter la critique et à \emph{justifier} ses propres choix de réécriture — compétence directement utile à l'épreuve de production d'écrits du Probatoire et du Baccalauréat technique, comme à la vie professionnelle réelle.

\newpage

\coursec{Méthodes et exercices résolus}

\coursec{Production d'écrits utilitaires : la lettre de motivation}

\courssub{Savoir-faire 1 : Rédiger l'en-tête d'une lettre de motivation}

\begin{definition}[Lettre de motivation]
La \cle{lettre de motivation} est un texte fonctionnel à visée professionnelle qui accompagne un curriculum vitae. Elle vise à convaincre un destinataire (employeur, responsable de formation) de l'adéquation entre le profil du candidat et le poste ou la formation visés. Elle obéit à des normes strictes de présentation (en-tête, corps, conclusion) et de registre (soutenu, vouvoiement).
\end{definition}

\begin{propriete}[Structure canonique de l'en-tête]
L'en-tête suit une disposition spatiale codifiée (norme AFNOR / usage administratif camerounais) :
\begin{itemize}
\item \textbf{Zone gauche (haut) :} Coordonnées de l'expéditeur (identité, adresse, téléphone, courriel).
\item \textbf{Zone droite (haut) :} Coordonnées du destinataire (fonction, organisme, adresse).
\item \textbf{Zone droite (sous destinataire) :} Lieu et date (ex. : « Fait à Yaoundé, le 15 mars 2026 »).
\item \textbf{Zone gauche (sous expéditeur) :} Références (facultatif) et Objet (obligatoire, phrase nominale, gras ou souligné).
\item \textbf{Zone gauche (sous objet) :} Formule d'appel (ex. : « Madame la Directrice, »).
\end{itemize}
\end{propriete}

\begin{methode}[Construire un en-tête conforme]
L'en-tête est la première partie lue par le destinataire : il doit être complet et rigoureusement disposé.
\begin{enumerate}
\item \textbf{Coordonnées de l'expéditeur} (en haut à gauche) : nom et prénom, adresse complète (boîte postale, ville), téléphone, adresse électronique.
\item \textbf{Coordonnées du destinataire} (en haut à droite, plus bas) : fonction (Monsieur le Directeur...), nom de l'entreprise ou de l'administration, adresse.
\item \textbf{Lieu et date} (à droite, sous le destinataire) : « Fait à Douala, le 15 mars 2026 ». Le mois s'écrit en toutes lettres, sans abréviation.
\item \textbf{Les références} éventuelles : « Réf. : votre annonce n° 12 du 10 mars 2026 ». Elles permettent d'identifier l'origine de la candidature.
\item \textbf{L'objet} (à gauche, en gras ou souligné) : une phrase nominale brève et précise : « Objet : demande de stage académique ». Pas de verbe conjugué, pas de point final.
\item \textbf{La formule d'appel} : « Monsieur le Directeur, » ou « Madame, Monsieur, » — jamais de prénom, jamais « Cher Monsieur » dans un courrier officiel. On reprend le titre exact de la fonction.
\end{enumerate}
\textbf{Réflexes} : vérifier que chaque élément est à sa place ; ne jamais mélanger les coordonnées ; l'objet annonce en une ligne le contenu de la lettre.
\end{methode}

\begin{savaistu}[Le courrier administratif au Cameroun]
L'administration camerounaise (fonction publique, établissements publics, grandes entreprises comme la SODECOTON, SONARA, CAMWATER) exige le respect scrupuleux de la présentation formelle. Une lettre mal présentée (en-tête incomplet, objet flou, formule d'appel incorrecte) est souvent écartée avant même la lecture du corps. Le \cle{Probatoire technique} évalue cette compétence en production d'écrits utilitaires.
\end{savaistu}

\begin{popfigure}
\begin{tikzpicture}[scale=0.85, every node/.style={font=\small}]
% En-tête schema
\draw[popblue, thick] (0,0) rectangle (14,7);
% Expéditeur
\node[align=left, anchor=north west, text width=6cm] at (0.3,6.7) {\textbf{EXPÉDITEUR (Gauche)}\newline Nom Prénom\newline Adresse (BP, Ville)\newline Tél. / Courriel};
% Destinataire
\node[align=left, anchor=north east, text width=6cm] at (13.7,6.7) {\textbf{DESTINATAIRE (Droite)}\newline Fonction (M. le Directeur...)\newline Organisme / Entreprise\newline Adresse (BP, Ville)};
% Lieu Date
\node[align=right, anchor=south east] at (13.7,4.5) {\textbf{Lieu, Date}\newline Fait à [Ville], le [Date]};
% Réf / Objet
\node[align=left, anchor=north west, text width=6cm] at (0.3,4.3) {\textbf{Références (si annonces)}\newline \textbf{Objet :} [Phrase nominale]};
% Appel
\node[align=left, anchor=south west] at (0.3,2.5) {\textbf{Formule d'appel}\newline Madame la Directeur,};
% Flèches de séparation
\draw[popline, dashed] (7,7) -- (7,0);
\draw[popline, dashed] (0,4.8) -- (14,4.8);
\end{tikzpicture}
\legende{Schéma de disposition de l'en-tête d'une lettre officielle (norme administrative).}
\end{popfigure}

\begin{exoresolu}[Rédiger un en-tête complet]
\textbf{Énoncé.} Tu es \textbf{Mbarga Alain}, élève en classe de Première technique au Lycée technique de Bafoussam (BP 45, Bafoussam, tél. 699 00 00 00, courriel : a.mbarga@mail.com). Tu réponds à l'annonce n° 8 parue le 2 avril 2026 pour un stage d'été à la \textbf{SODECOTON} de Garoua (BP 100, Garoua), dirigée par Madame la Directrice des ressources humaines. Rédige l'en-tête complet de ta lettre de motivation, datée du 10 avril 2026.
\tcblower
\textbf{Solution détaillée.}
\begin{enumerate}
\item \textbf{Expéditeur (haut gauche)} : on place d'abord l'identité, puis l'adresse, le téléphone et le courriel :
\begin{quote}
Mbarga Alain\\
Lycée technique de Bafoussam\\
BP 45, Bafoussam\\
Tél. : 699 00 00 00\\
a.mbarga@mail.com
\end{quote}
\item \textbf{Destinataire (haut droite)} : on indique la fonction exacte, l'entreprise et l'adresse :
\begin{quote}
Madame la Directrice des ressources humaines\\
SODECOTON\\
BP 100, Garoua
\end{quote}
\item \textbf{Lieu et date (à droite)} : « Fait à Bafoussam, le 10 avril 2026 ». On écrit le mois en toutes lettres, sans abréviation.
\item \textbf{Références et objet (à gauche)} :
\begin{quote}
Réf. : votre annonce n° 8 du 2 avril 2026\\
\textbf{Objet : demande de stage d'été}
\end{quote}
L'objet est une phrase nominale : pas de verbe conjugué, pas de point final obligatoire.
\item \textbf{Formule d'appel} : « Madame la Directrice, » — conforme à la fonction du destinataire.
\end{enumerate}
\textbf{Conclusion.} L'en-tête est complet : expéditeur, destinataire, lieu et date, références, objet et appel sont présents et correctement placés. Cette présentation soignée conditionne la première impression du lecteur.
\end{exoresolu}

\courssub{Savoir-faire 2 : Construire le corps de la lettre de motivation}

\begin{definition}[Schéma « Vous – Moi – Nous »]
C'est la structure standard du corps d'une lettre de motivation. Elle organise l'argumentation en trois temps logiques :
\begin{itemize}
\item \textbf{Vous} (L'entreprise) : Preuve de connaissance de la structure et du poste.
\item \textbf{Moi} (Le candidat) : Adéquation profil/poste (formation, expériences, qualités prouvées).
\item \textbf{Nous} (La collaboration) : Projection dans l'avenir, demande d'entretien, disponibilité.
\end{itemize}
\end{definition}

\begin{methode}[Organiser le corps en trois paragraphes]
Le corps de la lettre suit le schéma classique « vous – moi – nous ».
\begin{enumerate}
\item \textbf{Paragraphe 1 — « Vous »} : montrer qu'on connaît l'entreprise ou l'administration (activités, valeurs, poste proposé) et préciser l'annonce à laquelle on répond. \textit{Amorce type :} « Votre entreprise, spécialisée dans..., » ; « Suite à votre annonce n°..., » .
\item \textbf{Paragraphe 2 — « Moi »} : présenter sa formation, ses diplômes, ses expériences et ses qualités \emph{en lien avec le poste}. Chaque qualité affirmée doit être appuyée par un fait concret (stage, projet, note, responsabilité). \textit{Amorce type :} « Actuellement élève en Première technique, je... » ; « Mon parcours m'a permis de... » .
\item \textbf{Paragraphe 3 — « Nous »} : exprimer sa motivation, proposer une collaboration concrète, solliciter un entretien et se tenir à disposition. \textit{Amorce type :} « Intégrer vos équipes me permettrait... » ; « Je me permets de solliciter un entretien... » .
\item \textbf{Style} : phrases claires, registre soutenu, verbes d'action (maîtriser, réaliser, contribuer, assurer), aucune faute d'orthographe, phrases complexes correctement construites.
\end{enumerate}
\end{methode}

\begin{exoresolu}[Rédiger le corps d'une lettre]
\textbf{Énoncé.} À la suite de l'en-tête rédigé dans l'exercice précédent, rédige le corps de la lettre de motivation de Mbarga Alain pour le stage d'été à la SODECOTON, en respectant le schéma « vous – moi – nous » (trois paragraphes).
\tcblower
\textbf{Solution détaillée.}
\begin{enumerate}
\item \textbf{Paragraphe « Vous »} : on cite l'annonce et on montre sa connaissance de l'entreprise :
\begin{quote}
Suite à votre annonce n° 8 parue le 2 avril 2026, je me permets de vous adresser ma candidature pour un stage d'été au sein de votre société. La SODECOTON, acteur majeur de la filière cotonnière au Cameroun, offre un cadre idéal pour découvrir le monde industriel et les enjeux de la transformation agricole.
\end{quote}
\item \textbf{Paragraphe « Moi »} : formation et qualités justifiées par des faits :
\begin{quote}
Actuellement élève en classe de Première technique au Lycée technique de Bafoussam, je me prépare au Probatoire technique. Sérieux et ponctuel, j'ai obtenu d'excellents résultats en atelier de maintenance et j'ai participé activement à la révision des machines-outils de l'établissement. Je maîtrise également l'outil informatique (suite bureautique), compétence acquise lors de nos séances de travaux pratiques.
\end{quote}
Chaque qualité (sérieux, ponctualité, compétence technique, informatique) est appuyée par un élément concret (résultats, participation, séances TP).
\item \textbf{Paragraphe « Nous »} : motivation et demande d'entretien :
\begin{quote}
Intégrer vos équipes me permettrait de mettre en pratique mes connaissances théoriques tout en contribuant, à mon niveau, aux activités de maintenance de votre unité de production. Je reste à votre entière disposition pour un entretien à votre convenance.
\end{quote}
\end{enumerate}
\textbf{Conclusion.} Le corps respecte le schéma « vous – moi – nous » : trois paragraphes courts, un registre soutenu, des qualités prouvées et une demande claire d'entretien.
\end{exoresolu}

\courssub{Savoir-faire 3 : Employer les formules de conclusion}

\begin{propriete}[Règles de la formule de politesse finale]
\begin{enumerate}
\item \textbf{Concordance obligatoire} : La formule de politesse finale doit reprendre \textbf{exactement} le titre utilisé dans la formule d'appel (ex. : Appel « Madame la Directrice » $\rightarrow$ Conclusion « Madame la Directrice »).
\item \textbf{Formules consacrées (registre soutenu)} :
\begin{itemize}
\item Standard : « Je vous prie d'agréer, [Titre], l'expression de ma considération distinguée. »
\item Hiérarchie marquée : « Je vous prie d'agréer, [Titre], l'expression de mes sentiments respectueux. »
\item Variante : « Je vous prie de croire, [Titre], en l'assurance de ma considération distinguée. »
\end{itemize}
\item \textbf{Orthographe clé} : \cle{agréer} (deux \textbf{e} : a-g-r-é-e-r), \cle{considération} (féminin singulier $\rightarrow$ distingu\textbf{ée}).
\item \textbf{Interdits} : « Cordialement », « Bien à vous », « Amicalement », formules anglaises (« Best regards », « Sincerely »).
\end{enumerate}
\end{propriete}

\begin{methode}[Choisir et rédiger la formule de conclusion]
\begin{enumerate}
\item \textbf{Adapter la formule au destinataire} : reprendre dans la formule finale le titre exact de la formule d'appel.
\item \textbf{Rédiger la formule complète} : ajouter une phrase de transition courtoise (ex. : « Dans l'attente d'une réponse favorable, » ou « En espérant que ma candidature retiendra votre attention, »).
\item \textbf{Placer la signature} : à droite, sous la formule de politesse ; on signe, puis on écrit son nom et prénom lisiblement en dessous.
\item \textbf{Ajouter les pièces jointes} si nécessaire : « PJ : curriculum vitae, photocopie des relevés de notes ». La mention PJ se place à gauche, sous la signature.
\end{enumerate}
\end{methode}

\begin{exoresolu}[Conclure correctement la lettre]
\textbf{Énoncé.} Rédige la formule de conclusion complète de la lettre de Mbarga Alain à Madame la Directrice des ressources humaines de la SODECOTON, sachant qu'il joint son curriculum vitae et son relevé de notes.
\tcblower
\textbf{Solution détaillée.}
\begin{enumerate}
\item \textbf{Reprise du titre} : l'appel était « Madame la Directrice », donc la formule finale reprend exactement ce titre.
\item \textbf{Formule de politesse} :
\begin{quote}
Dans l'attente d'une réponse favorable, je vous prie d'agréer, Madame la Directrice, l'expression de ma considération distinguée.
\end{quote}
On vérifie l'orthographe : « agréer » (deux « e »), « considération distinguée » (accord avec « considération », féminin singulier).
\item \textbf{Signature} : à droite, le nom lisible :
\begin{flushright}
Mbarga Alain
\end{flushright}
\item \textbf{Pièces jointes} (à gauche) :
\begin{quote}
PJ : curriculum vitae ; relevé de notes de Première technique.
\end{quote}
\end{enumerate}
\textbf{Conclusion.} La conclusion est conforme : formule consacrée adaptée au destinataire, signature à droite et mention des pièces jointes. La lettre est prête à être envoyée.
\end{exoresolu}

\courssub{Savoir-faire 4 : Maîtriser la phrase complexe et l'emprunt}

\begin{definition}[Phrase complexe]
Une \cle{phrase complexe} est une phrase qui contient au moins deux verbes conjugués, donc au moins deux propositions. Elle s'oppose à la phrase simple (une seule proposition). Les propositions sont liées par :
\begin{itemize}
\item \textbf{La coordination} (mais, ou, et, donc, or, ni, car) : propositions indépendantes.
\item \textbf{La juxtaposition} (ponctuation forte : ; : ,) : propositions indépendantes.
\item \textbf{La subordination} (mots subordonnants : \emph{que, qui, parce que, si, lorsque, bien que, comme, puisque...}) : proposition dépendante (subordonnée) et proposition principale.
\end{itemize}
\end{definition}

\begin{definition}[Emprunt lexical]
Un \cle{emprunt} est un mot intégré au lexique français en provenance d'une langue étrangère (anglais, arabe, langues locales, etc.) sans adaptation morphologique initiale. Ex. : \emph{week-end, computer, business, marketing, software, leader, stress, cool, check-list}.
\end{definition}

\begin{propriete}[Types de propositions subordonnées (exigibles au Probatoire)]
\begin{tabularx}{\linewidth}{|l|X|X|}\hline
\textbf{Type} & \textbf{Mot subordonnant typique} & \textbf{Fonction / Question} \\ \hline
Relative & \emph{qui, que, où, dont, lequel...} & Complément du nom antécédent (remplaçable par « qui/que ») \\ \hline
Conjonctive (complétive) & \emph{que} (après verbe : dire, penser, vouloir, souhaiter...) & Objet direct de la principale (répond à « quoi ? ») \\ \hline
Circonstancielle de cause & \emph{parce que, comme, puisque, car (coord.)} & Répond à « pourquoi ? » \\ \hline
Circonstancielle de temps & \emph{quand, lorsque, dès que, avant que, après que} & Répond à « quand ? » \\ \hline
Circonstancielle de condition & \emph{si} (indicatif) / \emph{si} (imparfait + conditionnel) & Répond à « sous quelle condition ? » \\ \hline
Circonstancielle de concession & \emph{bien que, quoique, alors que} (+ subjonctif) & Répond à « malgré quoi ? » \\ \hline
Circonstancielle de but & \emph{afin que, pour que, de peur que} (+ subjonctif) & Répond à « dans quel but ? » \\ \hline
\end{tabularx}
\end{propriete}

\begin{methode}[Construire et analyser une phrase complexe / Gérer les emprunts]
\begin{enumerate}
\item \textbf{Identifier les propositions} : repérer tous les verbes conjugués. Compter les propositions = nombre de verbes conjugués.
\item \textbf{Distinguer les liens} : chercher les mots-outils (conjonctions, pronoms relatifs, ponctuation).
\item \textbf{Classifier les subordonnées} : identifier le mot subordonnant $\rightarrow$ déduire le type (relative, conjonctive, circonstancielle) $\rightarrow$ préciser la fonction (cause, temps, condition...).
\item \textbf{Test de la relative} : on peut supprimer la proposition relative ; la phrase principale reste grammaticalement correcte (ex. : « L'entreprise [qui est réputée] recrute » $\rightarrow$ « L'entreprise recrute »).
\item \textbf{Traiter les emprunts en lettre officielle} : identifier l'emprunt $\rightarrow$ chercher l'équivalent français standard $\rightarrow$ substituer.
\begin{itemize}
\item \emph{Computer} $\rightarrow$ \emph{ordinateur} / \emph{l'outil informatique} / \emph{les outils numériques}.
\item \emph{Week-end} $\rightarrow$ \emph{fin de semaine} / \emph{samedi et dimanche}.
\item \emph{Business} $\rightarrow$ \emph{affaires} / \emph{activité commerciale}.
\item \emph{Marketing} $\rightarrow$ \emph{mercatique} (terme officiel) / \emph{stratégie commerciale}.
\item \emph{Leader} $\rightarrow$ \emph{dirigeant} / \emph{responsable} / \emph{chef de file}.
\end{itemize}
\end{enumerate}
\textbf{Réflexe} : dans un courrier officiel, le registre soutenu impose le vocabulaire français standard. L'emprunt non adapté est une faute de registre.
\end{methode}

\begin{exoresolu}[Analyser et corriger des phrases]
\textbf{Énoncé.} 1) Analyse la phrase : « Je vous adresse ma candidature parce que je souhaite effectuer un stage dans votre entreprise qui est réputée pour son sérieux. » 2) Corrige la phrase suivante pour un registre officiel : « Je maîtrise le computer et je suis disponible chaque week-end. »
\tcblower
\textbf{Solution détaillée.}
\begin{enumerate}
\item \textbf{Dénombrement des propositions} : trois verbes conjugués : « adresse » (1), « souhaite » (2), « est » (3). Il y a donc trois propositions : c'est bien une phrase complexe.
\item \textbf{Analyse structurelle} :
\begin{itemize}
\item \textbf{P1 :} « Je vous adresse ma candidature » $\rightarrow$ \textbf{Proposition principale} (porte le sens central).
\item \textbf{P2 :} « parce que je souhaite effectuer un stage dans votre entreprise » $\rightarrow$ \textbf{Proposition subordonnée circonstancielle de cause}, introduite par \emph{parce que} (locution conjonctive). Elle répond à la question « pourquoi ? » par rapport à la principale.
\item \textbf{P3 :} « qui est réputée pour son sérieux » $\rightarrow$ \textbf{Proposition subordonnée relative}, introduite par le pronom relatif \emph{qui} (sujet de « est »). Son antécédent est « entreprise ». \textit{Test :} « Je vous adresse ma candidature parce que je souhaite effectuer un stage dans votre entreprise » (phrase correcte sans la relative).
\end{itemize}
\item \textbf{Correction des emprunts (registre officiel)} :
\begin{itemize}
\item \emph{computer} (emprunt à l'anglais \textit{computer}) $\rightarrow$ \textbf{je maîtrise l'outil informatique} (ou \emph{l'ordinateur}).
\item \emph{week-end} (emprunt à l'anglais \textit{week-end}, toléré à l'oral mais déconseillé à l'écrit officiel) $\rightarrow$ \textbf{chaque fin de semaine} (ou \emph{les samedis et dimanches}).
\end{itemize}
Phrase corrigée :
\begin{quote}
« Je maîtrise l'outil informatique et je suis disponible chaque fin de semaine. »
\end{quote}
\end{enumerate}
\textbf{Conclusion.} La phrase analysée comporte une principale, une subordonnée circonstancielle de cause et une subordonnée relative. La correction a remplacé deux emprunts à l'anglais par leurs équivalents français, exigence du registre officiel au Probatoire.
\end{exoresolu}

\courssub{Savoir-faire 5 : Identifier les tonalités et lire un texte iconique}

\begin{definition}[Tonalité (ou ton)]
La \cle{tonalité} est l'attitude affective ou intellectuelle que l'auteur adopte vis-à-vis de son sujet et de son lecteur. Elle se manifeste par des procédés lexicaux, syntaxiques, phoniques et rhétoriques. Au programme : \textbf{dramatique} et \textbf{comique}.
\end{definition}

\begin{propriete}[Procédés des tonalités dramatiques et comiques]
\begin{tabularx}{\linewidth}{|l|X|X|}\hline
\textbf{Tonalité} & \textbf{Procédés principaux} & \textbf{Effet visé} \\ \hline
\textbf{Dramatique} & \begin{itemize}\item Champ lexical : mort, douleur, sang, nuit, destin, fatalité, souffrance.\item Figures : gradation, hyperbole, antithèse, métaphore filée.\item Syntaxique : phrases nominales, exclamatives, interrogatives rhétoriques, ellipses, rythme saccadé ou oppressant.\item Ponctuation : points de suspension, points d'exclamation, tirets.\end{itemize} & Susciter l'émotion forte (pitié, terreur, angoisse), montrer la gravité d'une situation, l'impuissance face au destin. \\ \hline
\textbf{Comique} & \begin{itemize}\item \textbf{Comique de mots :} calembour, paronymie, contrepèterie, jeu sur la polysémie.\item \textbf{Comique de gestes :} gestuelle exagérée, chute, mimique, répétition mécanique.\item \textbf{Comique de situation :} quiproquo, inversion des rôles, répétition, décalage entre apparence et réalité, malentendu.\item \textbf{Comique de caractère :} exagération d'un trait (avare, vaniteux, pédant).\end{itemize} & Provoquer le rire ou le sourire, critiquer par le ridicule, détendre l'atmosphère, dénoncer les travers humains. \\ \hline
\end{tabularx}
\end{propriete}

\begin{methode}[Analyser un texte iconique (affiche, photo, caricature, dessin de presse)]
\begin{enumerate}
\item \textbf{Description objective (le « quoi »)} : identifier la nature du document (affiche, photo...), les éléments visuels (personnages, objets, décor, couleurs dominantes, composition : premier plan / arrière-plan), le texte associé (slogan, légende, bulles, titre), la source si connue.
\item \textbf{Interprétation (le « comment » et « pourquoi »)} : analyser le rapport texte/image (redondance, complémentarité, contradiction), identifier le message principal (informer, alerter, dénoncer, vendre, convaincre), déterminer l'intention de l'auteur (cible visée, appel à l'action).
\item \textbf{Identification de la tonalité} : annoncer la tonalité (dramatique, comique, ironique, pathétique...), citer \textbf{au moins deux procédés} précis (lexicaux, visuels, syntaxiques) qui la fondent.
\item \textbf{Rédaction de la réponse} : structurer en paragraphes (Description $\rightarrow$ Interprétation $\rightarrow$ Tonalité/Procédés).
\end{enumerate}
\end{methode}

\begin{savaistu}[L'affiche de sensibilisation au Cameroun]
Les campagnes de santé publique (paludisme, VIH/SIDA, vaccination, hygiène) ou de sécurité routière au Cameroun utilisent massivement l'affichage urbain. Elles combinent souvent une \textbf{tonalité dramatique} (images chocs : enfant malade, accident, cercueil) et un \textbf{slogan impératif} court pour provoquer un changement de comportement immédiat. L'analyse de ces textes iconiques est une compétence évaluée en réception des textes.
\end{savaistu}

\begin{exoresolu}[Analyser une affiche et un extrait]
\textbf{Énoncé.} 1) Une affiche de campagne contre le paludisme montre, sur fond sombre, un enfant endormi sous une moustiquaire déchirée, avec la légende : « Une nuit sans moustiquaire peut tout changer. » Analyse ce texte iconique et précise sa tonalité. 2) Dans une scène de théâtre, un personnage répète trois fois « Je ne suis pas malade ! » tout en éternuant à chaque réplique. Quelle tonalité et quels procédés ?
\tcblower
\textbf{Solution détaillée.}
\begin{enumerate}
\item \textbf{Description de l'affiche} :
\begin{itemize}
\item \textit{Visuel :} Fond sombre (couleurs froides : bleu nuit, noir), premier plan : un enfant vulnérable, endormi. Objet central : une moustiquaire déchirée (protection défaillante). Contraste fort entre l'innocence du sommeil et le danger visible (trous).
\item \textit{Textuel :} Légende unique, police blanche/jaune sur fond sombre : « Une nuit sans moustiquaire peut tout changer. » Conditionnel (\emph{peut}) + portée universelle (\emph{tout changer}).
\end{itemize}
\item \textbf{Interprétation} : Le message est une mise en garde sanitaire : la moustiquaire imprégnée (MILDA) est la protection principale. L'image et le texte se renforcent mutuellement (ancrage du sens). L'intention est persuasive : toucher le lecteur (parent, chef de famille) pour provoquer l'achat/utilisation correcte de la moustiquaire.
\item \textbf{Tonalité :} \textbf{Dramatique}.
\textbf{Procédés justificatifs} :
\begin{enumerate}
\item \textbf{Visuel / Chromatique :} Fond sombre, couleurs froides $\rightarrow$ atmosphère angoissante, menaçante (champ lexical visuel du danger, de la nuit).
\item \textbf{Lexical / Sémantique :} « déchirée » (violence, rupture), « tout changer » (hyperbole, conséquence irréversible : la mort), « nuit » (temporalité du danger, champ lexical de l'obscurité/mort).
\item \textbf{Syntaxique :} Phrase à l'indicatif présent (\emph{peut}) mais valeur prédictive menaçante ; structure nominale implicite « Une nuit [sans protection] = catastrophe ».
\end{enumerate}
\item \textbf{Scène de théâtre} : La tonalité est \textbf{comique}.
\textbf{Procédés} :
\begin{enumerate}
\item \textbf{Comique de situation / Décalage :} Le personnage affirme le contraire de ce que son corps manifeste (éternuements = symptôme évident). Ironie dramatique (le public sait, le personnage nie).
\item \textbf{Comique de gestes :} Les éternuements répétés (répétition mécanique, involontaire) contredisent la parole volontaire.
\item \textbf{Comique de répétition / Litanie :} La réplique dite trois fois renforce l'absurdité du déni.
\end{enumerate}
\end{enumerate}
\textbf{Conclusion.} L'affiche relève de la tonalité dramatique (menace, contraste visuel, hyperbole lexicale, couleurs sombres) tandis que la scène relève du comique de situation, de gestes et de répétition : dans les deux cas, la tonalité est justifiée par des procédés précis et identifiables.
\end{exoresolu}

\courssub{Savoir-faire 6 : Confronter et améliorer une production (Remédiation)}

\begin{methode}[Réviser une lettre de motivation en grille d'évaluation]
La confrontation (échange entre pairs) et la remédiation (correction individuelle) sont les étapes finales de la production d'écrits.
\begin{enumerate}
\item \textbf{Relecture formelle (Présentation - 4 pts)} : vérifier l'en-tête (places respectées : expéditeur G, destinataire D, lieu/date D, références/objet G, appel G), l'objet (phrase nominale), l'appel (titre exact), la formule de politesse (concordance titre, orthographe \emph{agréer}, \emph{distinguée}), la signature (Droite), les PJ (Gauche).
\item \textbf{Relecture du contenu (Argumentation - 6 pts)} : le schéma « vous – moi – nous » est-il respecté (3 paragraphes distincts) ? Chaque qualité est-elle prouvée par un fait (exemple, chiffre, expérience) ? La demande (entretien) est-elle explicite et courtoise ? L'annonce est-elle citée ?
\item \textbf{Relecture linguistique (Langue - 6 pts)} : orthographe (accords, homophones), grammaire (phrases complexes correctes, subordonnants variés), vocabulaire (registre soutenu, \textbf{pas d'emprunts inutiles}, pas de codes SMS : « slt », « svp », « j'attends ta réponse », « cdt »).
\item \textbf{Confrontation par les pairs} : échanger sa lettre avec un camarade. Utiliser une grille d'évaluation partagée (items ci-dessus). Noter les remarques constructives (points forts / points à améliorer).
\item \textbf{Réécriture (Version définitive)} : intégrer les corrections pertinentes. Produire la version finale propre.
\item \textbf{Compte rendu et remédiation} : noter dans son cahier les erreurs types corrigées (ex. : « J'ai remplacé \emph{computer} par \emph{outil informatique} », « J'ai corrigé l'accord de \emph{distinguée} ») pour ne plus les reproduire.
\end{enumerate}
\end{methode}

\begin{aretenir}[Codes de la communication par internet vs Lettre officielle]
\begin{center}
\begin{tabularx}{\linewidth}{|l|X|X|}\hline
\textbf{Code Internet (Interdit)} & \textbf{Équivalent Lettre Officielle (Exigé)} & \textbf{Règle} \\ \hline
Abréviations : \emph{slt, svp, bjr, cdt, pq, qqch} & Formules complètes : \emph{s'il vous plaît, bienvenue, cordialement (interdit ici), pourquoi, quelque chose} & \textbf{Zéro abréviation SMS} \\ \hline
Tutoiement (\emph{ton, ta, tu}) & Vouvoiement systématique (\emph{votre, votre, vous}) & \textbf{Vouvoiement obligatoire} \\ \hline
Émoticônes (\emph{:), ;), :D}) & Aucune image, aucun symbole graphique & \textbf{Texte pur} \\ \hline
Registre familier (\emph{je kiffe, c'est nul, genre}) & Registre soutenu / standard (\emph{j'apprécie, c'est insuffisant, notamment}) & \textbf{Registre soutenu} \\ \hline
Ponctuation expressive (\emph{!!!, ???, ...}) & Ponctuation standard (., ; :) & \textbf{Sobriété} \\ \hline
\end{tabularx}
\end{center}
\end{aretenir}

\begin{exoresolu}[Corriger une lettre défaillante]
\textbf{Énoncé.} Voici un extrait de lettre rédigé par un élève : « Slt Monsieur le Directeur, j'ai vu ton annonce sur internet. Je suis fort en computer et je veux un stage chez vous. Réponds-moi vite. Cordialement. » Releve cinq erreurs majeures et réécris cet extrait en version correcte.
\tcblower
\textbf{Solution détaillée.}
\begin{enumerate}
\item \textbf{Relevé des 5 erreurs majeures} :
\begin{itemize}
\item \textbf{Erreur 1 (Appel / Code SMS) :} « Slt » (abréviation de \emph{salut}) $\rightarrow$ \textbf{Formule d'appel complète} : « Monsieur le Directeur, ».
\item \textbf{Erreur 2 (Tutoiement) :} « ton annonce » $\rightarrow$ \textbf{Vouvoiement} : « votre annonce ».
\item \textbf{Erreur 3 (Emprunt + Registre familier) :} « fort en computer » $\rightarrow$ \textbf{Vocabulaire soutenu + Français} : « je maîtrise l'outil informatique » (ou « je possède de solides compétences en informatique »).
\item \textbf{Erreur 4 (Syntaxe / Tutoiement / Impératif) :} « Réponds-moi vite » $\rightarrow$ \textbf{Formule courtoise complexe} : « Dans l'attente de votre réponse, je vous prie de croire... » ou « Je me tiens à votre disposition pour tout entretien. ».
\item \textbf{Erreur 5 (Formule de politesse finale inadaptée) :} « Cordialement » (trop relâché, usage mail interne) $\rightarrow$ \textbf{Formule consacrée} : « Je vous prie d'agréer, Monsieur le Directeur, l'expression de ma considération distinguée. ».
\end{itemize}
\item \textbf{Version corrigée (Registre officiel)} :
\begin{quote}
Monsieur le Directeur,\\
J'ai pris connaissance de votre annonce publiée sur internet et je me permets de vous adresser ma candidature pour un stage au sein de votre entreprise. Je maîtrise l'outil informatique, compétence acquise au cours de ma formation en Première technique.\\
Dans l'attente d'une réponse favorable, je vous prie d'agréer, Monsieur le Directeur, l'expression de ma considération distinguée.
\end{quote}
\end{enumerate}
\textbf{Conclusion.} La version améliorée respecte le vouvoiement, le registre soutenu, le vocabulaire français standard et les formules consacrées : la confrontation et la relecture ont permis la remédiation.
\end{exoresolu}

\begin{exercice}[Entraînement complet : Lettre de motivation]
\textbf{Consigne.} Tu es \textbf{Ngo Bitjeck Marie}, élève en Première technique (spécialité Génie Civil) au Lycée technique de Douala-Bonabéri (BP 2000, Douala, Tél. 677 11 22 33, marie.ngo@mail.cm). Tu réponds à l'annonce n° 45 du 5 mai 2026 de l'entreprise \textbf{RAZEL-BEC} (BP 500, Douala), dirigée par \textbf{Monsieur le Directeur des Travaux}, pour un stage de fin d'année (juin-juillet).
\begin{enumerate}
\item Rédige l'\textbf{en-tête complet} de la lettre.
\item Rédige le \textbf{corps} de la lettre (schéma « Vous – Moi – Nous ») en valorisant ta spécialité (Génie Civil : lecture de plans, topographie, matériaux, chantiers-écoles).
\item Rédige la \textbf{conclusion} (formule de politesse, signature, PJ : CV + bulletins).
\item \textbf{Auto-évaluation :} Coche la grille : [ ] En-tête complet et bien placé [ ] Objet phrase nominale [ ] Appel = Titre conclusion [ ] 3 paragraphes Vous/Moi/Nous [ ] Qualités prouvées [ ] 0 emprunt / 0 SMS [ ] Orthographe \emph{agréer/distinguée} vérifiée.
\end{enumerate}
\end{exercice}

\begin{corrige}[Exercice n° 1 : Lettre de motivation Ngo Bitjeck]
\textbf{Proposition de correction (modèle).}
\begin{quote}
\textbf{Ngo Bitjeck Marie}\\
Lycée technique de Douala-Bonabéri\\
BP 2000, Douala\\
Tél. : 677 11 22 33\\
marie.ngo@mail.cm

\hfill \textbf{Monsieur le Directeur des Travaux}\\
\hfill \textbf{RAZEL-BEC}\\
\hfill \textbf{BP 500, Douala}

\hfill Fait à Douala, le 10 mai 2026

Réf. : votre annonce n° 45 du 5 mai 2026\\
\textbf{Objet : demande de stage de fin d'année (Génie Civil)}

\textbf{Monsieur le Directeur,}

\textbf{[Vous]} Suite à votre annonce n° 45 parue le 5 mai 2026, je vous adresse ma candidature pour un stage de fin d'année au sein de votre direction des travaux. Le groupe RAZEL-BEC, référence majeure dans les infrastructures routières et hydrauliques en Afrique centrale, représente pour moi l'environnement idéal pour confronter mes acquis scolaires à la réalité du chantier.

\textbf{[Moi]} Actuellement élève en Première technique, spécialité Génie Civil, au Lycée technique de Douala-Bonabéri, je prépare le Probatoire technique. Ma formation m'a permis d'acquérir des compétences solides en lecture de plans d'exécution, en topographie (nivellement, piquetage) et en technologie des matériaux (essais béton, granulats). Lors de mes stages en entreprise (chantier-école de l'établissement), j'ai fait preuve de rigueur dans le respect des normes de sécurité et de précision dans les relevés topographiques. Je maîtrise par ailleurs les logiciels de DAO (AutoCAD) utilisés en TP.

\textbf{[Nous]} Intégrer vos équipes serait une opportunité précieuse pour approfondir ma pratique du suivi de chantier et de la gestion des matériaux. Dynamique et respectueuse des consignes, je suis prête à m'investir pleinement. Je sollicite l'honneur d'un entretien à votre convenance pour vous exposer ma motivation.

Dans l'attente d'une réponse favorable, je vous prie d'agréer, Monsieur le Directeur des Travaux, l'expression de ma considération distinguée.

\hfill \textbf{Ngo Bitjeck Marie}

PJ : Curriculum vitae ; Bulletins de notes de Première technique ; Attestation de scolarité.
\end{quote}
\end{corrige}

\begin{attention}[Les 5 erreurs qui coûtent des points au Probatoire / Bac Tech]
\begin{enumerate}
\item \textbf{Mélanger les places de l'en-tête} : expéditeur à gauche, destinataire à droite ; inverser les deux ou oublier l'objet fait perdre des points de présentation (souvent 2 à 4 pts sur 20).
\item \textbf{Garder les codes d'internet dans la lettre} : abréviations (« svp », « slt », « cdt »), émoticônes, tutoiement : la lettre officielle exige le vouvoiement et le registre soutenu. \textbf{Sanction :} pénalité lourde sur la note de langue/registre.
\item \textbf{Affirmer une qualité sans la prouver} : « je suis sérieux », « je suis dynamique » seuls ne valent rien ; il faut un fait vérifiable : « j'ai été délégué de classe », « j'ai obtenu 16/20 en atelier », « j'ai géré le stock du club ». \textbf{Preuve = Points.}
\item \textbf{Écorcher la formule de politesse} : « agréé » (un seul « e »), « Cordialement » en conclusion, ou un titre différent de celui de l'appel (« Monsieur » à l'appel, « Madame la Directrice » en conclusion). \textbf{Concordance obligatoire.}
\item \textbf{Confondre les propositions dans la phrase complexe} : prendre une relative pour une conjonctive, ou oublier le mot subordonnant ; toujours vérifier en supprimant la subordonnée : la principale doit rester correcte. \textbf{Maîtrise syntaxique = Lisibilité.}
\end{enumerate}
\end{attention}

\newpage

\coursec{Exercices}

\coursec{Production d'écrits utilitaires : la lettre de motivation}

\courssub{Je vérifie mes connaissances}

\begin{exercice}[QCM : la lettre de motivation]
Pour chaque question, entoure la (ou les) bonne(s) réponse(s).
\begin{enumerate}
\item Une lettre de motivation s'écrit toujours :
\begin{enumerate}
\item[a)] au verso du curriculum vitae ;
\item[b)] sur une seule page, en réponse à une offre ou en candidature spontanée ;
\item[c)] en langage familier pour montrer qu'on est décontracté ;
\item[d)] avec une formule d'appel et une formule de politesse.
\end{enumerate}
\item L'en-tête d'une lettre de motivation comporte obligatoirement :
\begin{enumerate}
\item[a)] les coordonnées de l'expéditeur ;
\item[b)] la photo du candidat ;
\item[c)] les coordonnées du destinataire, le lieu et la date ;
\item[d)] l'objet de la lettre.
\end{enumerate}
\item La formule d'appel correcte pour s'adresser au directeur d'une entreprise est :
\begin{enumerate}
\item[a)] « Cher monsieur le boss » ;
\item[b)] « Monsieur le Directeur, » ;
\item[c)] « Salut Monsieur, » ;
\item[d)] « À qui de droit, ».
\end{enumerate}
\item Le corps de la lettre de motivation se compose classiquement de trois paragraphes :
\begin{enumerate}
\item[a)] « vous » – « moi » – « nous » ;
\item[b)] « moi » – « vous » – « eux » ;
\item[c)] « nous » – « vous » – « moi » ;
\item[d)] « eux » – « moi » – « vous ».
\end{enumerate}
\end{enumerate}
\end{exercice}

\begin{exercice}[Vrai ou faux justifié]
Dis si chaque affirmation est vraie ou fausse, puis justifie ta réponse en une ou deux phrases.
\begin{enumerate}
\item La phrase complexe contient au moins deux propositions.
\item Dans la phrase « Je vous adresse ma candidature afin que vous examiniez mon dossier », la subordonnée est une relative.
\item Le mot « job » est un emprunt à la langue anglaise.
\item La tonalité comique vise à émouvoir le lecteur jusqu'aux larmes.
\item Dans un courriel professionnel, il est acceptable d'écrire « slt » pour « salut ».
\end{enumerate}
\end{exercice}

\begin{exercice}[Définitions]
Définis précisément chacun des termes suivants en une ou deux phrases, puis donne un exemple.
\begin{enumerate}
\item La lettre de motivation.
\item La phrase complexe.
\item L'emprunt (en lexicologie).
\item La netiquette.
\item La tonalité dramatique.
\end{enumerate}
\end{exercice}

\begin{exercice}[Analyse grammaticale directe]
Pour chacune des phrases suivantes, indique : a) la nature de la subordonnée (relative, conjonctive, participiale) ; b) sa fonction (si relative : fonction dans la proposition ; si conjonctive : fonction de la proposition ou complément de conjonction ; si participiale : fonction circonstantielle).
\begin{enumerate}
\item « J'ai suivi une formation qui m'a permis de maîtriser la maintenance industrielle. »
\item « Je reste à votre disposition pour que nous convenions d'un entretien. »
\item « Titulaire d'un Baccalauréat technique, je souhaite mettre mes compétences à votre service. »
\item « L'offre que vous avez publiée correspond exactement à mon profil. »
\end{enumerate}
\end{exercice}

\courssub{Je m'entraîne}

\begin{exercice}[Rédiger l'en-tête]
Tu es Mballa Éric, demeurant au quartier Nkolndongo, BP 5123 Yaoundé, téléphone 6 90 12 34 56, adresse électronique mballa.eric@mail.cm. Tu réponds à une offre d'emploi de la société CAMBOIS SARL, située avenue Ahmadou Ahidjo, BP 208 Douala, pour un poste de technicien en maintenance. La référence de l'offre est « Offre n° 17/CB/2025 ». Rédige l'en-tête complet de ta lettre de motivation (coordonnées de l'expéditeur, coordonnées du destinataire, lieu et date, objet, références éventuelles).
\end{exercice}

\begin{exercice}[Construire des phrases complexes]
Transforme chaque groupe de phrases simples en une seule phrase complexe, selon la subordination demandée.
\begin{enumerate}
\item Subordination relative : « J'ai obtenu un Baccalauréat technique. Ce diplôme m'a donné de solides bases en électrotechnique. »
\item Subordination conjonctive de but : « Je me permets de vous écrire. Je veux vous proposer ma candidature. »
\item Proposition participiale : « Je suis motivé par le dynamisme de votre entreprise. Je postule au stage que vous proposez. »
\item Subordination conjonctive de conséquence : « Mon stage chez SODECOTON m'a beaucoup appris. Je peux désormais travailler en autonomie. »
\end{enumerate}
\end{exercice}

\begin{exercice}[Les emprunts dans un texte professionnel]
Lis le texte suivant : « Notre start-up recrute un community manager. Le candidat devra gérer le mailing, animer le chat de l'entreprise, mettre à jour le site web et assurer le suivi du feedback des clients. Une bonne maîtrise du smartphone et des logiciels de design est exigée. »
\begin{enumerate}
\item Relève cinq emprunts présents dans ce texte.
\item Pour chacun, indique la langue d'origine.
\item Propose pour chacun un équivalent français soutenu utilisable dans une lettre de motivation.
\end{enumerate}
\end{exercice}

\begin{exercice}[Formules de conclusion]
Voici quatre formules de politesse. Classe-les en deux colonnes : « correctes dans une lettre officielle » et « incorrectes dans une lettre officielle ». Corrige ensuite les formules incorrectes.
\begin{enumerate}
\item « Veuillez agréer, Monsieur le Directeur, l'expression de ma considération distinguée. »
\item « Salut et à bientôt j'espère ! »
\item « Dans l'attente d'une suite favorable, je vous prie de croire, Madame, à mes salutations distinguées. »
\item « Bisous et merci d'avance. »
\end{enumerate}
\end{exercice}

\courssub{Je me prépare au Bac}

\begin{exercice}[Type Baccalauréat : rédaction complète d'une lettre de motivation]
\textbf{Situation.} Tu lis dans \emph{Cameroon Tribune} du mardi 14 octobre 2025 l'offre d'emploi suivante :

\begin{center}
\begin{tabularx}{\linewidth}{|X|}
\hline
\rowcolor{popblueL} \textbf{LA SOCIÉTÉ AGROCAM SARL RECRUTE} \\
\hline
Un(e) technicien(ne) en agropastoral pour son site de Bafoussam. Profil : titulaire d'un Baccalauréat technique (série agropastorale) ou équivalent ; sens de l'organisation ; capacité à travailler en équipe. Dossier : lettre de motivation + CV, à adresser au Directeur des Ressources Humaines, AGROCAM SARL, BP 1450 Bafoussam, avant le 15 novembre 2025. Réf. : REC/AGRO/25/09. \\
\hline
\end{tabularx}
\end{center}

Tu es candidat(e) à ce poste. Tu t'appelleras Fotso Larissa, tu demeures à Bafoussam (BP 777), et tu es titulaire d'un Baccalauréat technique obtenu au Lycée technique de Bafoussam en juin 2025. Tu as effectué un stage de deux mois dans une exploitation agricole de Foumbot.

\textbf{Travail à faire :}
\begin{enumerate}
\item Rédige l'en-tête complet de ta lettre (coordonnées, lieu et date, destinataire, objet, référence). \textbf{(4 pts)}
\item Rédige le corps de la lettre en trois paragraphes : le paragraphe « vous » (l'entreprise et le poste), le paragraphe « moi » (ton profil, ta formation, ton expérience), le paragraphe « nous » (la collaboration envisagée, la demande d'entretien). Tu utiliseras au moins deux phrases complexes de types différents. \textbf{(8 pts)}
\item Rédige la formule d'appel et la formule de politesse finale adaptées. \textbf{(4 pts)}
\item Soigne l'orthographe, la syntaxe et la présentation. \textbf{(4 pts)}
\end{enumerate}
\end{exercice}

\begin{exercice}[Type Probatoire : réécriture d'une lettre fautive]
\textbf{Consigne.} La lettre ci-dessous, écrite par un élève de Première technique, contient \textbf{dix erreurs} : langage SMS, formule d'appel fautive, objet vague, emprunts inadaptés, formule de politesse familière, fautes de registre.

\begin{center}
\begin{tabularx}{\linewidth}{|X|}
\hline
\rowcolor{poppinkL} Lettre à corriger \\
\hline
Slt M. le boss,\\
Objet : truc\\
J'ai vu ton annonce sur le net et ça m'a grave branché. Je veux le job de technicien parce que je suis trop chaud pour bosser dans ta boîte. J'ai fait un stage dans une start-up où je gérais le mailing et le feedback des clients, c'était cool.\\
Appelle-moi vite pour qu'on se capte.\\
Biz,\\
Rodrigue\\
\hline
\end{tabularx}
\end{center}

\textbf{Travail à faire :}
\begin{enumerate}
\item Relève les dix erreurs et classe-les selon leur nature : registre de langue, formules protocolaires, emprunts inadaptés, objet de la lettre. \textbf{(5 pts)}
\item Réécris la formule d'appel et l'objet de manière correcte. \textbf{(2 pts)}
\item Réécris le paragraphe central en langue soutenue, en remplaçant les emprunts par des équivalents français et en construisant au moins une phrase complexe. \textbf{(8 pts)}
\item Propose une formule de politesse finale adaptée à un directeur d'entreprise. \textbf{(2 pts)}
\item Explique en trois lignes pourquoi la netiquette interdit le langage SMS dans une communication professionnelle par internet. \textbf{(3 pts)}
\end{enumerate}
\end{exercice}

\begin{exercice}[Type Bac : texte iconique et rédaction d'un courriel de candidature]
\textbf{Première partie : analyse d'un texte iconique.} Une affiche de recrutement publiée par une banque de Yaoundé montre, sur fond vert, un jeune diplômé en costume, souriant, gravissant une échelle géante dont le sommet porte l'inscription « VOTRE AVENIR COMMENCE ICI ». Au bas de l'échelle, une foule de candidats anonymes, dessinés en gris, attendent. En bas de l'affiche figurent le logo de la banque et la mention : « Envoyez votre lettre de motivation et votre CV avant le 30 novembre. »
\begin{enumerate}
\item Décris objectivement l'affiche (éléments visuels, couleurs, texte). \textbf{(3 pts)}
\item Dégage la tonalité dominante de cette affiche (comique ou dramatique) et identifie deux procédés qui la construisent (contraste des couleurs, symbolique de l'échelle, opposition individu/foule). \textbf{(4 pts)}
\item Explique en quelques lignes l'effet recherché sur le lecteur-candidat. \textbf{(3 pts)}
\end{enumerate}

\textbf{Deuxième partie : courriel de candidature.} Tu décides de postuler par voie électronique.
\begin{enumerate}
\setcounter{enumi}{3}
\item Rédige le courriel accompagnant ta lettre de motivation et ton CV placés en pièces jointes : objet du message, formule d'appel, corps bref et poli annonçant les pièces jointes, formule de politesse, signature. Tu respecteras la netiquette (pas de langage SMS, pas d'abréviations, ponctuation soignée). \textbf{(6 pts)}
\item Justifie en deux ou trois phrases le choix de l'objet de ton courriel. \textbf{(2 pts)}
\item Cite deux règles de la communication par internet que tu as appliquées dans ton message. \textbf{(2 pts)}
\end{enumerate}
\end{exercice}

\newpage

\coursec{Corrigés des exercices}

\coursec{Corrigés des exercices — Leçon 09 : Production d'une lettre de motivation}

\begin{corrige}[Exercice 1 — QCM : la lettre de motivation]
\textbf{Question 1 :} réponses \textbf{b)} et \textbf{d)}. Une lettre de motivation tient sur une seule page et accompagne le CV (elle ne s'écrit pas au verso) ; elle peut répondre à une offre ou constituer une candidature spontanée. Elle exige une \cle{formule d'appel} (« Monsieur le Directeur, ») et une \cle{formule de politesse} finale. Le langage familier (c) est proscrit : le registre soutenu est de rigueur.

\textbf{Question 2 :} réponses \textbf{a)}, \textbf{c)} et \textbf{d)}. L'\cle{en-tête} comporte les coordonnées de l'expéditeur (en haut à gauche), les coordonnées du destinataire (à droite), le lieu et la date, ainsi que l'\cle{objet} de la lettre. La photo (b) n'appartient pas à la lettre : elle figure éventuellement sur le CV, et n'est jamais obligatoire.

\textbf{Question 3 :} réponse \textbf{b)}. « Monsieur le Directeur, » est la formule d'appel protocolaire correcte. « Cher monsieur le boss » (a) et « Salut Monsieur » (c) relèvent du registre familier ; « À qui de droit » (d) est impersonnel et dévalorisant quand on connaît le destinataire.

\textbf{Question 4 :} réponse \textbf{a)}. Le plan classique du corps de la lettre est « \cle{vous} – \cle{moi} – \cle{nous} » : d'abord l'entreprise et le poste (vous), ensuite le profil du candidat (moi), enfin la collaboration envisagée et la demande d'entretien (nous).
\end{corrige}

\begin{corrige}[Exercice 2 — Vrai ou faux justifié]
\begin{enumerate}
\item \textbf{Vrai.} Par définition, la \cle{phrase complexe} contient au moins deux propositions : une proposition principale et au moins une proposition subordonnée (relative, conjonctive, participiale...). \emph{Attention :} les propositions juxtaposées ou coordonnées forment une \emph{phrase composée}, non une phrase complexe.
\item \textbf{Faux.} « afin que vous examiniez mon dossier » est une subordonnée \emph{conjonctive} introduite par la locution conjonctive « afin que » (suivie du subjonctif) ; elle exprime le but. Une relative serait introduite par un pronom relatif (qui, que, dont...).
\item \textbf{Vrai.} « Job » est un \cle{emprunt} à l'anglais (anglicisme), intégré au vocabulaire courant pour désigner un travail, un emploi.
\item \textbf{Faux.} C'est la \cle{tonalité dramatique} qui vise à émouvoir, à susciter la tristesse ou l'angoisse. La tonalité comique cherche à faire rire ou sourire.
\item \textbf{Faux.} La \cle{netiquette} (code de politesse de la communication par internet) interdit le langage SMS dans un courriel professionnel : il nuit à la clarté et donne une image négligée de l'expéditeur.
\end{enumerate}
\end{corrige}

\begin{corrige}[Exercice 3 — Définitions]
\begin{enumerate}
\item \textbf{La lettre de motivation} : lettre officielle, d'une page au maximum, dans laquelle un candidat expose ses raisons de postuler à un emploi, un stage ou une formation, et met en valeur son profil. \emph{Exemple :} la lettre adressée au Directeur des Ressources Humaines d'AGROCAM pour un poste de technicien.
\item \textbf{La phrase complexe} : phrase qui contient au moins deux propositions, dont l'une est subordonnée à l'autre (relative, conjonctive, participiale). \emph{Exemple :} « J'ai lu l'offre \emph{que vous avez publiée}. »
\item \textbf{L'emprunt} : mot qu'une langue prend à une langue étrangère pour nommer une réalité nouvelle ou par prestige. \emph{Exemple :} « week-end », emprunté à l'anglais.
\item \textbf{La netiquette} : ensemble des règles de politesse et de bon usage qui régissent la communication par internet (courriels, forums, réseaux sociaux). \emph{Exemple :} ne pas écrire en majuscules, ce qui équivaut à crier.
\item \textbf{La tonalité dramatique} : registre qui cherche à émouvoir le lecteur en présentant des situations de souffrance, de danger ou de tension. \emph{Exemple :} le récit d'un chômeur qui décrit sa détresse pour obtenir de l'aide.
\end{enumerate}
\end{corrige}

\begin{corrige}[Exercice 4 — Analyse grammaticale directe]
\begin{enumerate}
\item « qui m'a permis de maîtriser la maintenance industrielle » : subordonnée \textbf{relative} (pronom relatif « qui ») ; « qui » est \textbf{sujet} du verbe « a permis » ; la proposition est épithète de l'antécédent « une formation ».
\item « pour que nous convenions d'un entretien » : subordonnée \textbf{conjonctive} introduite par « pour que » (+ subjonctif) ; proposition \textbf{complément circonstanciel de but} de « je reste à votre disposition ».
\item « Titulaire d'un Baccalauréat technique » : subordonnée \textbf{participiale} (participe passé employé comme adjectif en tête de phrase, sans sujet exprimé propre) ; elle a une valeur \textbf{circonstancielle de cause} (= parce que je suis titulaire de...).
\item « que vous avez publiée » : subordonnée \textbf{relative} ; « que » est \textbf{complément d'objet direct} du verbe « avez publiée » (d'où l'accord du participe passé avec « l'offre ») ; épithète de l'antécédent « L'offre ».
\end{enumerate}
\textbf{Point de vigilance :} ne pas confondre « que » pronom relatif (remplaçable par « lequel ») et « que » conjonction de subordination (suivi d'un sujet exprimé différent).
\end{corrige}

\begin{corrige}[Exercice 5 — Rédiger l'en-tête]
En-tête attendu :

\begin{center}
\begin{tabularx}{\linewidth}{X X}
Mballa Éric & \\
Quartier Nkolndongo & \\
BP 5123 Yaoundé & \\
Tél. : 6 90 12 34 56 & \\
mballa.eric@mail.cm & \\
& CAMBOIS SARL\\
& Avenue Ahmadou Ahidjo\\
& BP 208 Douala\\
& \\
\multicolumn{2}{r}{Yaoundé, le 20 octobre 2025}\\
\end{tabularx}
\end{center}

\textbf{Objet :} Candidature au poste de technicien en maintenance.

\textbf{Réf. :} Offre n° 17/CB/2025.

\textbf{Points de vigilance :} coordonnées de l'expéditeur en haut à gauche, celles du destinataire à droite et légèrement plus bas ; lieu et date à droite, sous les coordonnées du destinataire ; objet précis (ni « truc » ni « candidature » seul) ; référence de l'offre reprise exactement (« Offre n° 17/CB/2025 ») ; date rédigée en toutes lettres (« le 20 octobre 2025 »).
\end{corrige}

\begin{corrige}[Exercice 6 — Construire des phrases complexes]
\begin{enumerate}
\item \textbf{Relative :} « J'ai obtenu un Baccalauréat technique \emph{qui} m'a donné de solides bases en électrotechnique. » (ou : « ...un Baccalauréat technique, \emph{lequel} m'a donné... »)
\item \textbf{Conjonctive de but :} « Je me permets de vous écrire \emph{pour que} je \emph{puisse} vous proposer ma candidature. » (\emph{Pour que} + subjonctif ; le sujet change ou non, la subordonnée conjonctive est requise. L'infinitif « afin de » est une alternative élégante mais ne constitue pas une proposition subordonnée conjonctive.)
\item \textbf{Participiale :} « \emph{Motivé par} le dynamisme de votre entreprise, je postule au stage que vous proposez. »
\item \textbf{Conjonctive de conséquence :} « Mon stage chez SODECOTON m'a \emph{tant} appris \emph{que} je peux désormais travailler en autonomie. »
\end{enumerate}
\textbf{Point de vigilance :} vérifier le mode verbal : subjonctif après « pour que », « afin que » ; indicatif après « si bien que », « tant... que » de conséquence.
\end{corrige}

\begin{corrige}[Exercice 7 — Les emprunts dans un texte professionnel]
\begin{center}
\begin{tabularx}{\linewidth}{|l|l|X|}
\hline
\rowcolor{popblueL} \textbf{Emprunt} & \textbf{Langue d'origine} & \textbf{Équivalent français soutenu} \\
\hline
start-up & anglais & jeune entreprise innovante \\
\hline
community manager & anglais & responsable des réseaux sociaux / animateur de communauté en ligne \\
\hline
mailing & anglais & courrier publicitaire / envoi groupé de courriels \\
\hline
chat & anglais & messagerie instantanée / clavardage \\
\hline
feedback & anglais & retour d'expérience / avis des clients \\
\hline
smartphone & anglais & téléphone intelligent / téléphone multifonction \\
\hline
design & anglais & conception graphique / stylisme \\
\hline
\end{tabularx}
\end{center}
(Cinq emprunts suffisent ; tous sont des \textbf{anglicismes}.) \textbf{Point de vigilance :} dans une lettre officielle, on préfère les équivalents français ; les anglicismes techniques admis (ex. « logiciel » est déjà français) ne posent pas de problème, mais « job », « boss », « cool » sont à bannir.
\end{corrige}

\begin{corrige}[Exercice 8 — Formules de conclusion]
\begin{center}
\begin{tabularx}{\linewidth}{|X|X|}
\hline
\rowcolor{popgreenL} \textbf{Correctes} & \textbf{Incorrectes} \\
\hline
1. « Veuillez agréer, Monsieur le Directeur, l'expression de ma considération distinguée. » & 2. « Salut et à bientôt j'espère ! » \\
\hline
3. « Dans l'attente d'une suite favorable, je vous prie de croire, Madame, à mes salutations distinguées. » & 4. « Bisous et merci d'avance. » \\
\hline
\end{tabularx}
\end{center}
\textbf{Corrections :}
\begin{itemize}
\item 2 $\rightarrow$ « Dans l'attente de votre réponse, je vous prie d'agréer, Monsieur, mes salutations distinguées. »
\item 4 $\rightarrow$ « Je vous remercie de l'attention portée à ma candidature et vous prie de croire, Monsieur le Directeur, à l'expression de ma considération distinguée. »
\end{itemize}
\textbf{Points de vigilance :} la formule de politesse reprend la civilité de la formule d'appel (« Madame » $\leftrightarrow$ « Madame ») ; « merci d'avance » est à éviter car il présuppose la réponse ; on écrit « l'expression de \emph{ma} considération » quand on est l'expéditeur.
\end{corrige}

\begin{corrige}[Exercice 9 — Type Baccalauréat : rédaction complète d'une lettre de motivation]
\textbf{Modèle de lettre (corrigé indicatif) :}

\begin{center}
\begin{tabularx}{\linewidth}{X X}
Fotso Larissa & \\
BP 777 Bafoussam & \\
& Monsieur le Directeur des Ressources Humaines\\
& AGROCAM SARL\\
& BP 1450 Bafoussam\\
\multicolumn{2}{r}{Bafoussam, le 20 octobre 2025}\\
\end{tabularx}
\end{center}

\textbf{Objet :} Candidature au poste de technicienne en agropastoral.

\textbf{Réf. :} REC/AGRO/25/09.

Monsieur le Directeur,

Votre société, qui figure parmi les entreprises agropastorales les plus dynamiques de la région de l'Ouest, recherche un technicien pour son site de Bafoussam. L'offre que vous avez publiée dans \emph{Cameroon Tribune} du 14 octobre 2025 correspond exactement au poste que je recherche.

Titulaire d'un Baccalauréat technique, série agropastorale, obtenu au Lycée technique de Bafoussam en juin 2025, j'ai acquis de solides connaissances en production agricole et en élevage. J'ai en outre effectué un stage de deux mois dans une exploitation agricole de Foumbot, qui m'a permis de développer le sens de l'organisation et le goût du travail en équipe.

Convaincue que mes compétences peuvent contribuer à la réussite de vos projets, je serais heureuse de vous exposer ma motivation lors d'un entretien, afin que nous envisagions ensemble une collaboration fructueuse.

Dans cette attente, je vous prie d'agréer, Monsieur le Directeur, l'expression de ma considération distinguée.

Fotso Larissa

\textbf{Barème indicatif (20 \%) :}
\begin{itemize}
\item En-tête complet et correctement disposé (coordonnées, lieu et date, destinataire, objet, référence) : 4 pts ;
\item Corps en trois paragraphes « vous – moi – nous », avec au moins deux phrases complexes de types différents : 8 pts ;
\item Formule d'appel et formule de politesse adaptées : 4 pts ;
\item Orthographe, syntaxe, présentation : 4 pts.
\end{itemize}
\textbf{Points de vigilance :} l'objet doit reprendre l'intitulé exact du poste ; la référence REC/AGRO/25/09 doit être recopiée sans erreur ; accorder « convaincue » avec le sujet féminin ; phrases complexes attendues, par exemple relative (« qui figure parmi... »), participiale (« Titulaire de... »), conjonctive de but (« afin que nous envisagions... », subjonctif).
\end{corrige}

\begin{corrige}[Exercice 10 — Type Probatoire : réécriture d'une lettre fautive]
\textbf{1. Relevé et classement des dix erreurs (5 pts) :}
\begin{center}
\begin{tabularx}{\linewidth}{|l|X|}
\hline
\rowcolor{popblueL} \textbf{Nature} & \textbf{Erreurs relevées} \\
\hline
Registre de langue / langage SMS & « Slt » ; « ça m'a grave branché » ; « trop chaud » ; « bosser » ; « c'était cool » ; « Biz » \\
\hline
Formules protocolaires & « M. le boss » (formule d'appel fautive) ; « Appelle-moi vite pour qu'on se capte » (impératif familier) ; « Biz » (formule finale familière) \\
\hline
Emprunts inadaptés & « le net », « le job », « ta boîte » (familier), « start-up », « mailing », « feedback » \\
\hline
Objet de la lettre & « Objet : truc » (vague et familier) \\
\hline
\end{tabularx}
\end{center}

\textbf{2. Formule d'appel et objet corrigés (2 pts) :} « Monsieur le Directeur, » — « Objet : candidature au poste de technicien. »

\textbf{3. Paragraphe central réécrit (8 pts) :} « J'ai pris connaissance de votre annonce publiée sur internet, et le poste de technicien que vous proposez correspond pleinement à mon profil. J'ai effectué un stage dans une jeune entreprise innovante, où j'assurais l'envoi des courriers publicitaires et le suivi des avis des clients, expérience qui m'a permis de développer rigueur et sens du contact. » (Emprunts remplacés : net $\rightarrow$ internet/site ; job $\rightarrow$ poste ; boîte $\rightarrow$ entreprise ; start-up $\rightarrow$ jeune entreprise innovante ; mailing $\rightarrow$ envoi de courriers publicitaires ; feedback $\rightarrow$ avis des clients. Phrase complexe : relative « qui m'a permis... ».)

\textbf{4. Formule de politesse (2 pts) :} « Dans l'attente d'une réponse favorable, je vous prie d'agréer, Monsieur le Directeur, l'expression de ma considération distinguée. »

\textbf{5. Netiquette (3 pts) :} la netiquette impose clarté, correction et respect du destinataire ; le langage SMS rend le message difficile à lire, multiplie les fautes et donne une image négligée et peu sérieuse du candidat, ce qui peut suffire à écarter sa candidature.

\textbf{Barème indicatif :} 5 + 2 + 8 + 2 + 3 = 20 \%. \textbf{Point de vigilance :} toute trace de langage SMS ou d'anglicisme non justifié dans la réécriture fait perdre les points du critère « registre soutenu ».
\end{corrige}

\begin{corrige}[Exercice 11 — Type Bac : texte iconique et courriel de candidature]
\textbf{Première partie : analyse de l'affiche.}

\textbf{1. Description (3 pts) :} sur fond vert (couleur de l'espérance et de la réussite), un jeune diplômé en costume, souriant, gravit une échelle géante au sommet de laquelle on lit « VOTRE AVENIR COMMENCE ICI » ; au pied de l'échelle, une foule grise et anonyme de candidats attend ; en bas figurent le logo de la banque et la mention invitant à envoyer lettre de motivation et CV avant le 30 novembre.

\textbf{2. Tonalité et procédés (4 pts) :} la tonalité dominante est \textbf{dramatique} (avec une visée incitative) : l'affiche met en scène une compétition angoissante. Deux procédés : le \emph{contraste des couleurs} (le vert lumineux de la réussite contre le gris terne de la foule des exclus) et la \emph{symbolique de l'échelle} (l'ascension sociale réservée à quelques-uns) ; on peut citer aussi l'\emph{opposition individu/foule} (un seul sourit et monte, les autres attendent).

\textbf{3. Effet recherché (3 pts) :} provoquer chez le candidat une inquiétude stimulante (« serai-je de ceux qui montent ? ») et le pousser à agir vite : soigner sa lettre de motivation pour se distinguer de la masse et postuler avant la date limite.

\textbf{Deuxième partie : courriel (6 pts).}

\textbf{Objet :} Candidature au poste de [intitulé du poste] — lettre de motivation et CV.

Monsieur le Directeur des Ressources Humaines,

Suite à votre annonce de recrutement, j'ai l'honneur de vous adresser ma candidature. Vous trouverez en pièces jointes ma lettre de motivation ainsi que mon curriculum vitae.

Me tenant à votre disposition pour tout renseignement complémentaire, je vous prie d'agréer, Monsieur le Directeur, l'expression de ma considération distinguée.

[Nom Prénom] — [téléphone] — [adresse électronique]

\textbf{5. Justification de l'objet (2 pts) :} l'objet est bref, précis et informatif : il annonce la nature du message (une candidature), le poste visé et le contenu (pièces jointes), ce qui permet au destinataire de classer et de traiter le courriel sans l'ouvrir.

\textbf{6. Deux règles de netiquette appliquées (2 pts) :} par exemple : ne pas utiliser le langage SMS ni d'abréviations ; rédiger un objet explicite ; soigner orthographe et ponctuation ; ne pas écrire en majuscules ; signer avec ses coordonnées complètes ; joindre réellement les fichiers annoncés.

\textbf{Barème indicatif :} 3 + 4 + 3 + 6 + 2 + 2 = 20 \%. \textbf{Points de vigilance :} la description doit rester objective (pas d'interprétation dans la question 1) ; la tonalité doit être justifiée par des procédés précis relevés dans l'image ; le courriel doit annoncer explicitement les pièces jointes et respecter les formules protocolaires d'une lettre officielle.
\end{corrige}

\newpage

\coursec{Fiche bilan}

\begin{popfigure}
\begin{center}\tcbox[colback=popink!80,colframe=popink!80,coltext=white,arc=9pt,boxsep=4pt,left=6pt,right=6pt]{\bfseries\small La lettre de motivation}\end{center}
\vspace{-2pt}
\begin{tcbraster}[raster columns=3,raster equal height=rows,raster column skip=6pt,raster row skip=6pt,enhanced,blankest]
\begin{tcolorbox}[enhanced,colback=white,colframe=popblue!70,boxrule=0.9pt,arc=3pt,title={Textes fonctionnels},fonttitle=\bfseries\scriptsize,coltitle=white,colbacktitle=popblue!70,left=4pt,right=4pt,top=2pt,bottom=2pt,boxsep=1pt,toptitle=1.5pt,bottomtitle=1.5pt,halign=left]
\begin{itemize}[nosep,leftmargin=*,itemsep=1pt,font=\scriptsize]
\item lettre, CV, mail
\item codes d'internet
\end{itemize}
\end{tcolorbox}
\begin{tcolorbox}[enhanced,colback=white,colframe=poporange!80,boxrule=0.9pt,arc=3pt,title={L'en-tête},fonttitle=\bfseries\scriptsize,coltitle=white,colbacktitle=poporange!80,left=4pt,right=4pt,top=2pt,bottom=2pt,boxsep=1pt,toptitle=1.5pt,bottomtitle=1.5pt,halign=left]
\begin{itemize}[nosep,leftmargin=*,itemsep=1pt,font=\scriptsize]
\item expéditeur / destinataire
\item lieu, date, objet
\end{itemize}
\end{tcolorbox}
\begin{tcolorbox}[enhanced,colback=white,colframe=popgreen!70,boxrule=0.9pt,arc=3pt,title={Phrase complexe},fonttitle=\bfseries\scriptsize,coltitle=white,colbacktitle=popgreen!70,left=4pt,right=4pt,top=2pt,bottom=2pt,boxsep=1pt,toptitle=1.5pt,bottomtitle=1.5pt,halign=left]
\begin{itemize}[nosep,leftmargin=*,itemsep=1pt,font=\scriptsize]
\item subordination
\item coordination
\end{itemize}
\end{tcolorbox}
\begin{tcolorbox}[enhanced,colback=white,colframe=popgold!85,boxrule=0.9pt,arc=3pt,title={L'emprunt},fonttitle=\bfseries\scriptsize,coltitle=white,colbacktitle=popgold!85,left=4pt,right=4pt,top=2pt,bottom=2pt,boxsep=1pt,toptitle=1.5pt,bottomtitle=1.5pt,halign=left]
\begin{itemize}[nosep,leftmargin=*,itemsep=1pt,font=\scriptsize]
\item mot étranger intégré
\end{itemize}
\end{tcolorbox}
\begin{tcolorbox}[enhanced,colback=white,colframe=poppurple!70,boxrule=0.9pt,arc=3pt,title={Le corps},fonttitle=\bfseries\scriptsize,coltitle=white,colbacktitle=poppurple!70,left=4pt,right=4pt,top=2pt,bottom=2pt,boxsep=1pt,toptitle=1.5pt,bottomtitle=1.5pt,halign=left]
\begin{itemize}[nosep,leftmargin=*,itemsep=1pt,font=\scriptsize]
\item motivation, atouts
\item tonalités : comique, dramatique
\end{itemize}
\end{tcolorbox}
\begin{tcolorbox}[enhanced,colback=white,colframe=poppink!80,boxrule=0.9pt,arc=3pt,title={Conclusion},fonttitle=\bfseries\scriptsize,coltitle=white,colbacktitle=poppink!80,left=4pt,right=4pt,top=2pt,bottom=2pt,boxsep=1pt,toptitle=1.5pt,bottomtitle=1.5pt,halign=left]
\begin{itemize}[nosep,leftmargin=*,itemsep=1pt,font=\scriptsize]
\item formules de politesse
\end{itemize}
\end{tcolorbox}
\begin{tcolorbox}[enhanced,colback=white,colframe=popmuted!80,boxrule=0.9pt,arc=3pt,title={Rédaction},fonttitle=\bfseries\scriptsize,coltitle=white,colbacktitle=popmuted!80,left=4pt,right=4pt,top=2pt,bottom=2pt,boxsep=1pt,toptitle=1.5pt,bottomtitle=1.5pt,halign=left]
\begin{itemize}[nosep,leftmargin=*,itemsep=1pt,font=\scriptsize]
\item brouillon, relecture
\item amélioration
\end{itemize}
\end{tcolorbox}
\end{tcbraster}

\legende{Carte mentale de la leçon : les sept savoirs essentiels autour de la lettre de motivation.}
\end{popfigure}

\begin{bilan}
\textbf{Définitions clés}
\begin{itemize}
  \item \cle{Texte fonctionnel} : texte utilitaire de la vie courante (lettre, CV, formulaire, courriel) visant une action précise.
  \item \cle{Lettre de motivation} : lettre officielle qui accompagne un CV pour solliciter un emploi, un stage ou une formation ; elle doit convaincre le destinataire.
  \item \cle{Codes de la communication par internet} : règles de politesse et de clarté du courriel (objet explicite, salutations, pas d'abréviations SMS, signature).
  \item \cle{Phrase complexe} : phrase contenant au moins deux verbes conjugués, reliés par coordination ou subordination.
  \item \cle{Emprunt} : mot pris à une langue étrangère et intégré au français (ex. \emph{business}, \emph{week-end}).
  \item \cle{Tonalité} : effet produit sur le lecteur ; la tonalité \cle{dramatique} suscite l'émotion, la tonalité \cle{comique} provoque le rire.
\end{itemize}

\textbf{Structure de la lettre de motivation (à connaître par cœur)}
\begin{center}
\begin{tabularx}{\linewidth}{|l|X|}
\hline
\rowcolor{popblueL} \textbf{Partie} & \textbf{Contenu obligatoire} \\
\hline
En-tête & Coordonnées de l'expéditeur (en haut à gauche), du destinataire (à droite), lieu et date, objet de la lettre \\
\hline
Corps & Accroche, présentation de soi, motivation, atouts (formation, expérience, qualités), phrase complexe valorisante \\
\hline
Conclusion & Formule de politesse adaptée, signature \\
\hline
\end{tabularx}
\end{center}

\textbf{Méthodes express}
\begin{itemize}
  \item \textbf{Rédiger} : brouillon $\rightarrow$ plan (en-tête, corps, conclusion) $\rightarrow$ rédaction $\rightarrow$ relecture (orthographe, formules) $\rightarrow$ confrontation et amélioration.
  \item \textbf{Phrase complexe} : repérer les verbes conjugués, puis le mot de liaison (\emph{que, qui, parce que, si}) : subordonnée ou juxtaposée ?
  \item \textbf{Formule de conclusion type} : « Dans l'attente d'une réponse favorable, je vous prie d'agréer, Madame, Monsieur, l'expression de mes salutations distinguées. »
\end{itemize}
\end{bilan}

\begin{aretenir}[Checklist avant le Bac]
$\square$ Je sais définir un texte fonctionnel et citer des exemples.\par
$\square$ Je connais les codes de politesse de la communication par internet.\par
$\square$ Je place correctement l'expéditeur, le destinataire, le lieu, la date et l'objet dans l'en-tête.\par
$\square$ Je rédige un objet de lettre court et précis.\par
$\square$ Je reconnais et je construis une phrase complexe (subordination, coordination).\par
$\square$ Je sais définir l'emprunt et donner des exemples.\par
$\square$ Je rédige un corps de lettre : accroche, motivation, atouts.\par
$\square$ Je distingue la tonalité dramatique de la tonalité comique.\par
$\square$ Je connais au moins une formule de politesse de conclusion par cœur.\par
$\square$ Je relis ma lettre : orthographe, ponctuation, présentation.\par
$\square$ Je sais améliorer ma production après confrontation avec celle des camarades.\par
$\square$ Je rédige une lettre de motivation complète en temps limité.
\end{aretenir}

\findecours

\end{document}
